\documentclass[a4paper,12pt]{article}

% Пакеты для русского языка и кодировки
\usepackage[utf8]{inputenc}
\usepackage[T2A]{fontenc}
\usepackage[russian]{babel}

% Пакеты для математики
\usepackage{amsmath}
\usepackage{amssymb}
\usepackage{amsthm}

% Пакеты для оформления
\usepackage{geometry}
\geometry{left=2.5cm, right=2.5cm, top=2.5cm, bottom=2.5cm}
\usepackage{hyperref}
\usepackage{enumitem}
\usepackage{booktabs}
\usepackage{array}

% Настройка заголовков
\title{2.3.5 Материалы для изучения по научной специальности \\
Математическое и программное обеспечение вычислительных систем, комплексов и компьютерных сетей \\
Гусев Виталий Евгеньевич 2.3.5 2026}
\date{}

\begin{document}

\maketitle

\tableofcontents
\newpage

\section{Раздел 3. Вычислительные машины, системы и сети. Языки программирования.}

\subsection{Архитектура современных компьютеров. Организация памяти и архитектура процессора современных вычислительных машин}

\noindent\textit{Основные источники: Орлов С.А., Цилькер Б.Я. Организация ЭВМ и систем: Учебник для вузов. 2-е изд. СПб.: Питер, 2011. Гл. 1, с. 20--53; гл. 2, с. 55--117; гл. 6, с. 230--331; гл. 9, с. 386--449.}

\subsubsection{Понятия архитектуры и организации вычислительной машины}

\paragraph{Архитектура и организация.}
Под \emph{архитектурой вычислительной машины} в широком смысле понимается концептуальная структура ВМ, определяющая способы обработки информации, методы преобразования информации в данные и принципы взаимодействия технических средств и программного обеспечения. Международный стандарт ISO/IEC 2382/1-93 трактует архитектуру как логическую структуру и функциональные характеристики ВМ, включая взаимосвязи между аппаратными и программными компонентами. В более узкой трактовке, восходящей к фирме IBM и семейству IBM 360, архитектура охватывает те средства, которые видны программисту: систему команд, форматы данных и команд, способы адресации, механизмы ввода/вывода и т.п.
\textit{(Источник: Орлов, Цилькер, гл. 1, с. 21--22.)}

При этом различают \emph{функциональную организацию} и \emph{структурную организацию}. Функциональная организация --- это абстрактная модель совокупности функциональных возможностей и услуг, предоставляемых пользователю; она отвечает на вопрос, \emph{что} умеет делать машина. Структурная организация --- физическая модель, устанавливающая состав, порядок и принципы взаимодействия основных функциональных частей; она отвечает на вопрос, \emph{как} эти возможности реализованы технически. Графическим отображением структурной организации служит структурная схема.
\textit{(Источник: Орлов, Цилькер, гл. 1, с. 21--22.)}

\paragraph{Уровни детализации структуры ВМ.}
При изучении современных компьютеров обычно выделяют четыре уровня детализации. На уровне «чёрного ящика» ВМ рассматривается как устройство, способное хранить и обрабатывать информацию и обмениваться данными с внешним миром. На уровне общей архитектуры выделяются центральный процессор (ЦП), основная память (ОП), устройства ввода/вывода (УВВ) и система шин. На уровне архитектуры центрального процессора выделяются арифметико-логическое устройство (АЛУ), блок обработки чисел с плавающей запятой (БПЗ), регистры, устройство управления (УУ) и внутренние шины. Наконец, на уровне устройства управления рассматриваются логика программной последовательности, регистры и дешифраторы, управляющая память и логика формирования управляющих сигналов.
\textit{(Источник: Орлов, Цилькер, гл. 1, с. 22--24.)}

\paragraph{Функциональная схема фон-неймановской ВМ.}
Типичная фон-неймановская ВМ содержит центральный процессор, основную память, устройство ввода/вывода и систему шин. В составе процессора выделяются устройство управления, арифметико-логическое устройство и регистры. Устройство управления организует автоматическое выполнение программы: формирует управляющие сигналы, обеспечивает выборку команд из памяти в порядке, определяемом программой, и их исполнение, а также синхронизирует работу внутренних и внешних устройств. АЛУ выполняет арифметические и логические операции над операндами и формирует признаки результата (флаги): равенство нулю, знак, перенос, переполнение и т.п. Регистры служат для кратковременного хранения команд, данных и адресов. Основная память представляет собой массив пронумерованных ячеек, доступ к которым возможен в произвольном порядке; команды и данные хранятся в одной и той же памяти и внешне неразличимы.
\textit{(Источник: Орлов, Цилькер, гл. 1, с. 39--42.)}

\subsubsection{Фон-неймановская архитектура и принципы построения}

\paragraph{Принципы фон Неймана.}
Классическая архитектура современных ВМ восходит к концепции машины с хранимой в памяти программой, сформулированной Джоном фон Нейманом в 1945 г. Её основные принципы таковы:
\begin{enumerate}[label=\arabic*)]
\item \emph{принцип двоичного кодирования}: все данные и команды кодируются двоичными цифрами 0 и 1; каждый тип информации имеет свой формат, а последовательность битов, имеющая определённый смысл, называется полем;
\item \emph{принцип программного управления}: все вычисления представляются в виде программы, состоящей из последовательности управляющих слов --- команд; команды хранятся в памяти и выполняются в естественном порядке, то есть в порядке их расположения в программе; при необходимости порядок может быть изменён специальными командами;
\item \emph{принцип однородности памяти}: команды и данные хранятся в одной и той же памяти и внешне неразличимы; распознать их можно только по способу использования;
\item \emph{принцип адресуемости памяти}: основная память состоит из пронумерованных ячеек, причём процессору в произвольный момент доступна любая ячейка; двоичные коды команд и данных разделяются на единицы информации, называемые словами, и хранятся в ячейках памяти, а для доступа к ним используются номера соответствующих ячеек --- адреса.
\end{enumerate}
\textit{(Источник: Орлов, Цилькер, гл. 1, с. 37--39.)}

\paragraph{Структуры вычислительных машин.}
В настоящее время примерно одинаковое распространение получили два способа построения вычислительных машин: с непосредственными связями и на основе шины. В машине с непосредственными связями между взаимодействующими устройствами (процессор, память, устройство ввода/вывода) имеются прямые связи; достоинством является возможность «развязки» узких мест путём улучшения характеристик только определённых связей, недостатком --- плохая податливость реконфигурации. В машине с общей шиной все устройства подключены к магистральной шине, служащей единственным трактом для потоков команд, данных и управления; это упрощает реализацию и позволяет легко менять конфигурацию, но в каждый момент передавать информацию по шине может только одно устройство, что ограничивает производительность. Поэтому чаще применяется архитектура с иерархией шин, где помимо магистральной имеются дополнительные шины для наиболее интенсивных обменов (например, между процессором и кэш-памятью) или для объединения однотипных устройств ввода/вывода.
\textit{(Источник: Орлов, Цилькер, гл. 1, с. 42--43.)}

\paragraph{Типы структур вычислительных систем.}
Вычислительные системы строятся по одному из двух подходов. В системах с общей памятью имеется общая основная память, совместно используемая всеми процессорами; связь процессоров с памятью обеспечивается коммуникационной сетью, чаще всего вырождающейся в общую шину. В распределённых системах общая память отсутствует, а каждый процессор обладает собственной локальной памятью; обмен информацией между составляющими системы обеспечивается с помощью коммуникационной сети посредством обмена сообщениями. Первый подход характерен для сильно связанных (tightly coupled) систем, второй --- для слабо связанных (loosely coupled) систем.
\textit{(Источник: Орлов, Цилькер, гл. 1, с. 43--44.)}

\subsubsection{Архитектура системы команд}

\paragraph{Понятие архитектуры системы команд.}
Системой команд вычислительной машины называют полный перечень команд, которые способна выполнять данная ВМ. Под \emph{архитектурой системы команд} (АСК) принято определять те средства вычислительной машины, которые видны и доступны программисту. АСК можно рассматривать как линию согласования нужд разработчиков программного обеспечения с возможностями создателей аппаратуры. Время выполнения программы приближённо оценивается выражением
\[
T_{\text{выч}}=N_{\text{ком}}\times CPI\times \tau_{\text{пр}},
\]
где \(N_{\text{ком}}\) --- число команд в программе, \(CPI\) --- среднее число тактов процессора на одну команду, \(\tau_{\text{пр}}\) --- длительность тактового периода. Каждая из составляющих зависит от одних аспектов АСК и влияет на другие, что требует ответственного подхода к выбору АСК.
\textit{(Источник: Орлов, Цилькер, гл. 2, с. 55--56.)}

Общая характеристика АСК складывается из ответов на следующие вопросы: какого вида данные будут представлены в ВМ и в какой форме; где эти данные могут храниться помимо основной памяти; каким образом будет осуществляться доступ к данным; какие операции могут быть выполнены над данными; сколько операндов может присутствовать в команде; как будет определяться адрес очередной команды; каким образом будут закодированы команды.
\textit{(Источник: Орлов, Цилькер, гл. 2, с. 56.)}

\paragraph{Классификация АСК по составу и сложности команд.}
Современная технология программирования ориентирована на языки высокого уровня (ЯВУ), что порождает проблему семантического разрыва: сложные операторы ЯВУ существенно отличаются от простых машинных операций. Для преодоления разрыва разработчики выбирают один из трёх подходов, соответствующих трём типам АСК:
\begin{enumerate}[label=\arabic*)]
\item \emph{CISC} (Complex Instruction Set Computer) --- архитектура с полным набором команд. Семантический разрыв преодолевается за счёт расширения системы команд, дополнения её сложными командами, семантически аналогичными операторам ЯВУ. Для CISC типичны: наличие в процессоре сравнительно небольшого числа регистров общего назначения; большое количество машинных команд, часть которых аппаратно реализуют сложные операторы ЯВУ; разнообразие способов адресации операндов; множество форматов команд различной разрядности; наличие команд, где обработка совмещается с обращением к памяти. К типу CISC можно отнести практически все ВМ, выпускавшиеся до середины 1980-х годов, и значительную часть производящихся в настоящее время.
\item \emph{RISC} (Reduced Instruction Set Computer) --- архитектура с сокращённым набором команд. Идея заключается в ограничении списка команд ВМ наиболее часто используемыми простейшими командами, оперирующими данными, размещёнными только в регистрах процессора. Обращение к памяти допускается лишь с помощью специальных команд чтения и записи. Резко уменьшено количество форматов команд и способов указания адресов операндов. Эти меры позволяют существенно упростить аппаратные средства ВМ и повысить их быстродействие. RISC-архитектура разрабатывалась так, чтобы уменьшить \(T_{\text{выч}}\) за счёт сокращения \(CPI\) и \(\tau_{\text{пр}}\).
\item \emph{VLIW} (Very Long Instruction Word) --- архитектура с командными словами сверхбольшой длины. Концепция базируется на RISC-архитектуре, но несколько простых RISC-команд объединяются в одну сверхдлинную команду и выполняются параллельно. Появился дополнительный уровень параллелизма вычислений, в силу чего архитектуру VLIW логичнее адресовать не к вычислительным машинам, а к вычислительным системам.
\end{enumerate}
\textit{(Источник: Орлов, Цилькер, гл. 2, с. 57--59.)}

\paragraph{Классификация АСК по месту хранения операндов.}
С позиций места хранения операндов различают следующие виды АСК: стековую, аккумуляторную, регистровую и с выделенным доступом к памяти.
\begin{itemize}
\item В \emph{стековой} АСК операнды перед обработкой помещаются в две верхние ячейки стековой памяти; результат операции заносится в стек. Стек --- это память, работающая по принципу «последним вошёл --- первым вышел» (LIFO). Достоинства: сокращение адресной части команд, компактный код программы, простота декодирования. Недостатки: отсутствие произвольного доступа к памяти, трудности создания эффективного программного кода, стек становится «узким местом» при повышении производительности.
\item В \emph{аккумуляторной} АСК для хранения одного из операндов и результата используется выделенный регистр --- аккумулятор. В командах обработки достаточно явно указать местоположение только второго операнда. Достоинства: короткие команды, простота декодирования. Недостаток: наличие всего одного регистра порождает многократные обращения к основной памяти.
\item В \emph{регистровой} АСК процессор включает массив регистров общего назначения (РОН). К любому регистру можно обратиться, указав его номер. Количество РОН в CISC-архитектурах обычно невелико (от 8 до 32), в RISC-архитектурах --- значительно больше (до нескольких сотен). Регистровая АСК допускает расположение операндов как в регистрах, так и в основной памяти, поэтому выделяют три формата команд обработки: «регистр--регистр», «регистр--память», «память--память». Достоинства: компактность кода, высокая скорость вычислений за счёт замены обращений к основной памяти на обращения к быстрым регистрам. Недостаток: более длинные команды по сравнению с аккумуляторной архитектурой.
\item В АСК \emph{с выделенным доступом к памяти} (Load/Store architecture) обращение к основной памяти возможно только с помощью двух специальных команд: load и store. Операнды во всех командах обработки информации могут находиться только в регистрах процессора; результат операции также заносится в регистр. Эта архитектура характерна для всех ВМ с RISC-архитектурой. Достоинства: простота декодирования и исполнения команд.
\end{itemize}
\textit{(Источник: Орлов, Цилькер, гл. 2, с. 59--67.)}

\paragraph{Типы и форматы операндов.}
Машинные команды оперируют данными, которые принято называть операндами. К основным (базовым) типам операндов относятся числа, символы и логические данные. Числовая информация классифицируется на числа с фиксированной запятой (целые, дробные, смешанные), числа с плавающей запятой и логические значения.
\begin{itemize}
\item Числа в форме с фиксированной запятой (ФЗ) представляются в виде \(A=\pm a_{m-1}\ldots a_1a_0,a_{-1}a_{-2}\ldots a_{-r}\), где \(m\) разрядов используются для записи целой части, \(r\) разрядов --- для дробной. Обычно применяется формат с фиксацией запятой после младшего разряда (целые числа). Диапазон представления целых чисел со знаком в дополнительном коде: от \(-2^{n-1}\) до \(2^{n-1}-1\).
\item Числа в форме с плавающей запятой (ПЗ) представляются в виде \(X=\pm m q^{p}\), где \(m\) --- мантисса, \(p\) --- порядок, \(q\) --- основание системы счисления. Мантисса обычно нормализована: \(1>|m|\ge 1/q\). Для упрощения операций над порядками применяется смещённый порядок: к истинному порядку добавляется целое положительное число --- смещение. Стандарт IEEE 754 определяет 32-битовый (одинарный) и 64-битовый (двойной) форматы; в них используются скрытая единица, смещённый порядок и специальные коды для нуля, бесконечности и NAN.
\item Упакованные целые числа применяются для обработки мультимедийной информации. Упаковка предполагает объединение в пределах достаточно длинного слова нескольких целых чисел меньшей разрядности, а соответствующие команды обрабатывают все эти числа параллельно (технологии MMX, SSE).
\item Десятичные числа представляются в двоично-десятичном коде (BCD). Используются два формата: упакованный (в каждом байте хранятся две десятичные цифры, правая тетрада младшего байта отводится под знак) и зонный (под каждую цифру выделяется один байт, где младшие четыре разряда отводятся под код цифры, а старшая тетрада содержит код «зоны»; исключение составляет байт, содержащий младшую цифру, где в поле зоны хранится знак числа).
\item Символьная информация кодируется с помощью таблиц кодировки. Исторически распространены EBCDIC и ASCII; в настоящее время широко применяется Unicode (UTF-8, UTF-16 и др.).
\item Логические данные представляются битом: единицей кодируют истинное значение, нулём --- ложное; над наборами логических переменных выполняются логические операции без переносов между разрядами.
\item Строки --- это непрерывная последовательность битов, байтов, слов или двойных слов. Для указания конца текстовой строки в последний байт заносится код-ограничитель (обычно нули) либо длина строки указывается в первом или двух первых байтах.
\end{itemize}
\textit{(Источник: Орлов, Цилькер, гл. 2, с. 67--90.)}

\paragraph{Типы команд.}
Несмотря на различие в системах команд разных ВМ, основные типы операций могут быть найдены в любой из них. Классификация типов команд включает: команды пересылки данных; команды арифметической и логической обработки; команды работы со строками; команды SIMD; команды преобразования; команды ввода/вывода; команды управления системой; команды управления потоком команд. Команды пересылки данных обеспечивают передачу информации между процессором и ОП, внутри процессора и между ячейками памяти. Команды арифметической и логической обработки выполняют операции над целыми числами, числами с плавающей запятой и логическими данными, формируя признаки результата (Z, N, V, C). Команды управления потоком команд включают безусловные переходы, условные переходы (ветвления), вызовы процедур и возвраты из процедур.
\textit{(Источник: Орлов, Цилькер, гл. 2, с. 90--97.)}

\paragraph{Форматы команд.}
Формат команды определяет её структуру, то есть количество двоичных разрядов, отводимых под всю команду, а также количество и расположение отдельных полей. Типовая команда должна указывать: подлежащую выполнению операцию; адреса исходных данных (операндов), над которыми выполняется операция; адрес, по которому должен быть помещён результат операции. В соответствии с этим команда состоит из двух частей: операционной и адресной. Обобщённый формат команды имеет вид: КОП, СА, адресная часть. Длина команды \(R_K\) описывается соотношением \(R_K=\sum_{i=1}^{l} R_{A_i}+R_{\text{КОП}}+R_{\text{СА}}\), где \(l\) --- количество адресов в команде, \(R_{A_i}\) --- количество разрядов для записи \(i\)-го адреса, \(R_{\text{КОП}}\) --- разрядность поля кода операции, \(R_{\text{СА}}\) --- разрядность поля способа адресации. Разрядность поля кода операции выбирается так, чтобы можно было представить любую из операций: \(R_{\text{КОП}}=\lceil \log_2 N_{\text{КОП}}\rceil\). Разрядность адресной части определяется количеством адресуемых ячеек памяти и количеством способов адресации.
\textit{(Источник: Орлов, Цилькер, гл. 2, с. 97--99.)}

По количеству адресов различают четырёхадресный, трёхадресный, двухадресный, одноадресный, полутораадресный и нульадресный форматы. В фон-неймановских ВМ надобность в четвёртом адресе (адресе следующей команды) отпадает, поскольку команды располагаются в памяти в порядке их выполнения. Трёхадресный формат позволяет указать адреса двух операндов и результата. Двухадресный формат предполагает, что результат записывается на место одного из операндов (обычно второго). Одноадресный формат использует для хранения первого операнда и результата специальный регистр --- аккумулятор. Полутораадресный (регистровый) формат содержит номер регистра и адрес второго операнда. Нульадресный формат не содержит адресной части; он используется в стековой архитектуре или для команд, не требующих операндов.
\textit{(Источник: Орлов, Цилькер, гл. 2, с. 99--101.)}

\paragraph{Способы адресации операндов.}
Способ адресации --- это способ формирования исполнительного адреса операнда по адресному коду команды. Исполнительным адресом операнда называется двоичный код номера ячейки памяти, служащей источником или приёмником операнда; адресный код команды --- это двоичный код в адресном поле команды, из которого необходимо сформировать исполнительный адрес. В современных ВМ исполнительный адрес и адресный код, как правило, не совпадают. Основные способы адресации:
\begin{itemize}
\item \emph{непосредственная адресация}: в адресном поле команды содержится сам операнд; сокращает время выполнения команды и экономит память, но размер операнда ограничен длиной адресного поля;
\item \emph{прямая (абсолютная) адресация}: адресный код прямо указывает номер ячейки памяти; простой способ, но ограничивает адресное пространство и не позволяет изменять адрес в процессе вычислений;
\item \emph{косвенная адресация}: адресный код указывает адрес ячейки, в которой хранится исполнительный адрес операнда; позволяет изменять адрес в процессе вычислений, но требует двух обращений к памяти;
\item \emph{регистровая адресация}: адресное поле указывает на регистр процессора; короткое адресное поле, исключаются обращения к памяти;
\item \emph{косвенная регистровая адресация}: исполнительный адрес операнда хранится в регистре процессора; требует одного обращения к памяти;
\item \emph{адресация со смещением}: исполнительный адрес формируется суммированием содержимого адресного поля команды с содержимым одного или нескольких регистров; к этой группе относятся относительная адресация (сложение с содержимым счётчика команд), базовая регистровая адресация, индексная адресация, страничная адресация; возможны автоинкрементная и автодекрементная адресации, а также адресация с масштабированием;
\item \emph{блочная адресация}: используется в командах, для которых единицей обработки служит блок данных, расположенных в смежных ячейках памяти;
\item \emph{стековая адресация}: операнды и результаты размещаются в стеке.
\end{itemize}
\textit{(Источник: Орлов, Цилькер, гл. 2, с. 103--114.)}

\subsubsection{Организация памяти}

\paragraph{Иерархия памяти.}
Память современной ВМ строится по иерархическому принципу. Верхний уровень образуют регистры процессора, далее следуют кэш-память первого, второго и третьего уровней, основная память и внешняя память. По мере движения вниз по иерархии возрастают ёмкость и время доступа, но снижается стоимость хранения бита. Эффективность иерархии основана на свойстве локальности по обращению: программа в каждый момент времени использует относительно небольшую часть своего адресного пространства. Различают пространственную локальность (обращение к соседним адресам), временную локальность (повторное обращение к тем же данным или командам) и локальность переходов. Правило «90/10» гласит, что около 90\% времени работы программы связано с доступом к 10\% адресного пространства. Иерархическая организация памяти позволяет сочетать высокое быстродействие верхних уровней с большой ёмкостью нижних уровней.
\textit{(Источник: Орлов, Цилькер, гл. 6, с. 230--235.)}

Основными характеристиками запоминающих устройств являются ёмкость, единица пересылки, метод доступа, быстродействие, физический тип, энергозависимость и стоимость. Для внутренней памяти характерны произвольный и ассоциативный методы доступа. Быстродействие оценивается временем выборки, временем хранения, циклом обращения и скоростью передачи данных. Иерархическая организация памяти позволяет сочетать высокое быстродействие верхних уровней с большой ёмкостью нижних уровней.
\textit{(Источник: Орлов, Цилькер, гл. 6, с. 230--235.)}

\paragraph{Основная память.}
Основная память представляет собой единственный вид памяти, к которой ЦП может обращаться непосредственно (исключение составляют лишь регистры центрального процессора). Информация, хранящаяся на внешних ЗУ, становится доступной процессору только после того, как будет переписана в основную память. Основную память образуют запоминающие устройства с произвольным доступом; такие ЗУ образованы как массив ячеек, а «произвольный доступ» означает, что обращение к любой ячейке занимает одно и то же время и может производиться в произвольной последовательности. Каждая ячейка содержит фиксированное число запоминающих элементов и имеет уникальный адрес, позволяющий отличать её от других ячеек.
\textit{(Источник: Орлов, Цилькер, гл. 6, с. 235--236.)}

Основная память может включать в себя два типа устройств: оперативные запоминающие устройства (ОЗУ) и постоянные запоминающие устройства (ПЗУ). Преимущественную долю основной памяти образует ОЗУ, допускающее как запись, так и считывание информации, причём обе операции выполняются однотипно, практически с одной и той же скоростью. Для большинства типов полупроводниковых ОЗУ характерна энергозависимость --- даже при кратковременном прерывании питания хранимая информация теряется. Вторую группу полупроводниковых ЗУ основной памяти образуют энергонезависимые микросхемы ПЗУ, обеспечивающие считывание информации, но либо вообще не допускающие её изменения, либо требующие для изменения значительно большего времени.
\textit{(Источник: Орлов, Цилькер, гл. 6, с. 236.)}

Емкость основной памяти современных ВМ слишком велика, чтобы её можно было реализовать на базе единственной интегральной микросхемы (ИМС). Необходимость объединения нескольких ИМС ЗУ возникает также, когда разрядность ячеек в микросхеме ЗУ меньше разрядности слов ВМ. Увеличение разрядности ЗУ реализуется за счёт объединения адресных входов объединяемых ИМС ЗУ; информационные входы и выходы микросхем являются входами и выходами модуля ЗУ увеличенной разрядности. Полученную совокупность микросхем называют модулем памяти. Модуль можно считать и единственную микросхему, если она уже имеет нужную разрядность. Один или несколько модулей образуют банк памяти.
\textit{(Источник: Орлов, Цилькер, гл. 6, с. 236--237.)}

При использовании блочной памяти, состоящей из \(B\) банков, адрес ячейки \(A\) преобразуется в пару \((b,w)\), где \(b\) --- номер банка, \(w\) --- адрес ячейки внутри банка. Известны три схемы распределения разрядов адреса \(A\) между \(b\) и \(w\): блочная (номер банка \(b\) определяет старшие разряды адреса); циклическая (\(b=A\bmod B\); \(w=A\operatorname{div} B\)); блочно-циклическая (комбинация двух предыдущих схем). Блочное построение памяти позволяет сократить время доступа к информации за счёт потенциального параллелизма. Одна из используемых для этого методик называется расслоением памяти; в её основе лежит чередование адресов, заключающееся в изменении системы распределения адресов между банками памяти. Приём чередования адресов базируется на свойстве локальности по обращению, согласно которому последовательный доступ в память обычно производится к ячейкам, имеющим смежные адреса.
\textit{(Источник: Орлов, Цилькер, гл. 6, с. 237--239.)}

\paragraph{Кэш-память.}
Кэш-память --- это небольшая по ёмкости, но высокобыстродействующая буферная память, размещаемая между процессором и основной памятью. В ней копируются наиболее часто используемые блоки основной памяти. При обращении процессора сначала проверяется наличие копии в кэш-памяти. Если копия есть, происходит попадание (hit) и данные берутся из кэша; если копии нет, происходит промах (miss), блок загружается из основной памяти в кэш, после чего требуемое слово передаётся процессору. Кэш-память состоит из памяти данных и памяти тегов; тег указывает, копия какого блока основной памяти находится в данном блоке кэша. Эффективность кэша оценивается коэффициентом попаданий \(H\) и потерями на промах. Среднее время доступа приближённо равно
\[
T_{\text{ср}}=T_{\text{hit}}+(1-H)T_{\text{miss}}.
\]
\textit{(Источник: Орлов, Цилькер, гл. 6, с. 277--280.)}

Известны три основных способа отображения основной памяти на кэш-память. При прямом отображении каждый блок основной памяти может быть помещён только в один вполне определённый блок кэш-памяти; проверка попадания сводится к сравнению тега в единственной ячейке памяти тегов. При полностью ассоциативном отображении любой блок основной памяти может быть помещён в любой блок кэш-памяти; для поиска нужного тега используется ассоциативная память. При частично-ассоциативном (множественно-ассоциативном) отображении кэш разбивается на модули, каждый из которых содержит несколько блоков; блок основной памяти отображается на определённый модуль, но внутри модуля может занимать любой блок. Частично-ассоциативное отображение является компромиссом между прямым и полностью ассоциативным и широко применяется в современных процессорах.
\textit{(Источник: Орлов, Цилькер, гл. 6, с. 280--283.)}

При заполненной кэш-памяти новый блок можно разместить только заменив один из старых. При прямом отображении выбор замещаемого блока однозначен; при ассоциативном и частично-ассоциативном используются алгоритмы замещения: LRU (наиболее давно использовавшийся), FIFO (первым вошёл --- первым вышел), LFU (наименее часто использовавшийся) и случайный выбор. Наиболее распространён алгоритм LRU. Для согласования содержимого кэша и основной памяти применяются методы сквозной записи (write through) и обратной записи (write back). При сквозной записи данные одновременно записываются в кэш и в основную память; при обратной записи изменения фиксируются только в кэше, а основная память обновляется при вытеснении блока. В многопроцессорных системах дополнительно решается проблема когерентности кэшей с помощью протоколов наблюдения (snoopy) или протоколов на основе справочника; широко известен протокол MESI.
\textit{(Источник: Орлов, Цилькер, гл. 6, с. 284--288; гл. 11, с. 473--492.)}

Современные процессоры используют многоуровневую кэш-память: L1 (раздельная для команд и данных), L2 (обычно общая для ядра или пары ядер) и L3 (общая для группы ядер). Многоуровневая организация позволяет сочетать малое время доступа L1 с большой ёмкостью L2 и L3.
\textit{(Источник: Орлов, Цилькер, гл. 6, с. 287--288.)}

\paragraph{Виртуальная память.}
Виртуальная память --- это метод автоматического управления иерархической памятью, при котором программисту кажется, что он располагает единой памятью большой ёмкости и высокого быстродействия. Программа пишется в виртуальных адресах; физическая основная память имеет физические адреса. Система виртуальной памяти автоматически отображает виртуальное пространство на физическое, загружает отсутствующие фрагменты с внешних запоминающих устройств и обеспечивает защиту памяти. Виртуальное пространство обычно больше физического; отсутствующая страница или сегмент при обращении вызывает страничный сбой (page fault), после чего операционная система загружает нужный фрагмент в основную память.
\textit{(Источник: Орлов, Цилькер, гл. 6, с. 288--290.)}

При страничной организации виртуальное и физическое адресные пространства разбиваются на блоки одинакового размера --- страницы. Виртуальный адрес делится на номер виртуальной страницы и смещение внутри страницы. Номер виртуальной страницы преобразуется в номер физической страницы (страничного кадра); смещение остаётся неизменным. Преобразование выполняется по страничной таблице. Каждая запись таблицы содержит номер физической страницы и признаки: \(V\) --- присутствие страницы в основной памяти, \(R\) --- признак использования, \(M\) --- признак модификации, \(A\) --- права доступа. Если \(V=0\), возникает страничный сбой, и страница загружается с диска. Для ускорения преобразования используется буфер быстрого преобразования адреса (TLB) --- небольшая ассоциативная кэш-память, хранящая недавние пары «номер виртуальной страницы --- номер физической страницы».
\textit{(Источник: Орлов, Цилькер, гл. 6, с. 290--293.)}

Страничная организация устраняет внешнюю фрагментацию, поскольку все страницы имеют одинаковый размер, но возможна внутренняя фрагментация последней страницы. Для сокращения объёма страничных таблиц применяются многоуровневые таблицы и инвертированные страничные таблицы. Признак \(M\) позволяет не записывать на диск неизменённые страницы при их вытеснении, а признак \(A\) используется для защиты и организации доступа.
\textit{(Источник: Орлов, Цилькер, гл. 6, с. 290--293.)}

При сегментной организации виртуальное пространство разбивается на сегменты переменной длины, каждый из которых является логически самостоятельной единицей: кодом, данными, стеком и т.п. Виртуальный адрес состоит из номера сегмента и смещения внутри сегмента. Преобразование выполняется по сегментной таблице, содержащей начальные адреса сегментов и их длины. Сегментная организация удобна для логической защиты и совместного использования памяти, но из-за переменной длины сегментов возможна внешняя фрагментация. В сегментно-страничной организации сегмент разбивается на страницы фиксированного размера; возникает иерархия «сегмент \(\to\) страница \(\to\) слово». Сегментная таблица указывает расположение страничных таблиц для каждого сегмента; страничная таблица определяет физические страницы. Такая схема сочетает логические преимущества сегментации с эффективностью страничного управления и уменьшает внешнюю фрагментацию. Для защиты памяти применяются граничные регистры, ключи защиты, кольца защиты и признаки прав доступа в таблицах страниц и сегментов.
\textit{(Источник: Орлов, Цилькер, гл. 6, с. 293--299.)}

\subsubsection{Архитектура процессора}

\paragraph{Конвейеризация вычислений.}
Совершенствование элементной базы уже не приводит к кардинальному росту производительности ВМ. Более перспективными представляются архитектурные приёмы, среди которых один из наиболее значимых --- конвейеризация. Идея конвейера состоит в разделении функционального блока на несколько последовательных ступеней с буферными регистрами между ними. Данные на вход конвейера могут подаваться с интервалом, равным времени обработки в самой медленной ступени, а последующие результаты появляются на выходе с этим же интервалом. По способу синхронизации конвейеры могут быть синхронными и асинхронными; для традиционных ВМ характерны синхронные конвейеры. Эффективность конвейера оценивается ускорением, эффективностью и производительностью. Теоретически наилучшее время обработки входного потока из \(N\) значений на конвейере с \(K\) ступенями и тактовым периодом \(T_K\) равно \(T_{NK}=(K+(N-1))T_K\). Ускорение \(S=\frac{NK}{K+N-1}\), при \(N\to\infty\) ускорение стремится к \(K\).
\textit{(Источник: Орлов, Цилькер, гл. 9, с. 386--389.)}

\paragraph{Конвейер команд.}
Конвейер команд реализует совмещение во времени этапов цикла команды. Типовые ступени конвейера соответствуют этапам: выборка команды (ВК), декодирование команды (ДК), вычисление исполнительных адресов (ВА), выборка операндов (ВО), исполнение операции (ИО), запись результата (ЗР), формирование адреса следующей команды (ФАСК). Пока одна команда исполняется, следующие могут выбираться и декодироваться. В идеальном случае конвейер из \(K\) ступеней даёт ускорение, стремящееся к \(K\). Однако реальная производительность снижается из-за конфликтов (рисков): структурных, по данным и по управлению.
\textit{(Источник: Орлов, Цилькер, гл. 9, с. 390--394.)}

Структурный риск возникает при одновременном обращении нескольких команд к одному ресурсу, например к памяти. Риск по данным связан с зависимостями команд по операндам: «чтение после записи» (RAW), «запись после чтения» (WAR), «запись после записи» (WAW). Для борьбы с RAW-конфликтами используется ускоренное продвижение данных (forwarding), позволяющее брать результат непосредственно из буферного регистра, минуя промежуточные ступени. Если это невозможно, конвейер приостанавливается. Риск по управлению возникает при выполнении команд перехода: адрес следующей команды становится известен не сразу. Для уменьшения потерь применяются вычисление адреса перехода на ступени декодирования, буфер адресов перехода (BTB), кэш команд точки перехода (BTIC), буфер цикла и предсказание переходов (статическое и динамическое).
\textit{(Источник: Орлов, Цилькер, гл. 9, с. 390--412.)}

\paragraph{Суперскалярная обработка и параллельное выполнение независимых команд.}
Суперскалярный процессор одновременно выполняет более одной скалярной команды. Для этого в его составе имеются несколько конвейеров и несколько функциональных блоков, специализированных по типам операций. Блок диспетчеризации распределяет команды по конвейерам и функциональным блокам. Для повышения параллелизма используются неупорядоченная выдача команд, неупорядоченное завершение, переименование регистров, окно команд (scoreboard или станции резервирования) и буфер восстановления последовательности (ROB). Переименование регистров устраняет ложные зависимости по данным; окно команд и станции резервирования позволяют выдавать команды на исполнение по мере готовности операндов; ROB обеспечивает правильный порядок завершения команд и записи результатов.
\textit{(Источник: Орлов, Цилькер, гл. 9, с. 414--430.)}

Гиперпотоковая обработка (Hyper-Threading) позволяет одному физическому суперскалярному ядру представляться операционной системе как два логических процессора и параллельно обрабатывать два потока команд. Ресурсы ядра при этом разделяются: счётчики команд, регистровые наборы и таблицы распределения регистров дублируются, а функциональные блоки и кэш-память в основном используются совместно. Эффективность гиперпотоковой обработки зависит от характера потоков: она особенно полезна, когда один поток часто простаивает в ожидании памяти или при выполнении команд с разными требованиями к ресурсам.
\textit{(Источник: Орлов, Цилькер, гл. 9, с. 430--434.)}

\paragraph{Арифметический конвейер и векторные команды.}
Арифметический конвейер строится для выполнения отдельных арифметических операций, прежде всего операций с плавающей запятой. Сложная операция разбивается на ряд более простых ступеней, каждая из которых выполняется за один такт; между ступенями размещаются буферные регистры. Пока очередной операнд проходит одну ступень, на предыдущих ступенях обрабатываются следующие операнды. В векторных процессорах арифметические конвейеры применяются к элементам векторов; при наличии нескольких конвейеров возможна параллельная обработка чётных и нечётных элементов вектора. Зацепление векторов (vector chaining) позволяет начать выполнение следующей векторной операции, как только становятся доступны первые результаты предыдущей операции, не дожидаясь завершения всего вектора.
\textit{(Источник: Орлов, Цилькер, гл. 13, с. 527--535; гл. 9, с. 386--389.)}

Векторный процессор оперирует векторами --- одномерными массивами данных, обычно числами с плавающей запятой. Одна векторная команда задаёт операцию сразу над всеми элементами векторов-операндов. В архитектуре «регистр--регистр» векторы предварительно загружаются в векторные регистры; в архитектуре «память--память» элементы векторов выбираются из памяти и результаты сразу записываются в память. Для управления векторными операциями используются регистр длины вектора, регистр максимальной длины, регистр маски и регистр индексов. Регистр маски разрешает или запрещает запись отдельных элементов результата; операции уплотнения/развертывания и сбора/рассеяния (gather/scatter) позволяют работать с нерегулярными элементами. В микропроцессорах векторные и групповые операции реализуются SIMD-командами (MMX, SSE и др.), обрабатывающими упакованные целые числа и числа с плавающей запятой.
\textit{(Источник: Орлов, Цилькер, гл. 13, с. 527--535; гл. 2, с. 93--94.)}

\paragraph{Многоядерные процессоры.}
Многоядерный процессор содержит два или более вычислительных ядра на одном кристалле или в одном корпусе. Под ядром принято понимать процессор и кэш-память (обычно первого уровня). При наличии нескольких ядер необходимо обеспечить возможность их взаимодействия с основной памятью и между собой. Эта задача решается либо путём подключения ядер и памяти к общей шине, либо с помощью коммуникационной сети. Выбор варианта зависит от способа организации основной памяти. Если она является совместно используемой, то применяется вариант с шиной. В случае распределённой основной памяти коммуникации может обеспечить лишь сеть. Шинная организация обусловливает ограничение на число ядер: пропускная способность шины ограничивает число ядер величиной 32, поскольку дальнейшее увеличение количества ядер ведёт к снижению производительности многоядерного процессора. В плане реализации многоядерные архитектуры различных производителей существенно различаются. Это касается не только системы коммуникаций, но и организации иерархии кэш-памяти. Во всех вариантах каждый процессор обладает кэш-памятью первого уровня. Кэш-память второго уровня может быть либо индивидуальной для каждого ядра, либо общей для группы ядер (обычно одна память L2 на пару ядер). В новейших многоядерных процессорах присутствует кэш-память третьего уровня L3, которая обычно является общей для всех ядер. Применение многоядерной технологии позволяет увеличить производительность ВМ и одновременно избежать роста потребления энергии, которое накладывает ограничения на развитие одноядерных процессоров.
\textit{(Источник: Орлов, Цилькер, гл. 9, с. 446--448; гл. 14, с. 557--565.)}

\paragraph{Краткое резюме подраздела.}
Архитектура ВМ охватывает функциональные возможности, видимые программисту, и структурную организацию, реализующую эти возможности. Фон-неймановская модель с хранимой в памяти программой остаётся основой современных компьютеров. Архитектура системы команд определяет состав и сложность команд, место хранения операндов, форматы данных и команд, способы адресации; исторически выделяются CISC-, RISC-, VLIW- и EPIC-подходы. Память организуется иерархически; кэш-память сокращает среднее время доступа за счёт локальности, а виртуальная память обеспечивает иллюзию большого адресного пространства и защиту. Страничная организация использует страницы фиксированного размера и страничные таблицы; сегментная организация --- логические сегменты переменной длины; сегментно-страничная схема сочетает оба подхода. Архитектура процессора включает конвейеризацию команд и арифметических операций, суперскалярную обработку, внеочередное выполнение, переименование регистров, предсказание переходов, гиперпотоковую обработку и многоядерность. Векторные и SIMD-команды обрабатывают группы данных одной командой. Эти принципы обеспечивают высокую производительность современных вычислительных машин.
\textit{(Обобщение по: Орлов, Цилькер, гл. 1--2, 6, 9, 13, 14.)}

\subsection{Страничная и сегментная организация виртуальной памяти}

\noindent\textit{Основные источники: Орлов С.А., Цилькер Б.Я. Организация ЭВМ и систем: Учебник для вузов. 2-е изд. СПб.: Питер, 2011. Гл. 6, с. 288--299.}

\subsubsection{Понятие виртуальной памяти}

\paragraph{Назначение виртуальной памяти.}
В большинстве типичных применений ВМ размещение всей программы в основной памяти невозможно из-за большого размера программы. Решение этой проблемы базируется на свойстве локальности по обращению, согласно которому в каждый момент времени «внимание» машины концентрируется на определённых сравнительно небольших участках программы. Таким образом, в основной памяти достаточно хранить только используемые в данный период фрагменты программ, а остальные их части могут располагаться на внешних запоминающих устройствах (ВЗУ). При таком подходе, однако, необходимо принимать во внимание проблему линейности адресации памяти. Загружаемые в основную память программы или их фрагменты имеют различный размер, но линейную адресацию, то есть должны располагаться в непрерывной области памяти. Для занесения в основную память очередного фрагмента нужно освободить для него непрерывную область, удалив другой фрагмент, но такой, размер которого равен или больше размера загружаемой части программы. Из-за различия в размерах фрагментов, как правило, приходится удалять больший по размеру фрагмент. В итоге свободной остаётся небольшая по размеру область памяти. При многократном повторе подобных действий в основной памяти остаётся множество свободных участков, недостаточных по размеру для загрузки полного непрерывного фрагмента программы. Возникает состояние, известное как \emph{фрагментация памяти}.
\textit{(Источник: Орлов, Цилькер, гл. 6, с. 288--289.)}

Базисом для решения перечисленных проблем служит идея \emph{виртуализации памяти}, под которой понимается такой метод автоматического управления иерархической памятью, при котором программисту кажется, что он имеет дело с единой памятью большой ёмкости и высокого быстродействия. Эту память называют \emph{виртуальной памятью} (ВП). По своей сути виртуализация памяти представляет собой способ аппаратной и программной реализации концепции иерархической памяти.
\textit{(Источник: Орлов, Цилькер, гл. 6, с. 289.)}

В рамках идеи виртуализации памяти основная память рассматривается как линейное пространство \(N\) адресов, называемое \emph{физическим пространством памяти}. Для задач, где требуется более чем \(N\) ячеек, предоставляется значительно большее пространство адресов (обычно равное общей ёмкости всех видов памяти), называемое \emph{виртуальным пространством}, в общем случае не обязательно линейное. Адреса виртуального пространства называют \emph{виртуальными}, а адреса физического пространства --- \emph{физическими}. Программа пишется в виртуальных адресах, но поскольку для её выполнения нужно, чтобы обрабатываемые команды и данные находились в основной памяти, требуется, чтобы каждому виртуальному адресу соответствовал физический. Таким образом, в процессе вычислений необходимо, прежде всего, переписать из ВЗУ в основную память ту часть информации, на которую указывает виртуальный адрес (отобразить виртуальное пространство на физическое), после чего преобразовать виртуальный адрес в физический.
\textit{(Источник: Орлов, Цилькер, гл. 6, с. 289.)}

Среди систем виртуальной памяти можно выделить два класса: системы с фиксированным размером блоков (страничная организация) и системы с переменным размером блоков (сегментная организация). Оба варианта обычно совмещают (сегментно-страничная организация). В современных ВМ наибольшее распространение получила страничная и сегментно-страничная организация.
\textit{(Источник: Орлов, Цилькер, гл. 6, с. 290.)}

\subsubsection{Страничная организация памяти}

\paragraph{Разбиение на страницы.}
Целям преобразования виртуальных адресов в физические служит страничная организация памяти. Её идея состоит в разбиении программы на части равной величины, называемые \emph{страницами}. Размер страницы обычно выбирают в пределах 4--8 Кбайт, но так, чтобы он был кратен ёмкости одного сектора магнитного диска. Виртуальное и физическое адресные пространства разбиваются на блоки размером в страницу. Блок основной памяти, соответствующий странице, часто называют \emph{страничным кадром} или \emph{фреймом} (page frame). Страницам виртуальной и физической памяти присваивают номера.
\textit{(Источник: Орлов, Цилькер, гл. 6, с. 290.)}

Виртуальный адрес в страничной организации состоит из двух частей: номера виртуальной страницы (старшие разряды) и смещения внутри страницы (младшие разряды). Аналогично физический адрес состоит из номера физической страницы и смещения. Поскольку размер страницы фиксирован, смещение в виртуальном и физическом адресах одинаково. Преобразованию подвергается только номер страницы. Процесс доступа к данным по их виртуальному адресу включает следующие шаги:
\begin{enumerate}[label=\arabic*)]
\item Центральный процессор обращается к ячейке, указывая её виртуальный адрес, состоящий из номера виртуальной страницы и смещения относительно её начала.
\item Этот адрес поступает в систему преобразования адресов с целью получения из него физического адреса ячейки в основной памяти.
\item Поскольку смещение в виртуальном и физическом адресах одинаково, преобразованию подвергается лишь номер страницы.
\item Если преобразователь обнаруживает, что нужная физическая страница отсутствует в основной памяти (произошёл промах или страничный сбой), то нужная страница считывается из внешней памяти и заносится в основную память.
\end{enumerate}
\textit{(Источник: Орлов, Цилькер, гл. 6, с. 290.)}

\paragraph{Страничная таблица.}
Преобразователь адресов --- это часть операционной системы, транслирующая номер виртуальной страницы в номер физической страницы, расположенной в основной памяти, а также аппаратура, обеспечивающая этот процесс и позволяющая ускорить его. Преобразование осуществляется с помощью так называемой \emph{страничной таблицы}. При отсутствии нужной страницы в основной памяти преобразователь адресов вырабатывает признак страничного сбоя, по которому операционная система приостанавливает вычисления, пока нужная страница не будет считана из вторичной памяти и помещена в основную.
\textit{(Источник: Орлов, Цилькер, гл. 6, с. 290--291.)}

Виртуальное пространство полностью описывается двумя таблицами: \emph{страничной таблицей} и \emph{картой диска} (будем считать, что вторичная память реализована на магнитных дисках). Таблица страниц определяет, какие виртуальные страницы находятся в основной памяти и в каких физических фреймах, а карта диска содержит информацию о секторах диска, где хранятся виртуальные страницы на диске.
\textit{(Источник: Орлов, Цилькер, гл. 6, с. 291.)}

Число записей в страничной таблице (СТ) в общем случае равно количеству виртуальных страниц. Каждая запись содержит поле номера физической страницы (НФС) и четыре признака:
\begin{itemize}
\item \(V\) (Valid) --- признак присутствия. Устанавливается в единицу, если виртуальная страница в данный момент находится в основной памяти. В этом случае в поле номера физической страницы находится соответствующий номер. Если \(V=0\), то при попытке обратиться к данной виртуальной странице преобразователь адреса генерирует сигнал страничного сбоя (page fault), и операционная система предпринимает действия по загрузке страницы с диска в основную память, обращаясь для этого к карте диска. Загрузка страницы с диска в основную память сопровождается записью в соответствующую строку страничной таблицы (указывается номер физической страницы, куда была загружена виртуальная страница).
\item \(R\) (Referenced) --- признак использования страницы. Устанавливается при обращении к данной странице. Эта информация используется в алгоритме замещения страниц для выбора той из них, которую можно наиболее безболезненно удалить из основной памяти, чтобы освободить место для новой.
\item \(M\) (Modified) --- признак модификации. Поскольку в основной памяти находятся лишь копии страниц, а их оригиналы хранятся на диске, необходимо обеспечить идентичность подлинников и копий. В ходе вычислений содержимое отдельных страниц может изменяться, что фиксируется путём установки в единицу признака модификации \(M\). При удалении страницы из основной памяти проверяется состояние признака \(M\). Если \(M=1\), то перед удалением страницы из основной памяти её необходимо переписать на диск, а при \(M=0\) этого можно не делать.
\item \(A\) (Access) --- признак прав доступа. Служит целям защиты информации и определяет, какой вид доступа к странице разрешён: только для чтения, только для записи, для чтения и для записи.
\end{itemize}
\textit{(Источник: Орлов, Цилькер, гл. 6, с. 291--292.)}

Когда программа загружается в основную память, она может быть направлена в любые свободные в данный момент страничные кадры, независимо от того, расположены ли они подряд или нет. Страничная организация позволяет сократить объём пересылок информации между внешней памятью и основной памятью, так как страницу не нужно загружать до тех пор, пока она действительно не понадобится.
\textit{(Источник: Орлов, Цилькер, гл. 6, с. 292.)}

\paragraph{Буфер быстрого преобразования адреса (TLB).}
Способ реализации страничной таблицы жизненно важен для эффективности техники виртуальной адресации, поскольку каждое обращение к памяти предполагает обращение к страничной таблице. Наиболее быстрый способ --- хранение таблицы в специально выделенных для этого регистрах, но от него приходится отказываться при большом объёме страничной таблицы. Остаётся практически один вариант --- выделение страничной таблице области основной памяти, несмотря на то что это приводит к двукратному увеличению времени доступа к информации. Чтобы сократить это время, в состав ВМ включают дополнительное запоминающее устройство, называемое \emph{буфером быстрого преобразования адреса} (TLB --- Translation Look-aside Buffer), или буфером ассоциативной трансляции, или буфером опережающей выборки, представляющее собой кэш-память.
\textit{(Источник: Орлов, Цилькер, гл. 6, с. 292.)}

При каждом преобразовании номера виртуальной страницы в номер физической страницы результат заносится в TLB: номер физической страницы в память данных, а виртуальной --- в память тегов. Таким образом, в TLB попадают результаты нескольких последних операций трансляции адресов. При каждом обращении к основной памяти преобразователь адресов сначала производит поиск в памяти тегов TLB номера требуемой виртуальной страницы. При попадании адрес соответствующей физической страницы берётся из памяти данных TLB. Если в TLB зафиксирован промах, то процедура преобразования адресов производится с помощью страничной таблицы, после чего осуществляется запись новой пары «номер виртуальной страницы --- номер физической страницы» в TLB.
\textit{(Источник: Орлов, Цилькер, гл. 6, с. 292.)}

Буфер преобразования адресов обычно реализуется в виде полностью ассоциативной или множественно-ассоциативной кэш-памяти с высокой степенью ассоциативности и временем доступа, сопоставимым с аналогичным показателем для кэш-памяти первого уровня (L1). Число входов у типовых TLB невелико (64--256). Так, TLB микропроцессора Pentium III имеет 64 входа при размере страницы 4 Кбайт, что позволяет получить быстрый доступ к адресному пространству в 256 Кбайт.
\textit{(Источник: Орлов, Цилькер, гл. 6, с. 292--293.)}

\paragraph{Многоуровневые и инвертированные страничные таблицы.}
Серьёзной проблемой в системе виртуальной памяти является большой объём страничных таблиц, который пропорционален числу виртуальных страниц. Таблица занимает значительную часть основной памяти, а на поиск уходит много времени, что крайне нежелательно. Один из способов сокращения длины таблиц основан на введении \emph{многоуровневой организации таблиц}. В этом варианте информация оформляется в виде нескольких страничных таблиц сравнительно небольшого объёма, которые образуют второй уровень. Первый уровень представлен таблицей с каталогом, где указано местоположение каждой из страничных таблиц (адрес начала таблицы в памяти) второго уровня. Сначала в каталоге определяется расположение нужной страничной таблицы и лишь затем производится обращение к нужной таблице. О достигаемом эффекте можно судить из следующего примера. Пусть адресная шина ВМ имеет ширину 32 бита, а размер страницы равен 4 Кбайт. Тогда количество виртуальных страниц, а следовательно, и число входов в единой страничной таблице составит \(2^{20}\). При двухуровневой организации можно обойтись одной страницей первого уровня на \(1024\) (\(2^{10}\)) входов и \(1024\) страничными таблицами на такое же число входов.
\textit{(Источник: Орлов, Цилькер, гл. 6, с. 293.)}

Другой подход называют способом \emph{обращённых} или \emph{инвертированных} страничных таблиц. Такую таблицу в каком-то смысле можно рассматривать как увеличенный эквивалент TLB, отличающийся тем, что она хранится в основной памяти и реализуется не аппаратурой, а программными средствами. Число входов в таблицу определяется ёмкостью основной памяти и равно числу страниц, которое может быть размещено в основной памяти. Одновременно с этим имеется и традиционная страничная таблица, но хранится она не в основной памяти, а на диске. Для поиска нужной записи в инвертированной таблице используется хеширование, когда номер записи в таблице вычисляется в соответствии с некой хеш-функцией. Аргументом этой функции служит номер искомой виртуальной страницы. Хеширование позволяет ускорить операцию поиска. Если нужная страница в основной памяти отсутствует, производится обращение к основной таблице на диске, и после загрузки страницы в основную память корректируется и инвертированная таблица.
\textit{(Источник: Орлов, Цилькер, гл. 6, с. 293.)}

\paragraph{Алгоритмы замещения страниц.}
При заполненной основной памяти новая страница может быть загружена только после вытеснения одной из находящихся в ней страниц. Выбор вытесняемой страницы производится по алгоритму замещения. Наиболее распространены следующие алгоритмы:
\begin{itemize}
\item \textbf{FIFO} (First In First Out) --- вытесняется страница, дольше всех находившаяся в основной памяти. Алгоритм прост в реализации (очередь), но не учитывает частоту обращений.
\item \textbf{LRU} (Least Recently Used) --- вытесняется страница, к которой дольше всего не было обращения. Считается достаточно эффективным, но требует аппаратной поддержки (счётчики или очередь ссылок).
\item \textbf{LFU} (Least Frequently Used) --- вытесняется страница, к которой было меньше всего обращений. Требует счётчиков обращений.
\item \textbf{NRU} (Not Recently Used) --- приближённый вариант LRU, использующий биты \(R\) и \(M\): сначала вытесняются страницы с \(R=0,M=0\), затем с \(R=0,M=1\), затем с \(R=1,M=0\) и т.д.
\item Случайный выбор --- страница выбирается случайно; прост, но малоэффективен.
\end{itemize}
Эффективность алгоритма оценивается по числу страничных сбоев: чем их меньше, тем лучше. На практике в современных ОС часто используется усовершенствованный LRU (например, с использованием «часов» --- clock algorithm).
\textit{(Источник: Орлов, Цилькер, гл. 6, с. 291--292; общие сведения об алгоритмах замещения.)}

\paragraph{Защита памяти.}
Страничная организация обеспечивает также защиту памяти. Каждая запись страничной таблицы содержит признак прав доступа \(A\), который определяет, разрешено ли чтение, запись или выполнение для данной страницы. При попытке нарушить права доступа возникает исключение (нарушение защиты). Дополнительно могут использоваться кольца защиты (уровни привилегий) и ключи защиты. Защита на уровне страниц удобна тем, что она не требует проверки границ для каждой ячейки, а применяется к целой странице.
\textit{(Источник: Орлов, Цилькер, гл. 6, с. 292, 295--299.)}

\subsubsection{Сегментная организация памяти}

\paragraph{Понятие сегмента.}
При страничной организации предполагается, что виртуальная память --- это непрерывный массив со сквозной нумерацией слов, что не всегда можно признать оптимальным. Обычно программа состоит из нескольких частей --- кодовой, информационной и стековой. Так как заранее неизвестны длины этих составляющих, то удобно, чтобы при программировании каждая из них имела собственную нумерацию слов, отсчитываемых с нуля. Для этого организуют систему \emph{сегментированной памяти}, выделяя в виртуальном пространстве независимые линейные пространства переменной длины, называемые \emph{сегментами}. Каждый сегмент представляет собой отдельную логическую единицу информации, содержащую совокупность данных или программный код и расположенную в адресном пространстве пользователя. В каждом сегменте устанавливается своя собственная нумерация слов, начиная с нуля. Виртуальная память также разбивается на сегменты, с независимой адресацией слов внутри сегмента. Каждой составляющей программы выделяется сегмент памяти.
\textit{(Источник: Орлов, Цилькер, гл. 6, с. 293--294.)}

Виртуальный адрес при сегментной организации состоит из двух частей: номера сегмента и смещения внутри сегмента. Преобразование виртуального адреса в физический выполняется с помощью \emph{сегментной таблицы}. Каждая запись сегментной таблицы содержит начальный физический адрес сегмента (базу) и его длину (границу). При обращении к памяти проверяется, что смещение не превышает длины сегмента; в противном случае возникает исключение (нарушение границ). Физический адрес вычисляется как сумма базы сегмента и смещения.
\textit{(Источник: Орлов, Цилькер, гл. 6, с. 294.)}

\paragraph{Достоинства и недостатки сегментной организации.}
Сегментная организация обладает следующими достоинствами:
\begin{itemize}
\item логическая структура программы отражается в структуре сегментов, что упрощает программирование, защиту и совместное использование памяти;
\item сегменты могут иметь разную длину, что удобно для размещения кода, данных и стека;
\item каждый сегмент может защищаться независимо (права доступа, режим только чтение и т.п.);
\item возможна разделяемая память: несколько процессов могут отображать один и тот же сегмент в свои адресные пространства.
\end{itemize}
Недостатком сегментной организации является то, что неодинаковый размер сегментов приводит к неэффективному использованию основной памяти. Так, если основная память заполнена, то при замещении одного из сегментов требуется вытеснить такой, размер которого равен или больше размера нового. При многократном повторе подобных действий в основной памяти остаётся множество свободных участков, недостаточных по размеру для загрузки полного сегмента (внешняя фрагментация). Решением проблемы служит сегментно-страничная организация памяти.
\textit{(Источник: Орлов, Цилькер, гл. 6, с. 294.)}

\subsubsection{Сегментно-страничная организация памяти}

\paragraph{Принцип построения.}
В сегментно-страничной организации памяти размер сегмента выбирается не произвольно, а задаётся кратным размеру страницы. Сегмент может содержать то или иное, но обязательно целое число страниц, даже если одна из страниц заполнена частично. Возникает определённая иерархия в организации доступа к данным, состоящая из трёх ступеней: сегмент \(\to\) страница \(\to\) слово. Этой структуре соответствует иерархия таблиц, служащих для перевода виртуальных адресов в физические.
\textit{(Источник: Орлов, Цилькер, гл. 6, с. 294.)}

В сегментной таблице программы перечисляются все сегменты данной программы с указанием начальных адресов страничных таблиц, относящихся к каждому сегменту. Количество страничных таблиц равно числу сегментов, и любая из них определяет расположение каждой из страниц сегмента в памяти, которые могут располагаться не подряд --- часть страниц может находиться в основной памяти, остальные --- во внешней памяти.
\textit{(Источник: Орлов, Цилькер, гл. 6, с. 294--295.)}

Структуру виртуального адреса и процесс преобразования его в физический адрес иллюстрирует рис. 6.38 (см. учебник). Виртуальный адрес содержит номер сегмента, номер страницы внутри сегмента и смещение внутри страницы. Преобразование выполняется в два этапа: сначала по номеру сегмента из сегментной таблицы находится адрес страничной таблицы для данного сегмента; затем по номеру страницы из этой страничной таблицы находится номер физической страницы, к которому добавляется смещение. Для получения физического адреса необходим доступ к сегментной и одной из страничных таблиц, поэтому преобразование адреса может занимать много времени. Для ускорения используются TLB (кэш преобразования адресов) и многоуровневые таблицы.
\textit{(Источник: Орлов, Цилькер, гл. 6, с. 294--295.)}

\paragraph{Защита памяти в сегментно-страничной организации.}
Сегментно-страничная организация позволяет применять защиту на двух уровнях: на уровне сегментов (логическая защита, например запрет записи в кодовый сегмент) и на уровне страниц (права доступа \(A\) в страничной таблице). Дополнительно могут использоваться кольца защиты и ключи. Такая многоуровневая защита обеспечивает гибкое управление доступом и повышает надёжность многозадачных систем.
\textit{(Источник: Орлов, Цилькер, гл. 6, с. 295--299.)}

\subsubsection{Краткое резюме подраздела}
Виртуальная память обеспечивает программисту иллюзию большого непрерывного адресного пространства, автоматически отображая его на физическую память и внешние запоминающие устройства. Страничная организация использует страницы фиксированного размера и страничные таблицы; для ускорения преобразования адресов применяется TLB; для сокращения объёма таблиц --- многоуровневые и инвертированные таблицы. Сегментная организация использует сегменты переменной длины и сегментные таблицы, обеспечивая логическую структуризацию и защиту. Сегментно-страничная организация сочетает оба подхода, разбивая сегменты на страницы и используя иерархию таблиц. Алгоритмы замещения страниц (FIFO, LRU, LFU, NRU) определяют, какая страница будет вытеснена при переполнении. Защита памяти реализуется через признаки прав доступа в таблицах, граничные регистры, ключи и кольца защиты. Виртуальная память является ключевым механизмом современных многозадачных ОС, обеспечивающим эффективное использование памяти и изоляцию процессов.
\textit{(Обобщение по: Орлов, Цилькер, гл. 6, с. 288--299.)}

\subsection{Кэш-память}

\noindent\textit{Основные источники: Орлов С.А., Цилькер Б.Я. Организация ЭВМ и систем: Учебник для вузов. 2-е изд. СПб.: Питер, 2011. Гл. 6, с. 277--288; гл. 11, с. 473--492.}

\subsubsection{Назначение и общая структура кэш-памяти}

\paragraph{Причины введения кэш-памяти.}
В большинстве ВМ основная память строится на базе микросхем динамических ОЗУ, на порядок уступающих по быстродействию центральному процессору. Замена динамических запоминающих устройств на статические ведёт к весьма существенному удорожанию ОП. Экономически приемлемое решение этой проблемы, предложенное М. Уилксом, заключается в том, что между ОП и процессором размещается небольшая, но быстродействующая буферная память, куда в процессе работы копируются те участки ОП, к которым производится обращение со стороны процессора. Выигрыш достигается за счёт ранее рассмотренного свойства локальности: если скопировать содержимое участка ОП в более быстродействующую буферную память и переадресовать на неё все обращения в пределах скопированного участка, то можно существенно сократить среднее время доступа к информации. Упомянутая буферная память получила название \emph{кэш-память} (от английского слова cache --- убежище, тайник), поскольку она обычно скрыта от программиста в том смысле, что он не может её адресовать и может даже вообще не знать о её существовании.
\textit{(Источник: Орлов, Цилькер, гл. 6, с. 277.)}

\paragraph{Структура системы с кэш-памятью.}
В общем виде использование кэш-памяти (КП) поясним следующим образом. Когда ЦП пытается прочитать слово из основной памяти, сначала осуществляется поиск копии этого слова в КП. Если такая копия существует, обращение к ОП не производится, а в ЦП передаётся слово, извлечённое из кэш-памяти. Данную ситуацию принято называть успешным обращением или \emph{попаданием} (hit). При отсутствии слова в кэш-памяти, то есть при неуспешном обращении (\emph{промахе}, miss), блок данных, содержащий это слово, сначала пересылается из ОП в кэш-память, а затем уже из КП затребованное слово передаётся в ЦП.
\textit{(Источник: Орлов, Цилькер, гл. 6, с. 277.)}

ОП ёмкостью \(2^n\) слов при взаимодействии с кэш-памятью рассматривается как состоящая из \(M\) блоков фиксированной длины по \(K\) слов в каждом (\(M=2^n/K\)). Кэш-память также условно разбивается на \(C\) блоков аналогичного размера, причём \(C\ll M\). При обращении к какой-либо ячейке ОП блок, содержащий данную ячейку, копируется в какой-то из блоков кэш-памяти, и все последующие обращения к блоку переадресовываются на его копию в КП. Ввиду малой ёмкости в кэш-памяти в разное время могут находиться копии различных блоков ОП. По этой причине каждую копию необходимо снабдить так называемым \emph{тегом} (признаком), указывающим на то, копия какого блока ОП в настоящий момент хранится в данном блоке КП. В результате кэш-память представляет собой совокупность двух ЗУ, одно из которых предназначено для хранения копий данных из ОП (память данных), а второе --- для хранения тегов этих копий (память тегов). Каждому блоку в памяти данных соответствует одна ячейка в памяти тегов.
\textit{(Источник: Орлов, Цилькер, гл. 6, с. 277--278.)}

\paragraph{Эффективность кэш-памяти.}
Эффективность применения кэш-памяти в иерархической системе памяти зависит от целого ряда факторов, к наиболее существенным из которых следует отнести:
\begin{itemize}
\item ёмкость кэш-памяти;
\item размер блока;
\item способ отображения основной памяти на кэш-память;
\item алгоритм замещения информации в заполненной кэш-памяти;
\item алгоритм согласования содержимого основной и кэш-памяти;
\item число уровней кэш-памяти.
\end{itemize}
\textit{(Источник: Орлов, Цилькер, гл. 6, с. 278.)}

Среднее время доступа \(T_{\text{ср}}\) к системе «кэш + основная память» приближённо оценивается выражением
\[
T_{\text{ср}}=T_{\text{hit}}+(1-H)T_{\text{miss}},
\]
где \(T_{\text{hit}}\) --- время доступа при попадании, \(T_{\text{miss}}\) --- потери на промах (время загрузки блока из ОП и выдачи слова), \(H\) --- коэффициент попаданий. Чем выше \(H\) и меньше \(T_{\text{hit}}\), тем ближе среднее время доступа к времени доступа кэш-памяти.
\textit{(Источник: Орлов, Цилькер, гл. 6, с. 277--280.)}

\subsubsection{Ёмкость кэш-памяти и размер блока}

\paragraph{Выбор ёмкости.}
Выбор ёмкости кэш-памяти --- это всегда определённый компромисс. С одной стороны, кэш-память должна быть достаточно мала, чтобы её стоимостные показатели были близки к величине, характерной для ОП. С другой --- она должна быть достаточно большой, чтобы среднее время доступа в системе, состоящей из основной и кэш-памяти, определялось временем доступа к кэш-памяти. Кроме того, чем вместительнее кэш-память, тем больше логических схем должно участвовать в её адресации. Как следствие, микросхемы повышенной ёмкости работают медленнее по сравнению с микросхемами меньшей ёмкости, даже если они выполнены по одной и той же технологии. Реальная эффективность использования кэш-памяти зависит от характера решаемых задач, и невозможно заранее определить, какая её ёмкость будет действительно оптимальной. По мере увеличения ёмкости кэш-памяти вероятность промахов сначала существенно снижается, но при достижении определённого значения эффект сглаживается и становится несущественным. В настоящее время считается, что для большинства задач близкой к оптимальной является ёмкость кэш-памяти от 1 до 512 Кбайт.
\textit{(Источник: Орлов, Цилькер, гл. 6, с. 278--279.)}

\paragraph{Размер блока.}
Ещё одним важным фактором, влияющим на эффективность использования кэш-памяти, служит размер блока. Когда в кэш-память помещается блок, вместе с требуемым словом туда попадают и соседние слова. По мере увеличения размера блока вероятность промахов сначала падает, так как в кэш, согласно принципу локальности, попадает всё больше данных, которые понадобятся в ближайшее время. Однако когда размер блока становится излишне большим, вероятность промахов начинает расти. Объясняется это тем, что:
\begin{itemize}
\item большие размеры блока уменьшают общее количество блоков, которые можно загрузить в кэш-память, а малое число блоков приводит к необходимости частой их смены;
\item по мере увеличения размера блока каждое дополнительное слово оказывается дальше от запрошенного, поэтому такое дополнительное слово менее вероятно понадобится в ближайшем будущем.
\end{itemize}
Зависимость между размером блока и вероятностью промахов во многом определяется характеристиками конкретной программы, из-за чего трудно рекомендовать определённое значение величины блока. Исследования показывают, что наиболее близким к оптимальному можно признать размер блока, равный 4--8 адресуемым единицам (словам или байтам). На практике его обычно выбирают сравнительно небольшим, равным ширине шины данных, связывающей кэш-память с основной памятью.
\textit{(Источник: Орлов, Цилькер, гл. 6, с. 279.)}

\subsubsection{Способы отображения оперативной памяти на кэш-память}

\paragraph{Постановка задачи отображения.}
Сущность отображения блока основной памяти на кэш-память состоит в его копировании в какой-либо из блоков кэш-памяти, после чего все обращения к блоку в ОП должны переадресовываться на копию в кэш-памяти. Задача отображения сводится к выбору подходящего места в кэш-памяти для очередного копируемого блока ОП. Удачным может быть признан лишь такой способ отображения, который одновременно отвечает трём требованиям: обеспечивает быструю проверку кэш-памяти на наличие в ней копии блока основной памяти; обеспечивает быстрое преобразование адреса блока ОП в адрес блока КП; реализует достижение первых двух требований наиболее экономными средствами. Известные варианты отображения основной памяти на кэш можно свести к трём видам: прямому, полностью ассоциативному и частично-ассоциативному.
\textit{(Источник: Орлов, Цилькер, гл. 6, с. 279--280.)}

\paragraph{Прямое отображение.}
При прямом отображении множество блоков основной памяти условно представляется в виде матрицы, в которой количество строк равно числу блоков в кэш-памяти \(m\). В блок кэш-памяти с номером \(i\) может быть помещён любой блок ОП, но только из \(i\)-й строки матрицы. Номер строки (\(i\)) и столбца (\(k\)) матрицы, на пересечении которых располагается блок основной памяти с адресом \(j\), определяются выражениями \(i=j\bmod m\) и \(k=j\operatorname{div} m\). В двоичной записи адреса блока основной памяти номер строки представлен \(\log_2 m\) младшими разрядами, а номер столбца --- оставшимися старшими разрядами. Такая система позволяет по адресу блока основной памяти, поступившему из ЦП, сразу же однозначно определить адрес блока кэш-памяти, куда должен быть отображён данный блок ОП (если его копия в КП отсутствует) или где следует искать соответствующую копию (если она присутствует в КП). В качестве указателя того, какой из блоков \(i\)-й строки отображён или должен быть отображён на \(i\)-й блок КП (тега), используется номер столбца матрицы, где расположен отображённый (отображаемый) блок ОП.
\textit{(Источник: Орлов, Цилькер, гл. 6, с. 280--281.)}

При обращении к кэш-памяти сначала проверяется наличие в ней копии блока, содержащего искомую ячейку. Для этого из исполнительного адреса ячейки выделяется адрес блока ОП. Младшие разряды адреса указывают на тот единственный блок кэш-памяти \(i\), где может находиться копия, и одновременно на соответствующую ячейку памяти тегов. Тег сравнивается со старшими разрядами поступившего из ЦП адреса. Совпадение означает, что в кэш-памяти находится копия затребованного блока ОП, и процессор адресуется к кэш-памяти. Несовпадение порождает копирование затребованного блока в кэш-память с одновременной записью в память тегов номера колонки, из которой был скопирован блок.
\textit{(Источник: Орлов, Цилькер, гл. 6, с. 281.)}

Прямое отображение --- простой и недорогой в реализации способ отображения. К его достоинствам можно отнести малую разрядность тега и то, что при проверке наличия в КП нужной копии достаточно проверить лишь одну определённую ячейку памяти тегов. Основной его недостаток --- жёсткое закрепление за определёнными блоками ОП одного блока в КП. Поэтому если программа поочерёдно обращается к словам из двух различных блоков, отображаемых на одну и ту же строку кэш-памяти, постоянно станет происходить обновление данной строки и вероятность попадания будет низкой.
\textit{(Источник: Орлов, Цилькер, гл. 6, с. 281.)}

\paragraph{Полностью ассоциативное отображение.}
Полностью ассоциативное отображение позволяет преодолеть недостаток прямого, разрешая загрузку любого блока ОП в любой блок кэш-памяти. В этих условиях тегом служит полный адрес блока в основной памяти, то есть длина тега равна разрядности адреса блока ОП. Для проверки наличия копии блока в кэш-памяти необходимо сравнить теги всех блоков кэш-памяти на совпадение со старшими разрядами адреса, поступившего из ЦП. Такое сравнение желательно выполнять одновременно для всех ячеек памяти тегов. Этому требованию наилучшим образом отвечает ассоциативная память, то есть память тегов должна быть ассоциативной. Полностью ассоциативное отображение обеспечивает гибкость при выборе строки для вновь записываемого блока. Принципиальный недостаток этого способа --- необходимость использования дорогостоящей ассоциативной памяти, причём разрядность ячейки этой памяти достаточно велика.
\textit{(Источник: Орлов, Цилькер, гл. 6, с. 281--282.)}

\paragraph{Частично-ассоциативное отображение.}
Частично-ассоциативное отображение, известное также как множественно-ассоциативное (set-associative), является одним из возможных компромиссов, сочетающим достоинства прямого и полностью ассоциативного способов отображения. В определённой мере оно свободно от их недостатков. Кэш-память разбивается на \(v\) подмножеств (модулей), каждое из которых содержит \(k\) блоков. Иными словами, память тегов и память данных условно представляется матрицами, имеющими \(v\) строк и \(k\) столбцов. Кэш-память с такой организацией принято называть \(k\)-канальной (k-way) или кэш-памятью со степенью ассоциативности \(k\). Основная память также рассматривается как матрица с тем же числом строк \(v\). Зависимость между отдельным модулем и блоками ОП такая же, как и при прямом отображении: на блоки, входящие в модуль \(i\), могут быть отображены только вполне определённые блоки основной памяти, в соответствии с соотношением \(i=j\bmod v\), где \(j\) --- адрес блока ОП. Это означает, что блоки из \(i\)-й строки матрицы блоков ОП могут отображаться лишь на блоки из \(i\)-й строки матрицы блоков кэш-памяти, причём на любой из этих блоков. Так как копия блока ОП может находиться в любом из блоков модуля кэш-памяти, для ускорения поиска нужного блока в пределах модуля используется ассоциативный принцип.
\textit{(Источник: Орлов, Цилькер, гл. 6, с. 282.)}

При обращении к кэш-памяти младшие разряды адреса блока ОП указывают на конкретную ячейку памяти тегов (это соответствует прямому отображению). Далее производится параллельное сравнение каждого из \(k\) тегов, хранящихся в этой ячейке, с полем тега поступившего адреса, то есть поиск нужного тега среди \(k\) возможных осуществляется ассоциативно. В предельных случаях, когда \(v=m, k=1\), частично-ассоциативное отображение сводится к прямому, а при \(v=1, k=m\) --- к полностью ассоциативному. Наиболее общий вид организации частично-ассоциативного отображения --- использование двух блоков на модуль (\(v=m/2, k=2\)). Четырёхканальная частично-ассоциативная кэш-память (\(v=m/4, k=4\)) даёт дополнительное улучшение за сравнительно небольшую дополнительную цену. Дальнейшее увеличение числа блоков в модуле существенного эффекта не даёт. Именно этот способ отображения наиболее широко распространён в современных микропроцессорах. Наряду с 2- и 4-канальной кэш-памятью первого уровня в последних моделях микропроцессоров применяется и 16-канальная организация, в частности при построении кэш-памяти второго уровня.
\textit{(Источник: Орлов, Цилькер, гл. 6, с. 282--283.)}

\subsubsection{Алгоритмы замещения информации в заполненной кэш-памяти}

\paragraph{Необходимость замещения.}
Ввиду малой ёмкости в процессе вычислений все блоки кэш-памяти обычно заняты. Занесение в кэш-память копии ещё одного блока ОП возможно лишь путём замещения содержимого одного из занятых блоков кэш-памяти. В этой ситуации возникает вопрос оптимального выбора замещаемого блока. Естественно, что наилучшим решением была бы замена блока, обращение к которому произойдёт в более отдалённом времени или вообще не случится. При прямом отображении каждому блоку основной памяти соответствует только один определённый блок в кэш-памяти, и никакой иной выбор удаляемого блока здесь невозможен. При полностью и частично ассоциативных способах отображения требуется какой-либо алгоритм замещения (выбора удаляемого из кэш-памяти блока). К сожалению, однозначно предсказать, какой из блоков кэш-памяти в будущем будет наименее востребован, практически нереально, и приходится привлекать алгоритмы, уступающие оптимальному. Среди возможных алгоритмов замещения наиболее распространёнными являются четыре, рассматриваемые в порядке уменьшения их относительной эффективности.
\textit{(Источник: Орлов, Цилькер, гл. 6, с. 284.)}

\paragraph{Алгоритм LRU (Least Recently Used).}
Наиболее эффективным является алгоритм замещения на основе наиболее давнего использования (LRU), при котором замещается тот блок кэш-памяти, к которому дольше всего не было обращения. Проводившиеся исследования показали, что алгоритм LRU, который «смотрит» назад, работает достаточно хорошо в сравнении с оптимальным алгоритмом, «смотрящим» вперёд. Наиболее известны два способа аппаратурной реализации этого алгоритма. В первом из них с каждым блоком кэш-памяти ассоциируют счётчик. К содержимому всех счётчиков через определённые интервалы времени добавляется единица. При обращении к блоку её счётчик обнуляется. Таким образом, наибольшее число будет в счётчике того блока, к которому дольше всего не было обращений, и этот блок --- первый кандидат на замещение. Второй способ реализуется с помощью очереди, куда в порядке заполнения блоков кэш-памяти заносятся ссылки на эти блоки. При каждом обращении к блоку ссылка на него перемещается в конец очереди. В итоге первой в очереди каждый раз оказывается ссылка на блок, к которому дольше всего не было обращений. Именно этот блок прежде всего и замещается.
\textit{(Источник: Орлов, Цилькер, гл. 6, с. 284--285.)}

\paragraph{Алгоритм FIFO (First In First Out).}
Другой возможный алгоритм замещения --- алгоритм, работающий по принципу «первый вошёл, первый вышел» (FIFO). Здесь заменяется блок, дольше всего находившийся в кэш-памяти. Алгоритм легко реализуется с помощью рассмотренной ранее очереди, с той лишь разницей, что после обращения к блоку положение соответствующей ссылки в очереди не меняется.
\textit{(Источник: Орлов, Цилькер, гл. 6, с. 284--285.)}

\paragraph{Алгоритм LFU (Least Frequently Used).}
Ещё один алгоритм --- замена наименее часто использовавшегося блока (LFU). Замещается тот блок в кэш-памяти, к которому было меньше всего обращений. Принцип можно воплотить на практике, связав каждый блок кэш-памяти со счётчиком обращений, к содержимому которого после каждого обращения добавляется единица. Главным претендентом на замещение является блок, счётчик которого содержит наименьшее число.
\textit{(Источник: Орлов, Цилькер, гл. 6, с. 284--285.)}

\paragraph{Случайный выбор.}
Простейший алгоритм --- произвольный выбор блока для замены. Замещаемый блок выбирается случайным образом. Реализовано это может быть, например, с помощью счётчика, содержимое которого увеличивается на единицу с каждым тактовым импульсом, вне зависимости от того, имело место попадание или промах. Значение в счётчике определяет заменяемый блок в полностью ассоциативной кэш-памяти или блок в пределах модуля для частично-ассоциативной кэш-памяти. Данный алгоритм используется крайне редко. Среди известных в настоящее время систем с кэш-памятью наиболее встречаемым является алгоритм LRU.
\textit{(Источник: Орлов, Цилькер, гл. 6, с. 285.)}

\subsubsection{Алгоритмы согласования содержимого кэш-памяти и основной памяти}

\paragraph{Проблема согласования.}
В процессе вычислений ЦП может не только считывать имеющуюся информацию, но и записывать новую, обновляя тем самым содержимое кэш-памяти. С другой стороны, многие устройства ввода/вывода умеют напрямую обмениваться информацией с основной памятью. В обоих вариантах возникает ситуация, когда содержимое блока кэш-памяти и соответствующего блока ОП перестаёт совпадать. В результате на связанное с основной памятью устройство вывода может быть выдана «устаревшая» информация, поскольку все изменения в ней, сделанные процессором, фиксируются только в кэш-памяти, а процессор будет использовать старое содержимое кэш-памяти вместо новых данных, загруженных в ОП из устройства ввода. Для разрешения первой из рассмотренных ситуаций (когда процессор выполняет операцию записи) в системах с кэш-памятью предусмотрены методы обновления основной памяти, которые можно разбить на две большие группы: метод сквозной записи (write through, store through) и метод обратной записи (write back, copy back).
\textit{(Источник: Орлов, Цилькер, гл. 6, с. 285.)}

\paragraph{Метод сквозной записи.}
По методу сквозной записи прежде всего обновляется слово, хранящееся в основной памяти. Если в кэш-памяти существует копия этого слова, то она также обновляется. Если же в кэш-памяти отсутствует нужная копия, то либо из основной памяти в кэш-память пересылается блок, содержащий обновлённое слово (сквозная запись с отображением), либо этого не делается (сквозная запись без отображения). Главное достоинство метода сквозной записи состоит в том, что если блок кэш-памяти освобождается для хранения другого блока, то удаляемый блок можно не возвращать в основную память, поскольку его копия там уже имеется. Метод достаточно прост в реализации. К сожалению, эффект от использования кэш-памяти (сокращение времени доступа) в отношении к операциям записи здесь отсутствует. Определённый выигрыш даёт его модификация, известная как метод буферизированной сквозной записи. Информация сначала записывается в кэш-память и в специальный буфер, работающий по схеме FIFO. Запись в основную память производится уже из буфера, а процессор, не дожидаясь её окончания, может сразу же продолжать свою работу. Конечно, соответствующая логика управления должна заботиться о том, чтобы своевременно «опустошать» заполненный буфер. При использовании буферизации процессор полностью освобождается от работы с ОП.
\textit{(Источник: Орлов, Цилькер, гл. 6, с. 285.)}

\paragraph{Метод обратной записи.}
Согласно методу обратной записи, слово заносится только в кэш-память. Если соответствующего блока в кэш-памяти нет, то нужный блок сначала пересылается из ОП, после чего запись всё равно выполняется исключительно в кэш-память. При замещении блока его необходимо предварительно переслать в соответствующее место основной памяти. Для метода обратной записи, в отличие от алгоритма сквозной записи, характерно то, что при каждом чтении из основной памяти осуществляются две пересылки между основной и кэш-памятью. У рассматриваемого метода есть разновидность --- метод флаговой обратной записи. Когда в каком-то блоке кэш-памяти производится изменение, устанавливается связанный с этим блоком бит изменения (флажок). При замещении блок из кэш-памяти переписывается в ОП только тогда, когда его флажок установлен в 1. Ясно, что эффективность флаговой обратной записи несколько выше. В среднем обратная запись на 10\% эффективнее сквозной записи, но для её реализации требуются и повышенные аппаратные затраты. С другой стороны, практика показывает, что операции записи составляют небольшую долю от общего количества обращений к памяти. Так, в литературе приводится величина 16\%. Другие авторы оценивают долю операций записи величинами в диапазоне от 5 до 34\%. Таким образом, различие по быстродействию между рассмотренными методами невелико.
\textit{(Источник: Орлов, Цилькер, гл. 6, с. 285--286.)}

\paragraph{Согласование при вводе из устройства ввода/вывода.}
Теперь рассмотрим ситуацию, когда в основную память из устройства ввода, минуя процессор, заносится новая информация, и неверной становится копия, хранящаяся в кэш-памяти. Предотвратить подобную несогласованность позволяют два приёма. В первом случае система строится так, чтобы ввод любой информации в ОП автоматически сопровождался соответствующими изменениями в кэш-памяти. Для второго подхода «прямой» доступ к основной памяти допускается только через кэш-память.
\textit{(Источник: Орлов, Цилькер, гл. 6, с. 286.)}

\subsubsection{Смешанная и разделённая кэш-память}

\paragraph{Смешанная кэш-память.}
Когда в микропроцессорах впервые стали применять внутреннюю кэш-память, её обычно использовали как для команд, так и для данных. Такую кэш-память принято называть \emph{смешанной}, а соответствующую архитектуру --- Принстонской (Princeton architecture), по названию университета, где разрабатывались ВМ с единой памятью для команд и данных, то есть соответствующие классической архитектуре фон Неймана. Смешанная кэш-память обладает тем преимуществом, что при заданной ёмкости ей свойственна более высокая вероятность попаданий, так как в ней оптимальный баланс между командами и данными устанавливается автоматически. Так, если в выполняемом фрагменте программы обращения к памяти связаны в основном с выборкой команд, а доля обращений к данным относительно мала, кэш-память имеет тенденцию насыщаться командами, и наоборот.
\textit{(Источник: Орлов, Цилькер, гл. 6, с. 286.)}

\paragraph{Разделённая кэш-память.}
Сравнительно недавно стало обычным разделять кэш-память на две --- отдельно для команд и отдельно для данных. Подобная архитектура получила название Гарвардской (Harvard architecture), поскольку именно в Гарвардском университете был создан компьютер «Марк-1» (1950 год), имевший раздельные ЗУ для команд и данных. При раздельной кэш-памяти выборка команд и данных может производиться одновременно, при этом исключаются возможные конфликты. Последнее обстоятельство существенно в системах, использующих конвейеризацию команд, где процессор извлекает команды с опережением и заполняет ими буфер или конвейер. Раздельная схема обычно характерна для кэш-памяти первого уровня, причём кэш-память команд и кэш-память данных могут быть организованы по-разному, например в них могут отличаться методы отображения или замещения информации. Следует добавить, что в некоторых ВМ, помимо кэш-памяти команд и кэш-памяти данных, может использоваться и адресная кэш-память (в устройствах управления памятью и при преобразовании виртуальных адресов в физические).
\textit{(Источник: Орлов, Цилькер, гл. 6, с. 286--287.)}

\subsubsection{Одноуровневая и многоуровневая кэш-память}

\paragraph{Многоуровневая организация.}
Логично предположить, что с увеличением ёмкости кэш-памяти можно ожидать и соответствующего повышения быстродействия ВМ. Современные технологии позволяют разместить на общем с процессором кристалле кэш-память достаточно большой ёмкости. С другой стороны, попытки увеличения ёмкости обычно приводят к снижению быстродействия, главным образом из-за усложнения схем управления и дешифрации адреса. По этой причине общую ёмкость кэш-памяти ВМ увеличивают за счёт иерархической организации, при которой кэш-память состоит из нескольких уровней запоминающих устройств, отличающихся ёмкостью и быстродействием. Каждый последующий уровень имеет большую ёмкость по сравнению с предыдущим, но и меньшее быстродействие.
\textit{(Источник: Орлов, Цилькер, гл. 6, с. 287.)}

Первый уровень (L1) иерархии образует наиболее скоростная кэш-память сравнительно небольшой ёмкости (не более 128 Кбайт). На кристалле её располагают по возможности ближе к процессору, чтобы минимизировать длину соединяющей их шины и тем самым способствовать ускорению обмена информацией. Этот уровень в перспективных микропроцессорах обычно строится по разделённой схеме. Так, в многоядерном процессоре Nehalem фирмы Intel каждому ядру придаётся кэш-память команд ёмкостью 32 Кбайт (4-входовая, модульно-ассоциативная) и кэш-память данных ёмкостью 32 Кбайт (8-входовая, модульно-ассоциативная). В многоядерном микропроцессоре Phenom фирмы AMD каждое ядро оснащено кэш-памятью команд и кэш-памятью данных ёмкостью по 64 Кбайт каждая. Ёмкость кэш-памяти второго уровня (L2) обычно больше, чем у L1, а быстродействие и стоимость --- несколько ниже. В упоминавшихся микропроцессорах каждому ядру придаётся кэш-память L2, но уже смешанного типа. В Nehalem она составляет 256 Кбайт, а в Phenom --- 512 Кбайт. В современных микропроцессорах кэш-память второго уровня также размещают на одном кристалле с процессором, за счёт чего сокращается длина связей и повышается быстродействие. Для новейших многоядерных микропроцессоров типично размещение на том же кристалле ещё и кэш-памяти третьего уровня (L3). Обычно этот уровень является общим для всех ядер (или группы ядер). В Nehalem, где количество ядер в зависимости от модели может меняться от 2 до 8, ёмкость L3 составляет от 2 до 8 Мбайт, а в Phenom --- 2 Мбайт.
\textit{(Источник: Орлов, Цилькер, гл. 6, с. 287.)}

При доступе к памяти ЦП сначала обращается к кэш-памяти первого уровня. В случае промаха производится обращение к кэш-памяти второго уровня и т. д. (вплоть до основной памяти). Среднее время доступа \(T_{\text{av}}\) к одноуровневой кэш-памяти можно оценить как
\[
T_{\text{av}}=T_{L1h}+K_{L1m}T_{L1m},
\]
где \(T_{L1h}\) --- время обращения при попадании; \(T_{L1m}\) --- потери на промах; \(K_{L1m}\) --- коэффициент промахов. Для двухуровневой кэш-памяти имеем:
\[
T_{\text{av}}=T_{L1h}+(K_{L1m}T_{L2h})+(K_{L2m}T_{L2m}).
\]
Потенциальная экономия за счёт применения L2 зависит от вероятности попаданий как в L1, так и в L2. Ряд исследований показывает, что использование кэш-памяти второго уровня существенно улучшает производительность.
\textit{(Источник: Орлов, Цилькер, гл. 6, с. 288.)}

\paragraph{Инклюзивная и эксклюзивная иерархия.}
При построении иерархии кэш-памяти используется один из двух принципов --- инклюзивный или эксклюзивный. При инклюзивном подходе вся информация из верхних уровней иерархии содержится и в последующих нижних уровнях. Это означает, что, например, при замещении блока в L1 удаляемый блок переносится в L2. Так организована кэш-память L3 в микропроцессоре Nehalem. В эксклюзивной схеме блок, удаляемый из кэш-памяти более высокого уровня, не дублируется в памяти следующего уровня. Вместо этого замещаемый и замещающий блоки взаимно меняются местами. Например, место блока из L2, замещающего какой-то блок в L1, занимает блок, удалённый при этом из L1. По такой схеме происходит взаимодействие уровня L3 с уровнем L2 в микропроцессоре Phenom. Каждый из вариантов имеет свои достоинства. Так, инклюзивная схема потенциально обеспечивает более высокое быстродействие, а эксклюзивная --- более эффективное использование ёмкости второго и последующих уровней кэш-памяти. На практике же сделанные утверждения достаточно условны, и однозначно отдать предпочтение той или иной схеме трудно.
\textit{(Источник: Орлов, Цилькер, гл. 6, с. 288.)}

\subsubsection{Когерентность кэш-памяти в многопроцессорных системах}

\paragraph{Проблема когерентности.}
Мультипроцессорная система с разделяемой памятью состоит из двух или более независимых процессоров, каждый из которых выполняет либо часть большой программы, либо независимую программу. Все процессоры обращаются к командам и данным, хранящимся в разделяемой основной памяти. Поскольку память является общим ресурсом, при обращении к ней между процессорами возникает соперничество, в результате чего средняя задержка на доступ к памяти увеличивается. Для сокращения такой задержки каждому процессору придаётся локальная кэш-память, которая, обслуживая локальные обращения к памяти, во многих случаях предотвращает необходимость доступа к разделяемой основной памяти. В свою очередь, оснащение каждого процессора локальной кэш-памятью приводит к так называемой проблеме когерентности или обеспечения согласованности кэш-памяти. Согласно литературе, система является когерентной, если каждая операция чтения по какому-либо адресу, выполненная любым из процессоров, возвращает значение, занесённое в ходе последней операции записи по этому адресу, вне зависимости от того, какой из процессоров производил запись последним.
\textit{(Источник: Орлов, Цилькер, гл. 11, с. 473--474.)}

В простейшей форме проблему когерентности кэш-памяти можно пояснить следующим образом. Пусть два процессора ПР1 и ПР2 связаны с общей памятью посредством шины. Сначала оба процессора читают переменную \(x\). Копии блоков, содержащих эту переменную, пересылаются из основной памяти в локальные кэши обоих процессоров КП1 и КП2. Далее процессор ПР1 выполняет операцию увеличения значения переменной \(x\) на единицу. Так как копия переменной уже находится в кэш-памяти данного процессора, произойдёт кэш-попадание и значение \(x\) будет изменено только в КП1. Если теперь процессор ПР2 вновь выполнит операцию чтения \(x\), то также произойдёт кэш-попадание и ПР2 получит хранящееся в его кэш-памяти «старое» значение \(x\). Поддержание согласованности требует, чтобы при изменении элемента данных одним из процессоров соответствующие изменения были проведены в кэш-памяти остальных процессоров, где есть копия изменённого элемента данных, а также в общей памяти. Схожая проблема возникает и в однопроцессорных системах, где присутствует несколько уровней кэш-памяти. Здесь требуется согласовать содержимое кэшей разных уровней.
\textit{(Источник: Орлов, Цилькер, гл. 11, с. 474.)}

\paragraph{Аппаратные протоколы когерентности.}
Большинство из предложенных способов борьбы с некогерентностью ориентированы на динамическое (в процессе вычислений) распознавание и устранение несогласованности копий совместно используемых данных с помощью специальной аппаратуры. Аппаратные методы обеспечивают более высокую производительность, поскольку издержки, связанные с некогерентностью, имеют место только при возникновении ситуации некогерентности. Кроме того, этот подход прозрачен для программиста и пользователя. Аппаратные механизмы преодоления проблемы когерентности принято называть \emph{протоколами когерентности кэш-памяти}. В общем случае для поддержания когерентности в мультипроцессорных системах имеются следующие возможности: разделяемая кэш-память; некэшируемые данные; широковещательная запись; протоколы наблюдения; протоколы на основе справочника.
\textit{(Источник: Орлов, Цилькер, гл. 11, с. 475--477.)}

\paragraph{Протоколы наблюдения.}
В протоколах наблюдения (snoopy protocols) ответственность за поддержание когерентности всех кэшей многопроцессорной системы возлагается на контроллеры кэшей. В системах, где реализованы протоколы наблюдения, контроллер каждой локальной кэш-памяти содержит блок слежения за шиной, который следит за всеми транзакциями на общей шине и, в частности, контролирует все операции записи. Процессоры должны широковещательно передавать на шину любые запросы на доступ к памяти, потенциально способные изменить состояние когерентности совместно используемых блоков данных. Локальный контроллер кэш-памяти каждого процессора затем определяет, присутствует ли в его кэш-памяти копия модифицируемого блока, и если это так, то такой блок аннулируется или обновляется. Протоколы наблюдения характерны для мультипроцессорных систем на базе шины, поскольку общая шина достаточно просто обеспечивает как наблюдение, так и широковещательную передачу сообщений. Однако здесь необходимо принимать меры, чтобы повышенная нагрузка на шину, связанная с наблюдением и трансляцией сообщений, не «съела» преимуществ локальных кэшей. В большинстве протоколов стратегия обеспечения когерентности кэш-памяти расценивается как смена состояний в конечном автомате. При таком подходе предполагается, что любой блок в локальной кэш-памяти может находиться в одном из фиксированных состояний. Обычно число таких состояний не превышает четырёх, поэтому в теге каждого блока кэш-памяти имеются два бита, называемые битами состояния (SB, Status Bit). Среди известных протоколов наблюдения наиболее популярным является протокол MESI (Modified/Exclusive/Shared/Invalid).
\textit{(Источник: Орлов, Цилькер, гл. 11, с. 478--484.)}

\paragraph{Протокол MESI.}
Протокол MESI широко распространён в коммерческих микропроцессорных системах, например на базе микропроцессоров Pentium и PowerPC. Протокол был разработан для кэш-памяти с обратной записью. Одной из основных задач протокола MESI является откладывание на максимальный срок операции обратной записи кэшированных данных в основную память ВС. Это позволяет улучшить производительность системы за счёт минимизации ненужных пересылок информации между кэшами и основной памятью. Протокол MESI приписывает каждому блоку кэш-памяти одно из четырёх состояний, которые контролируются двумя битами состояния MESI в теге данного блока:
\begin{itemize}
\item \textbf{M} (Modified) --- данные в блоке кэш-памяти были модифицированы, но изменённая информация пока не переписана в основную память. Информация, содержащаяся в рассматриваемом блоке, достоверна только в данном кэше, а в основной памяти и остальных кэшах --- недостоверна.
\item \textbf{E} (Exclusive) --- данный блок в кэш-памяти не подвергался изменению посредством запроса на запись, совпадает с аналогичным блоком в основной памяти, но отсутствует в любом другом локальном кэше. Информация достоверна в этой локальной кэш-памяти и недостоверна в любой другой.
\item \textbf{S} (Shared) --- блок в кэш-памяти совпадает с аналогичным блоком в основной памяти (данные достоверны) и может присутствовать в одном или нескольких из прочих кэшей.
\item \textbf{I} (Invalid) --- блок кэша, помеченный как недействительный, не содержит достоверных данных и становится логически недоступным.
\end{itemize}
Порядок перехода блока кэш-памяти из одного состояния в другое зависит от текущего статуса блока, выполняемой операции (чтение или запись), результата обращения в кэш (попадание или промах) и от того, является ли блок разделяемым или нет. Протокол MESI обеспечивает корректную работу с разделяемыми данными и является одним из наиболее эффективных протоколов наблюдения.
\textit{(Источник: Орлов, Цилькер, гл. 11, с. 484--488.)}

\paragraph{Протоколы на основе справочника.}
Протоколы обеспечения когерентности на основе справочника характерны для сложных мультипроцессорных систем с совместно используемой памятью, где процессоры объединены многоступенчатой иерархической коммуникационной сетью. Сложность конфигурации сети приводит к тому, что применение протоколов наблюдения с их механизмом широковещания становится дорогостоящим и неэффективным. Протоколы на основе справочника предполагают сбор и отслеживание информации о содержимом всех локальных кэшей. Такие протоколы обычно реализуются с помощью централизованного контроллера, физически представляющего собой часть контроллера основной памяти. Собственно справочник хранится в основной памяти. Когда контроллер локальной кэш-памяти делает запрос, контроллер справочника обнаруживает такой запрос и формирует команды, необходимые для пересылки данных из основной памяти либо из другой локальной кэш-памяти, содержащей последнюю версию запрошенных данных. Центральный контроллер отвечает за обновление информации о состоянии локальных кэшей, поэтому он должен быть извещён о любом локальном действии, способном повлиять на состояние блока данных. Известны три способа реализации протоколов на основе справочника: полный справочник, ограниченные справочники и сцепленные справочники.
\textit{(Источник: Орлов, Цилькер, гл. 11, с. 488--492.)}

\subsubsection{Краткое резюме подраздела}
Кэш-память является ключевым элементом иерархии памяти, сокращающим среднее время доступа за счёт использования свойства локальности. Она состоит из памяти данных и памяти тегов; эффективность кэша зависит от ёмкости, размера блока, способа отображения, алгоритмов замещения и согласования. Основные способы отображения --- прямое, полностью ассоциативное и частично-ассоциативное. Алгоритмы замещения (LRU, FIFO, LFU, случайный) определяют, какой блок вытесняется при заполнении. Согласование содержимого кэша и основной памяти выполняется методами сквозной и обратной записи. В многопроцессорных системах для поддержания когерентности используются протоколы наблюдения (например, MESI) и протоколы на основе справочника. Многоуровневая кэш-память (L1, L2, L3) позволяет сочетать малое время доступа с большой ёмкостью.
\textit{(Обобщение по: Орлов, Цилькер, гл. 6, с. 277--288; гл. 11, с. 473--492.)}

\subsection{Командный и арифметический конвейеры, параллельное выполнение независимых команд, векторные команды}

\noindent\textit{Основные источники: Орлов С.А., Цилькер Б.Я. Организация ЭВМ и систем: Учебник для вузов. 2-е изд. СПб.: Питер, 2011. Гл. 9, с. 386--449; гл. 13, с. 526--535; гл. 2, с. 93--94.}

\subsubsection{Командный конвейер}

\paragraph{Идея конвейеризации вычислений.}
Совершенствование элементной базы уже не приводит к кардинальному росту производительности ВМ. Более перспективными в этом плане представляются архитектурные приёмы, среди которых один из наиболее значимых --- конвейеризация. Идея конвейера состоит в разделении функционального блока на несколько последовательных ступеней с буферными регистрами между ними. Данные на вход конвейера могут подаваться с интервалом, равным времени обработки в самой медленной ступени, а последующие результаты появляются на выходе с этим же интервалом. По способу синхронизации конвейеры могут быть синхронными и асинхронными; для традиционных ВМ характерны синхронные конвейеры. Эффективность конвейера оценивается ускорением, эффективностью и производительностью. Теоретически наилучшее время обработки входного потока из \(N\) значений на конвейере с \(K\) ступенями и тактовым периодом \(T_K\) равно \(T_{NK}=(K+(N-1))T_K\). Ускорение \(S=\frac{NK}{K+N-1}\), при \(N\to\infty\) ускорение стремится к \(K\).
\textit{(Источник: Орлов, Цилькер, гл. 9, с. 386--389.)}

\paragraph{Структура конвейера команд.}
В архитектуре вычислительных машин можно найти множество объектов, где конвейеризация обеспечивает существенный прирост производительности ВМ, например операционные устройства и память, однако наиболее ощутимый эффект достигается при конвейеризации этапов машинного цикла. Как известно, цикл команды представляет собой последовательность этапов. Возложив реализацию каждого из них на самостоятельное устройство и последовательно соединив такие устройства, мы получим классическую схему конвейера команд. На рис. 9.3 (см. учебник) показан конвейер с семью ступенями, соответствующими семи этапам цикла команды:
\begin{itemize}
\item выборка команды (ВК);
\item декодирование команды (ДК);
\item вычисление исполнительных адресов (ВА);
\item выборка операндов (ВО);
\item исполнение операции (ИО);
\item запись результата (ЗР);
\item формирование адреса следующей команды (ФАСК).
\end{itemize}
В диаграмме предполагается, что каждая команда обязательно проходит все семь ступеней, хотя этот случай не совсем типичен. Так, команда загрузки регистра не требует этапа ЗР. Кроме того, здесь принято, что все этапы могут выполняться одновременно. Без конвейеризации выполнение девяти команд заняло бы \(9\times7=63\) единицы времени. Использование конвейера позволяет сократить время обработки до 15 единиц.
\textit{(Источник: Орлов, Цилькер, гл. 9, с. 390.)}

\paragraph{Конфликты в конвейере команд.}
Полученное в примере число 15 характеризует лишь потенциальную производительность конвейера команд. На практике, в силу возникающих в конвейере конфликтных ситуаций, достичь такой производительности не удаётся. Конфликтные ситуации в конвейере принято обозначать термином \emph{риск} (hazard), а обусловлены они могут быть тремя причинами:
\begin{enumerate}[label=\arabic*)]
\item попыткой нескольких команд одновременно обратиться к одному и тому же ресурсу ВМ (\emph{структурный риск});
\item взаимосвязью команд по данным (\emph{риск по данным});
\item неоднозначностью при выборке следующей команды в случае команд, изменяющих последовательность вычислений (\emph{риск по управлению}).
\end{enumerate}
\textit{(Источник: Орлов, Цилькер, гл. 9, с. 390--391.)}

\paragraph{Структурный риск.}
Структурный риск (конфликт по ресурсам) имеет место, когда несколько команд, находящихся на разных ступенях конвейера, пытаются одновременно использовать один и тот же ресурс, чаще всего --- память. Так, в типовом цикле команды сразу три этапа (ВК, ВО и ЗР) связаны с обращением к памяти. Подобных конфликтов частично удаётся избежать за счёт блочной организации основной памяти и использования кэш-памяти --- имеется вероятность того, что команды будут обращаться либо к разным банкам ОП, либо одна из них станет обращаться к основной памяти, а другая --- к кэш-памяти. С этих позиций кэш-память выгоднее разделять на кэш-память команд и кэш-память данных. Конфликты из-за одновременного обращения к памяти могут и не возникать, поскольку для многих команд этапы выборки операнда и записи результата не требуются. В целом, влияние структурного риска на производительность конвейера по сравнению с другими видами рисков сравнительно невелико.
\textit{(Источник: Орлов, Цилькер, гл. 9, с. 391.)}

\paragraph{Риск по данным.}
Риск по данным, в противоположность структурному риску, --- типичная и регулярно возникающая ситуация. Для пояснения сущности взаимосвязи команд по данным положим, что две команды в конвейере (\(i\) и \(j\)) предусматривают обращение к одной и той же переменной \(x\), причём команда \(i\) предшествует команде \(j\). В общем случае между \(i\) и \(j\) ожидаются три типа конфликтов по данным:
\begin{itemize}
\item \textbf{«чтение после записи» (ЧПЗ, RAW)}: команда \(j\) читает \(x\) до того, как команда \(i\) успела записать новое значение \(x\), то есть команда \(j\) ошибочно получит старое значение \(x\) вместо нового;
\item \textbf{«запись после чтения» (ЗПЧ, WAR)}: команда \(j\) записывает новое значение \(x\) до того, как команда \(i\) успела прочитать \(x\), то есть команда \(i\) ошибочно получит новое значение \(x\) вместо старого;
\item \textbf{«запись после записи» (ЗПЗ, WAW)}: команда \(j\) записывает новое значение \(x\) прежде, чем команда \(i\) успела записать в качестве \(x\) своё значение, и окончательно \(x\) ошибочно будет содержать значение, определённое командой \(i\), а не командой \(j\).
\end{itemize}
Возможен и четвёртый случай, когда команда \(j\) читает \(x\) прежде команды \(i\). Этот случай не вызывает никаких конфликтов, поскольку как \(i\), так и \(j\) получат верное значение \(x\).
\textit{(Источник: Орлов, Цилькер, гл. 9, с. 391--392.)}

Наиболее частый вид конфликтов по данным --- ЧПЗ, поскольку операция чтения в цикле команды (этап ВО) предшествует операции записи (этап ЗР). По той же причине конфликты типа ЗПЧ большой проблемы не представляют. Сложности появляются, только если структура конвейера допускает запись прежде чтения или если команды в конвейере обрабатываются в последовательности, отличной от предписанной программой. Такое возможно, если командам в конвейере разрешается «догонять» предшествующие им команды, приостановленные из-за какого-то конфликта. Конфликт типа ЗПЗ также не вызывает особых проблем в конвейерах, где команды следуют в порядке, определённом программой, и могут производить запись только на этапе ЗР. В худшем случае, когда одна команда догоняет другую из-за приостановки последней, имеет место конфликт по ресурсу --- попытка одновременного доступа к одной и той же ячейке.
\textit{(Источник: Орлов, Цилькер, гл. 9, с. 392.)}

В борьбе с конфликтами по данным выделяют два аспекта: своевременное обнаружение потенциального конфликта и его устранение. Признаком возникновения конфликта по данным между двумя командами \(i\) и \(j\) служит невыполнение хотя бы одного из трёх условий Бернстейна (Bernstein's Conditions):
\begin{itemize}
\item для ЧПЗ: \(O(i)\cap I(j)=\varnothing\);
\item для ЗПЧ: \(I(i)\cap O(j)=\varnothing\);
\item для ЗПЗ: \(O(i)\cap O(j)=\varnothing\),
\end{itemize}
где \(O(k)\) --- множество ячеек, изменяемых командой \(k\); \(I(l)\) --- множество ячеек, читаемых командой \(l\); \(\varnothing\) --- пустое множество; \(\cap\) --- операция пересечения множеств. Условия можно распространить и на большее число команд: для трёх команд подобных уравнений будет 9, для четырёх команд --- 18 (по три на каждую пару). Соблюдение соотношений является достаточным, но не необходимым условием, поскольку в ряде случаев конфликты могут и не возникать.
\textit{(Источник: Орлов, Цилькер, гл. 9, с. 392.)}

Для борьбы с конфликтами по данным применяются как программные, так и аппаратные методы. Программные методы ориентированы на устранение самой возможности конфликтов ещё на стадии компиляции программы. Оптимизирующий компилятор пытается создать такой объектный код, чтобы между командами, склонными к конфликтам, находилось достаточное количество нейтральных в этом плане команд программы. Если такое не удаётся, то между конфликтующими командами компилятор вставляет необходимое количество команд типа «Нет операции». Фактическое разрешение конфликтов возлагается на аппаратные методы. Наиболее очевидным решением является остановка команды \(j\) на несколько тактов с тем, чтобы команда \(i\) успела завершиться или, по крайней мере, миновать ступень конвейера, на которой возникает конфликт. Соответственно задерживаются и команды, следующие в конвейере за \(j\)-й командой. Данную ситуацию называют «пузырьком» в конвейере. Иногда приостанавливают только команду \(j\), не задерживая следующие за ней команды. Это более эффективный приём, но его реализация усложняет конвейер. Поскольку наиболее частые конфликты по данным --- это ЧПЗ, основные усилия тратятся на противодействие именно этому типу конфликтов. Среди известных методов борьбы с ЧПЗ наибольшее распространение получил приём \emph{ускоренного продвижения информации} (forwarding). Обычно между двумя соседними ступенями конвейера располагается буферный регистр, через который предшествующая ступень передаёт результат своей работы на последующую ступень, то есть передача информации возможна лишь между соседними ступенями конвейера. При ускоренном продвижении, когда для выполнения команды требуется операнд, уже вычисленный предыдущей командой, этот операнд может быть получен непосредственно из соответствующего буферного регистра, минуя все промежуточные ступени конвейера. С данной целью в конвейере предусматриваются дополнительные тракты пересылки информации (тракты опережения, тракты обхода), снабжённые средствами коммутации трактов.
\textit{(Источник: Орлов, Цилькер, гл. 9, с. 392--393.)}

\paragraph{Риск по управлению и предсказание переходов.}
Наибольшие проблемы при создании эффективного конвейера обусловлены командами, изменяющими естественный порядок вычислений. Простейший конвейер ориентирован на линейные программы. В нём ступень выборки извлекает команды из ячеек памяти с последовательными адресами, используя для этого счётчик команд (СК). Адрес очередной команды в линейной программе формируется автоматически, за счёт прибавления к содержимому СК числа, равного длине текущей команды в байтах. Реальные программы практически никогда не бывают линейными. В них обязательно присутствуют команды управления, изменяющие последовательность вычислений: безусловный и условный переход, вызов процедуры, возврат из процедуры и т.п. Доля подобных команд в программе оценивается как 10--20\%. Выполнение команд, изменяющих последовательность вычислений (в дальнейшем будем их называть командами перехода), может приводить к приостановке конвейера на несколько тактов, из-за чего производительность процессора снижается. Приостановки конвейера при выполнении команд перехода обусловлены двумя факторами.
\textit{(Источник: Орлов, Цилькер, гл. 9, с. 393--394.)}

Первый фактор характерен для любой команды перехода и связан с выборкой команды из точки перехода (по адресу, указанному в команде перехода). То, что текущая команда относится к командам перехода, становится ясным только после прохождения ступени декодирования. За это время на первые ступени конвейера уже поступят новые команды, извлечённые в предположении, что естественный порядок вычислений не будет нарушен. В случае перехода эти начальные ступени нужно очистить и загрузить в конвейер команду, расположенную по адресу перехода. Вторая причина нарушения ритмичности работы конвейера имеет отношение только к командам условного перехода: до завершения команды невозможно определить, какая из команд должна выполняться следующей. Для сокращения задержек, обусловленных выборкой команды из точки перехода, применяются следующие подходы:
\begin{itemize}
\item вычисление исполнительного адреса перехода на ступени декодирования команды;
\item использование буфера адресов перехода (BTB, Branch Target Buffer);
\item использование кэш-памяти для хранения команд, расположенных в точке перехода (BTIC);
\item использование буфера цикла.
\end{itemize}
\textit{(Источник: Орлов, Цилькер, гл. 9, с. 394--396.)}

Для устранения или частичного сокращения издержек от условных переходов применяется \emph{предсказание переходов}. Идея заключается в том, что ещё до момента выполнения команды условного перехода делается предположение о наиболее вероятном исходе такой команды. Последующие команды подаются на конвейер в соответствии с предсказанием. При ошибочном предсказании конвейер необходимо вернуть к состоянию, с которого началась выборка «ненужных» команд, что по эффекту эквивалентно приостановке конвейера. При правильных предсказаниях конвейер функционирует ритмично без остановок. Различают статическое и динамическое предсказание.
\textit{(Источник: Орлов, Цилькер, гл. 9, с. 396--397.)}

При \emph{статическом} предсказании решение принимается на этапе компиляции и не меняется в процессе вычислений. Основные стратегии: переход происходит всегда (ПВ); переход не происходит никогда (ПН); предсказание по результатам профилирования; предсказание по коду операции; предсказание по направлению перехода (назад --- переход есть, вперёд --- перехода нет); при первом выполнении переход всегда происходит. Например, для команд, выполняющих переход «назад» (типичных для циклов), фактический переход имеет место в 85\% случаев. Для переходов «вперёд» эта цифра составляет 65\%.
\textit{(Источник: Орлов, Цилькер, гл. 9, с. 397--399.)}

В \emph{динамических} стратегиях решение принимается в ходе вычислений на основе информации о предшествующих переходах (истории переходов). История фиксируется в таблице предыстории переходов (PHT, Pattern History Table), каждый элемент которой состоит из \(m\) битов. Нужный элемент выбирается с помощью \(k\)-разрядного шаблона. В качестве шаблона могут использоваться адрес команды перехода, содержимое регистра глобальной истории (GHR) или регистра локальной истории (LHR). Наибольшее распространение получили двухразрядные счётчики (автомат A3), работающие в режиме с насыщением (алгоритм Смита). В этом случае каждому исходу соответствуют четыре состояния: 00 --- перехода не будет (сильное предсказание), 01 --- перехода не будет (слабое), 10 --- переход будет (слабое), 11 --- переход будет (сильное). Основанием для предсказания служит состояние старшего разряда счётчика. Существуют также гибридные и асимметричные схемы, комбинирующие несколько предикторов.
\textit{(Источник: Орлов, Цилькер, гл. 9, с. 399--412.)}

\subsubsection{Арифметический конвейер}

\paragraph{Принцип построения.}
Арифметический конвейер строится для выполнения отдельных арифметических операций, прежде всего операций с плавающей запятой. Сложная операция (например, сложение или умножение чисел с плавающей запятой) разбивается на ряд более простых ступеней, каждая из которых выполняется за один такт; между ступенями размещаются буферные регистры. Пока очередной операнд проходит одну ступень, на предыдущих ступенях обрабатываются следующие операнды. В векторных процессорах арифметические конвейеры применяются к элементам векторов; при наличии нескольких конвейеров возможна параллельная обработка чётных и нечётных элементов вектора. Зацепление векторов (vector chaining) позволяет начать выполнение следующей векторной операции, как только становятся доступны первые результаты предыдущей операции, не дожидаясь завершения всего вектора. В современных процессорах арифметические конвейеры для операций с плавающей запятой обычно имеют глубину 4--8 ступеней; для целочисленных операций --- 2--4 ступени.
\textit{(Источник: Орлов, Цилькер, гл. 9, с. 386--389; гл. 13, с. 533--535.)}

\subsubsection{Параллельное выполнение независимых команд}

\paragraph{Суперскалярная обработка.}
Суперскалярным называется центральный процессор, который одновременно выполняет более одной скалярной команды. Это достигается за счёт включения в состав ЦП нескольких самостоятельных функциональных блоков, каждый из которых отвечает за свой класс операций и может присутствовать в процессоре в нескольких экземплярах. Суперскалярность предполагает параллельную работу нескольких функциональных блоков, что возможно лишь при одновременном выполнении нескольких скалярных команд. Последнее условие хорошо сочетается с конвейерной обработкой, при этом желательно, чтобы таких конвейеров было несколько. В этом случае ступень выборки команд, общая для всех конвейеров, должна в каждом такте извлекать из памяти сразу несколько команд. За этой ступенью располагается блок диспетчеризации, отвечающий за распределение команд по конвейерам. По разным оценкам, применение суперскалярного подхода приводит к повышению производительности ВМ в пределах от 1,8 до 8 раз.
\textit{(Источник: Орлов, Цилькер, гл. 9, с. 414--416.)}

\paragraph{Стратегии выдачи и завершения команд.}
В режиме параллельного выполнения нескольких команд процессор должен определить, в какой очередности ему следует выбирать команды из памяти, выполнять эти команды и позволять командам изменять содержимое регистров и ячеек памяти. В самом общем виде стратегии выдачи и завершения команд можно сгруппировать в три категории:
\begin{enumerate}[label=\arabic*)]
\item упорядоченная выдача и упорядоченное завершение;
\item упорядоченная выдача и неупорядоченное завершение;
\item неупорядоченная выдача и неупорядоченное завершение.
\end{enumerate}
Первая стратегия наиболее проста в реализации, но в ней любая задержка в одном конвейере останавливает и другой. Вторая стратегия даёт возможность одному из конвейеров продолжать работать в случае «затора» в другом. Третья стратегия развивает предыдущую концепцию, разрешая процессору нарушать предписанный программой порядок выдачи команд на исполнение. Для обеспечения неупорядоченной выдачи команд в конвейере необходимо максимально развязать ступени декодирования и исполнения; это обеспечивается с помощью буферной памяти, называемой окном команд. Каждая декодированная команда сначала помещается в окно команд; выдача команд из буфера на исполнение определяется не последовательностью их поступления, а мерой готовности. При неупорядоченной выдаче в любой момент времени в стадии исполнения может находиться любое число команд, а степень параллелизма ограничена только числом функциональных блоков.
\textit{(Источник: Орлов, Цилькер, гл. 9, с. 419--421.)}

\paragraph{Переименование регистров.}
Когда команды выдаются и завершаются упорядоченно, каждый регистр в любой точке программы содержит именно то значение, которое диктуется программой. При неупорядоченной выдаче и завершении команд запись в регистры может происходить также неупорядоченно, и отдельные команды, обратившись к какому-то регистру, вместо нужного получат «устаревшее» или «опережающее» значение. Для устранения зависимости команд по данным (типа ЧПЗ и ЗПЗ) используется приём, известный как \emph{переименование регистров} (register renaming). Основная идея переименования регистров состоит в том, что каждый новый результат записывается в один из свободных в данный момент дополнительных регистров, при этом ссылки на заменяемый регистр во всех последующих командах соответствующим образом корректируются. Программист, составляющий программу, имеет дело с именами логических регистров. Число физических регистров аппаратного регистрового файла (АРФ) обычно больше числа логических. «Лишние» регистры АРФ используются в процедуре переименования для временного хранения результатов до момента разрешения конфликтов по данным, после чего значение из регистра временного хранения переписывается на своё «штатное» место. В некоторых процессорах «лишние» регистры в АРФ отсутствуют, а для поддержки переименования предусмотрены специальные структуры, например буфер переименования.
\textit{(Источник: Орлов, Цилькер, гл. 9, с. 421--426.)}

\paragraph{Переупорядочивание команд.}
После декодирования команд и переименования регистров команды передаются на исполнение. Поскольку порядок выполнения команд может отличаться от предписанного программой, необходимо обеспечить корректность их операндов и правильную последовательность занесения результатов в физические регистры. Одним из наиболее распространённых приёмов решения этих проблем служит переупорядочивание команд. В его основе лежит использование окна команд --- буферной памяти, куда помещаются все команды, прошедшие этапы декодирования и переименования регистров. Окно команд обеспечивает отсрочку передачи команд на исполнение до момента готовности операндов, а также нужную очередность завершения команд и загрузки их результатов в физические регистры. Различают централизованное окно команд (scoreboard, впервые предложено в 1964 году фирмой Cray и реализовано в CDC 6600) и распределённое окно команд (станции резервирования, разработаны Р.Л. Томасуло в 1967 году и впервые воплощены в IBM 360/91). В качестве инструмента для поддержания правильной последовательности исполнения команд (при наличии нескольких параллельно работающих функциональных блоков) используется буфер восстановления последовательности (ROB, Reorder Buffer). ROB представляет собой кольцевой буфер; команды заносятся в него в порядке, определяемом программой, и удаляются только тогда, когда они завершены и все предшествующие команды уже удалены. Это гарантирует, что команды покидают ROB строго по порядку и заносят свои результаты в память или регистры только после того, как покинули буфер.
\textit{(Источник: Орлов, Цилькер, гл. 9, с. 426--430.)}

\paragraph{Гиперпотоковая обработка.}
В основе гиперпотоковой технологии (HTT, Hyper-Threading Technology), разработанной фирмой Intel и впервые реализованной в микропроцессоре Intel Xeon MP, лежит тот факт, что современные процессоры в большинстве своём являются суперскалярными и многоконвейерными. Выполняемая программа разбивается на два параллельных потока (нити). Операционная система, поддерживающая гиперпотоковую технологию, воспринимает физический суперскалярный процессор как два логических процессора и организует поступление на эти два процессора двух независимых потоков команд. Процессор с поддержкой HTT эмулирует работу двух одинаковых логических процессоров, принимая команды, направленные для каждого из них. Это не означает, что в процессоре имеются два вычислительных ядра --- оба логических процессора конкурируют за ресурсы единственного вычислительного ядра. Следствием конкуренции является более эффективная загрузка всех ресурсов процессора. Ни один из потоков не считается приоритетным. При остановке одного из потоков (в ожидании какого-либо события или в результате зацикливания) процессор полностью переключается на второй поток. По информации Intel, поддержка гиперпотоковой технологии в процессоре Pentium 4 Xeon потребовала 5\% дополнительной площади на кристалле, в то же время прирост скорости вычислений может достигать 25--35\%.
\textit{(Источник: Орлов, Цилькер, гл. 9, с. 430--434.)}

\subsubsection{Векторные команды}

\paragraph{Понятие векторной обработки.}
Средства векторной обработки оперируют векторами --- одномерными массивами данных (обычно в форме с плавающей запятой), размещённых в памяти ВС. Количество элементов массива называется длиной вектора. Действия над многомерными массивами учитывают специфику их размещения. Способ размещения многомерного массива влияет на шаг изменения адреса элемента, выбираемого из памяти; для обозначения шага, с которым элементы вектора извлекаются из памяти, применяют термин \emph{шаг по индексу} (stride). Векторный процессор --- это процессор, в котором операндами некоторых команд могут выступать массивы данных --- векторы. В архитектуре средств векторной обработки используется один из двух подходов --- векторно-параллельный или векторно-конвейерный. В векторно-параллельном процессоре одновременные операции над элементами векторов проводятся с помощью нескольких функциональных блоков с плавающей запятой, каждый из которых отвечает за одну пару элементов. В векторно-конвейерном варианте обработка элементов векторов производится одним конвейерным функциональным блоком. Наибольшее распространение получили векторно-конвейерные ВС с одним или несколькими функциональными блоками.
\textit{(Источник: Орлов, Цилькер, гл. 13, с. 527--529.)}

\paragraph{Архитектуры «память-память» и «регистр-регистр».}
Принципиальным моментом в архитектуре векторных процессоров является способ доступа к операндам. В векторных процессорах с архитектурой «память-память» элементы векторов поочерёдно извлекаются из памяти и сразу же направляются в функциональный блок; по мере обработки элементы вектора результата сразу же заносятся в память. В архитектуре «регистр-регистр» операнды сначала загружаются из памяти в векторные регистры. Векторный регистр представляет собой совокупность скалярных регистров, объединённых в очередь типа FIFO, способную хранить 50--100 чисел с плавающей запятой (чаще всего --- 64). Операция выполняется над векторами, размещёнными в векторных регистрах операндов, а её результат сначала заносится в векторный регистр результата, а уже из него переписывается в память. Преимущество векторных процессоров «память-память» состоит в возможности обработки длинных векторов, в то время как в процессорах «регистр-регистр» приходится разбивать длинные векторы на сегменты фиксированной длины. Однако за гибкость режима «память-память» приходится расплачиваться относительно большими издержками, известными как время запуска. В настоящее время ВП «регистр-регистр» доминируют на компьютерном рынке.
\textit{(Источник: Орлов, Цилькер, гл. 13, с. 529--531.)}

\paragraph{Регистры управления векторными операциями.}
В состав устройства управления векторными операциями входят несколько регистров специального назначения. \emph{Регистр длины вектора} определяет, сколько элементов фактически содержит обрабатываемый вектор. \emph{Регистр максимальной длины вектора} фиксирует максимальное число элементов вектора, которое может быть одновременно обработано аппаратурой процессора; этот регистр используется при разделении очень длинных векторов на сегменты. \emph{Регистр маски вектора} позволяет выполнять операции не над всеми элементами векторов: в нём каждому элементу вектора соответствует один бит; установка бита в единицу разрешает запись соответствующего элемента вектора результата в выходной векторный регистр, а сброс в ноль --- запрещает. Регистр маски используется также в операциях уплотнения/развертывания (compress/expand). \emph{Регистр индексов} используется в операциях сбора/рассеяния (gather/scatter), в которых вектор в векторном регистре формируется из «нерегулярных» элементов или элементы вектора результата записываются в память нерегулярно.
\textit{(Источник: Орлов, Цилькер, гл. 13, с. 532--533.)}

\paragraph{SIMD-команды в микропроцессорах.}
Название SIMD-команд представляет собой аббревиатуру от Single Instruction Multiple Data --- буквально «одна команда --- много данных». В отличие от обычных команд, оперирующих двумя числами, SIMD-команды обрабатывают сразу две группы чисел; операнды таких команд обычно представлены в одном из упакованных форматов. Идея SIMD-обработки была выдвинута в Институте точной механики и вычислительной техники им. С.А. Лебедева в 1978 году в рамках проекта «Эльбрус-1». С 1992 года команды типа SIMD становятся неотъемлемым элементом АСК микропроцессоров фирм Intel и AMD. Первой на мультимедийный бум отреагировала фирма Intel, добавив в систему команд своего микропроцессора Pentium MMX 57 SIMD-команд. Название MMX (MultiMedia eXtention) разработчики обосновывали тем, что при выборе состава новых команд были проанализированы алгоритмы, применяемые в различных мультимедийных приложениях. Команды MMX обеспечивали параллельную обработку упакованных целых чисел. В дальнейшем наборы SIMD-команд стали обозначать аббревиатурой SSE (Streaming SIMD Extension). В последних разработках фирм Intel и AMD технология обозначена как SSE4, при этом у каждой из фирм она имеет свои особенности. В микропроцессорах Penryn фирмы Intel реализован набор из 47 SIMD-команд (SSE4.1), в микропроцессорах Nehalem (Core i7) той же фирмы --- дополнительно ещё 7 команд (SSE4.2).
\textit{(Источник: Орлов, Цилькер, гл. 2, с. 93--94.)}

\paragraph{Зацепление векторов.}
Интересной особенностью некоторых векторных процессоров типа «регистр-регистр», например ВС фирмы Cray Research Inc., является так называемое зацепление векторов (vector chaining или vector linking), когда векторный регистр результата одной векторной операции используется в качестве входного регистра для последующей векторной операции. Такая комбинация из последовательности умножения и суммирования характерна для операции свёртки и встречается во многих векторных и матричных вычислениях. Сущность зацепления векторов в том, что исполнение векторной команды начинается сразу, как только образуются компоненты участвующих в ней векторных операндов, не дожидаясь завершения вычисления полного вектора операнда и занесения его в соответствующий векторный регистр. Образуются цепочки операций, что даёт существенное ускорение вычислений.
\textit{(Источник: Орлов, Цилькер, гл. 13, с. 534--535.)}

\subsubsection{Краткое резюме подраздела}
Конвейеризация является одним из важнейших архитектурных приёмов повышения производительности. Командный конвейер совмещает во времени этапы цикла команды; его эффективность ограничивается структурными, данными и управленческими рисками. Для борьбы с ними применяются разделение кэшей, ускоренное продвижение данных, остановки конвейера, предсказание переходов (статическое и динамическое), буферы адресов и команд перехода. Арифметические конвейеры используются для выполнения сложных операций с плавающей запятой. Суперскалярные процессоры одновременно выполняют несколько независимых команд за счёт нескольких конвейеров и функциональных блоков; для повышения параллелизма используются неупорядоченная выдача и завершение, переименование регистров, окна команд и буферы восстановления последовательности. Гиперпотоковая обработка позволяет одному физическому ядру обрабатывать два потока команд. Векторные команды и SIMD-расширения обеспечивают параллельную обработку групп данных, что особенно эффективно в мультимедийных и научных приложениях.
\textit{(Обобщение по: Орлов, Цилькер, гл. 2, с. 93--94; гл. 9, с. 386--449; гл. 13, с. 526--535.)}

\subsection{Специализированные процессоры}

\noindent\textit{Основные источники: Орлов С.А., Цилькер Б.Я. Организация ЭВМ и систем: Учебник для вузов. 2-е изд. СПб.: Питер, 2011. Гл. 13, с. 526--556; гл. 8, с. 380--385; гл. 9, с. 434--446.}

\subsubsection{Общая характеристика специализированных процессоров}

\paragraph{Назначение и место в архитектуре ВМ.}
Специализированные процессоры ориентированы на эффективное выполнение определённого класса задач. В отличие от универсальных процессоров, они имеют более узкий набор операций и/или специальную структуру, оптимизированную под конкретный тип вычислений. К специализированным процессорам относятся векторные и векторно-конвейерные процессоры, матричные (array) процессоры, ассоциативные процессоры, систолические структуры, процессоры ввода/вывода и каналы ввода/вывода, а также различные проблемно-ориентированные ускорители (обработка сигналов, мультимедиа, графика, криптография). В учебнике Орлова и Цилькера подробно рассматриваются векторные, матричные, ассоциативные и систолические вычислительные системы, а также каналы и процессоры ввода/вывода.
\textit{(Источник: Орлов, Цилькер, гл. 13, с. 526--556; гл. 8, с. 380--385.)}

\paragraph{Классификация по уровню параллелизма.}
Специализированные процессоры можно классифицировать по уровню параллелизма, который они реализуют. Векторные процессоры используют параллелизм на уровне элементов вектора (SIMD); матричные процессоры — параллелизм на уровне процессорных элементов (SIMD); ассоциативные процессоры — параллелизм поиска и обработки по признаку; систолические структуры — конвейерно-параллельную обработку потоков данных. Процессоры ввода/вывода и каналы обеспечивают параллелизм между работой центрального процессора и операциями ввода/вывода.
\textit{(Источник: Орлов, Цилькер, гл. 13, с. 526--556; гл. 8, с. 380--385.)}

\subsubsection{Матричные процессоры}

\paragraph{Общая структура.}
Матричная вычислительная система (array processor) состоит из массива процессорных элементов (МПЭ) и контроллера массива процессорных элементов (КМП). КМП выполняет последовательную часть программы, декодирует команды и транслирует в массив управляющие сигналы, команды и данные. Каждый процессорный элемент (ПЭ) имеет локальную память и обрабатывает свои данные по общей команде, поступающей из КМП. Для связи с пользователем используется фронтальная ВМ (front-end computer), на которой производится подготовка, компиляция и отладка программ. Перед выполнением программа загружается из фронтальной ВМ в КМП, который затем исполняет последовательную часть и распределяет параллельные команды по ПЭ.
\textit{(Источник: Орлов, Цилькер, гл. 13, с. 535--537.)}

\paragraph{Структура процессорного элемента.}
В большинстве матричных SIMD-систем в качестве процессорных элементов применяются простые RISC-процессоры с локальной памятью ограниченной ёмкости. Например, каждый ПЭ системы MasPar MP-1 состоит из четырёхразрядного процессора с памятью ёмкостью 64 Кбайт; в системе MPP используются одноразрядные процессоры с памятью 1 Кбит; в CM-2 процессорный элемент представляет собой одноразрядный процессор с 64 Кбит локальной памяти. Неотъемлемыми компонентами ПЭ являются арифметико-логическое устройство (АЛУ), регистры данных, сетевой интерфейс, номер процессора, регистр флага разрешения маскирования \(F\) и локальная память. Процессорные элементы, управляемые командами из КМП, могут выбирать данные из своей локальной памяти и регистров, обрабатывать их в АЛУ и сохранять результаты. Обмен данными между ПЭ производится по коммуникационной сети; для передачи результатов в КМП служит шина результата.
\textit{(Источник: Орлов, Цилькер, гл. 13, с. 539--541.)}

\paragraph{Маскирование.}
В процессе вычислений в ряде операций должны участвовать только определённые ПЭ, в то время как остальные остаются бездействующими. Разрешение и запрет работы ПЭ могут исходить от контроллера массива (глобальное маскирование) и реализуются с помощью схем маскирования. В этом случае решение о необходимости маскирования принимается на этапе компиляции кода. Решение о маскировании может также приниматься во время выполнения программы (маскирование, определяемое данными), при этом опираются на хранящийся в ПЭ флаг разрешения маскирования \(F\). При использовании схемы глобального маскирования КМП вместе с командами посылает во все ПЭ глобальную маску; каждый ПЭ декодирует её и выясняет, должен ли он выполнять данную команду. Различают маску в виде \(N\)-разрядного вектора (каждый бит отражает состояние одного ПЭ) и маску с адресом ПЭ (позиции соответствуют разрядам двоичного адреса ПЭ, допускаются значения 0, 1, X).
\textit{(Источник: Орлов, Цилькер, гл. 13, с. 541--542.)}

\paragraph{Сети взаимосвязей.}
Эффективность матричной системы во многом определяется топологией сети взаимосвязей процессорных элементов. Наиболее распространены решётчатые и гиперкубические топологии. Так, в ILLIAC IV, MPP и CM-2 каждый ПЭ соединён с четырьмя соседними; в MP-1 и MP-2 каждый ПЭ связан с восемью смежными ПЭ. В ряде систем (MP-1, MP-2, GF11) реализуются многоступенчатые динамические сети соединений. При синхронной работе ПЭ обмен информацией должен производиться по согласованной схеме; для предотвращения потери данных при одновременной передаче от нескольких источников к одному приёмнику используются специальные механизмы (например, объединение сообщений в CM-2, возможность записи в разные ячейки локальной памяти в MP-1).
\textit{(Источник: Орлов, Цилькер, гл. 13, с. 542--543.)}

\paragraph{Достоинства и недостатки.}
Матричные процессоры обеспечивают высокий уровень параллелизма при выполнении однотипных операций над большими массивами данных. Они эффективны для задач, хорошо распараллеливаемых по данным (обработка изображений, решение сеточных задач, матричные вычисления). Основные недостатки: сложность программирования, необходимость согласования данных, ограниченная применимость для задач с нерегулярной структурой, а также высокая стоимость при большом числе ПЭ.
\textit{(Источник: Орлов, Цилькер, гл. 13, с. 535--543.)}

\subsubsection{Ассоциативные процессоры}

\paragraph{Принцип ассоциативной обработки.}
Ассоциативный процессор строится на базе ассоциативного запоминающего устройства (АЗУ) и осуществляет доступ к данным не по адресу, а по отличительным признакам, содержащимся в самих данных или в приданных им дополнительных разрядах (тегах). В отличие от АЗУ традиционного применения, ассоциативный процессор (АП) обладает средствами обработки данных и возможностью параллельной записи во все ячейки, для которых зафиксировано совпадение с ассоциативным признаком (мультизапись). Идентификация ячеек, отвечающих признаку поиска, обработка содержащихся в них данных и занесение результатов в идентифицированные ячейки осуществляется параллельно по всем словам запоминающего массива.
\textit{(Источник: Орлов, Цилькер, гл. 13, с. 543--544.)}

\paragraph{Структура ассоциативного процессора.}
Запоминающий массив (ЗМ) ассоциативного процессора состоит из \(n\) ячеек. Каждое из \(m+r\)-разрядных слов, хранящихся в ячейках, состоит из двух частей: \(m\)-разрядного слова данных \(Д\) и \(r\)-разрядного признака (тега) \(T\). Тег позволяет отличить данное слово от множества других. При доступе к ЗМ шаблон для поиска читаемого или записываемого слова (ассоциативный признак) заносится в регистр признака \(РгП\), выходы которого поразрядно связаны со схемами сравнения \(СхСр_1,\ldots,СхСр_n\). По второму входу каждая из схем сравнения соединена с теговой частью запоминающего массива. При совпадении ассоциативного признака в \(РгП\) с теговой частью \(i\)-й ячейки вырабатывается сигнал \(c_i\), который идентифицирует ячейку и разрешает выборку из неё данных, выполнение над ними в процессорном элементе \(ПЭ_i\) операции, определённой командой из процессора управления, и занесение результата в \(i\)-ю ячейку. Процессорные элементы АП позволяют выполнять арифметические и логические операции.
\textit{(Источник: Орлов, Цилькер, гл. 13, с. 543--544.)}

\paragraph{Классификация ассоциативных процессоров.}
Способ выполнения операций над словами позволяет определить четыре класса ассоциативных процессоров:
\begin{itemize}
\item \textbf{параллельные} — используют многоразрядные процессорные элементы, каждый из которых работает со своим модулем ассоциативной памяти; обеспечивают наиболее высокое быстродействие, но требуют больших аппаратурных затрат;
\item \textbf{поразрядно-последовательные} — запоминающий массив представляет собой матрицу \(n\times n\) одноразрядных запоминающих элементов; считывание и запись могут производиться по срезу (все разряды одного слова или один разряд всех слов). Каждый разряд среза снабжён одноразрядным процессорным элементом, что позволяет параллельно выполнять операции над разрядами. Наиболее распространённый класс;
\item \textbf{пословно-последовательные} — обработка ведётся последовательно по словам;
\item \textbf{блочно-ориентированные} — ориентированы на задачи информационного поиска.
\end{itemize}
Два последних класса не слишком перспективны для универсальных вычислений.
\textit{(Источник: Орлов, Цилькер, гл. 13, с. 544--545.)}

\paragraph{Ассоциативные вычислительные системы.}
Ассоциативная вычислительная система (АВС) представляет собой многопроцессорную ВС, объединяющую множество ассоциативных процессоров, процессор управления, процессор ввода/вывода и основную память. Программа для АВС хранится в основной памяти. В процессе реализации программы процессор управления выбирает из ОП очередную команду, декодирует её и, если команда предполагает обработку ассоциативными процессорами, передаёт её для параллельного выполнения во все ассоциативные процессоры, готовые к обработке. Процессор управления также обеспечивает обмен данными между ассоциативной памятью каждого АП и основной памятью, а также выполняет функции, аналогичные функциям фронтальной ВМ в матричных системах.
\textit{(Источник: Орлов, Цилькер, гл. 13, с. 546--547.)}

\subsubsection{Систолические структуры}

\paragraph{Идея систолической обработки.}
В фон-неймановских машинах данные, считанные из памяти, однократно обрабатываются в процессорном элементе, после чего снова возвращаются в память. Авторы идеи систолической матрицы Кунг и Лейзерсон предложили организовать вычисления так, чтобы данные на своём пути от считывания из памяти до возвращения обратно пропускались через как можно большее число ПЭ. Если проводить аналогию между структурами ВС и живого организма, то памяти можно отвести роль сердца, множеству ПЭ — роль тканей, а поток данных следует рассматривать как циркулирующую кровь. Отсюда и происходит название «систолическая матрица» (систола — сокращение предсердий и желудочков сердца, при котором кровь нагнетается в артерии). Систолические структуры эффективны при выполнении матричных вычислений, обработке сигналов, сортировке данных.
\textit{(Источник: Орлов, Цилькер, гл. 13, с. 547.)}

\paragraph{Процессорный элемент и принцип работы.}
Основу систолической структуры образует однородная вычислительная среда из процессорных элементов. Типичный ПЭ реализует операцию «умножение с накоплением» (IPS-элемент). На вход ПЭ подаются операнды \(a_{вх}\), \(b_{вх}\); на выходе формируются \(a_{вых}\), \(b_{вых}\) и частичная сумма \(c_{вых}\). На \(n\)-м шаге работы систолической системы ПЭ выполняет операцию
\[
c^{(n)}_{вых}=c^{(n-1)}_{вх}+a^{(n-1)}_{вх}\times b^{(n-1)}_{вх},
\]
при этом операнды на входе и выходе ПЭ одинаковы: \(a^{(n)}_{вых}=a^{(n-1)}_{вх}\), \(b^{(n)}_{вых}=b^{(n-1)}_{вх}\). Частичная сумма поступает на вход ПЭ либо с данного процессорного элемента, либо с соседнего ПЭ матрицы. Вычислительный процесс в систолических структурах представляет собой непрерывную и регулярную передачу данных от одного ПЭ к другому без запоминания промежуточных результатов; каждый элемент входных данных выбирается из памяти однократно и используется столько раз, сколько необходимо по алгоритму. Ввод данных осуществляется в крайние ПЭ матрицы.
\textit{(Источник: Орлов, Цилькер, гл. 13, с. 548--549.)}

\paragraph{Классификация систолических структур.}
Систолические структуры классифицируются по нескольким признакам:
\begin{itemize}
\item по степени гибкости: специализированные (ориентированы на определённый алгоритм), алгоритмически ориентированные (на класс алгоритмов), программируемые (с возможностью программирования ПЭ и конфигурации связей);
\item по разрядности ПЭ: одноразрядные и многоразрядные;
\item по характеру локально-пространственных связей: одномерные (линейные), двухмерные, трёхмерные.
\end{itemize}
Наиболее распространены линейные матрицы (ULA, BLA, TLA). ULA (Unidirectional Linear Array) — однонаправленный линейный массив, где поток данных перемещается в одном направлении; BLA (Bidirectional Linear Array) — двунаправленный массив, в котором два потока данных движутся навстречу друг другу; TLA (Three-path communication Linear Array) — массив с тремя коммуникационными трактами, в котором по разным направлениям перемещаются три потока данных. Выбор конфигурации зависит от вида обрабатываемой информации: одномерные схемы применяются при обработке векторов, двухмерные — матриц, трёхмерные — коллекций иного типа.
\textit{(Источник: Орлов, Цилькер, гл. 13, с. 549--552.)}

\paragraph{Пример вычислений.}
Организацию вычислительного процесса в систолических массивах удобно пояснить на примере умножения матрицы \(A=\{a_{ij}\}\) на вектор \(X=\{x_1,x_2,\ldots,x_n\}\). Элементы вектора произведения \(Y=\{y_1,y_2,\ldots,y_n\}\) могут быть получены периодически повторяющимися операциями
\[
y_i(1)=0;\qquad y_i(k+1)=y_i(k)+a_i(k)\times x_i(k);\qquad y_i=y_i(n+1),
\]
где \(k\) — номер шага вычислений. Компоненты \(y_i\) вектора \(Y\), имеющие в начальный момент нулевое значение, поступают на вход массива и продвигаются через ПЭ справа налево, в то время как компоненты вектора \(X\) движутся слева направо. Элементы матрицы \(A=\{a_{ij}\}\) в порядке, определяемом схемой, вводятся в ПЭ сверху вниз. Промежуточные результаты \(y_i(k)\) накапливаются по мере продвижения от одного ПЭ к другому. В каждый момент времени активны только \(\frac{p+q-1}{2}\) процессорных элементов (где \(p+q-1\) — ширина полосы ненулевых элементов матрицы), поэтому каждый выходной результат формируется за два такта. Для вычисления всех \(n\) элементов выходного вектора \(Y\) необходимо \(2n+p+q-1\) тактов.
\textit{(Источник: Орлов, Цилькер, гл. 13, с. 553--555.)}

\paragraph{Достоинства и недостатки.}
Систолические структуры обладают следующими особенностями:
\begin{itemize}
\item вычислительный процесс представляет собой непрерывную и регулярную передачу данных от одного ПЭ к другому без запоминания промежуточных результатов;
\item каждый элемент входных данных выбирается из памяти однократно и используется многократно;
\item образующие структуру ПЭ однотипны и могут быть менее универсальными, чем процессоры обычных многопроцессорных систем;
\item потоки данных и управляющих сигналов обладают регулярностью, что позволяет объединять ПЭ локальными связями минимальной длины;
\item алгоритмы функционирования позволяют совместить параллелизм с конвейерной обработкой данных;
\item производительность матрицы можно улучшить за счёт добавления определённого числа ПЭ, причём коэффициент повышения производительности при этом линеен.
\end{itemize}
К недостаткам можно отнести сложность программирования и ограниченную применимость для задач с нерегулярной структурой.
\textit{(Источник: Орлов, Цилькер, гл. 13, с. 549.)}

\subsubsection{Процессоры и каналы ввода/вывода}

\paragraph{Назначение.}
Процессоры ввода/вывода (каналы ввода/вывода) освобождают центральный процессор от управления периферийными устройствами. Канал ввода/вывода — это специализированный процессор, который выполняет канальную программу и управляет обменом между периферийными устройствами и основной памятью без участия центрального процессора. Все функции ЦП сводятся к запуску и остановке операций в канале, а также проверке состояния канала и подключённых к нему ПУ. Для этих целей ЦП использует лишь несколько (от 4 до 8) команд ввода/вывода.
\textit{(Источник: Орлов, Цилькер, гл. 8, с. 380--381.)}

\paragraph{Канальная программа и управляющие слова канала.}
КВВ реализует операции ввода/вывода путём выполнения канальной программы. Канальные программы для каждого ПУ, с которым предполагается обмен информацией, описывают нужную последовательность операций ввода-вывода и хранятся в основной памяти ВМ либо локальной памяти канала. Роль команд в канальных программах выполняют управляющие слова канала (УСК). Типовое УСК содержит:
\begin{itemize}
\item код операции, определяющий для КВВ и ПУ тип операции: «Записать» (вывод информации из ОП в ПУ), «Прочитать» (ввод информации из ПУ в ОП), «Управление» (перемещение головок НМД, магнитной ленты и т.п.);
\item указатели — дополнительные предписания, задающие более сложную последовательность операций В/ВЫВ;
\item адрес данных, указывающий область памяти, используемую в операции ввода-вывода;
\item счётчик данных, хранящий значение длины передаваемого блока данных.
\end{itemize}
Кроме того, в УСК может содержаться идентификатор ПУ и информация о его уровне приоритета, указания по действиям при возникновении ошибок и т.п.
\textit{(Источник: Орлов, Цилькер, гл. 8, с. 381--382.)}

\paragraph{Структура ВМ с канальной системой ввода/вывода.}
Обмен информацией между КВВ и основной памятью осуществляется посредством системной шины ВМ. ПУ подключаются к каналу не непосредственно, а через блоки управления периферийными устройствами (БУПУ). БУПУ принимает от канала приказы по управлению периферийным устройством (чтение, запись, перемещение носителя или магнитной головки и т.п.) и преобразует их в сигналы управления, свойственные данному типу ПУ. Обычно один БУПУ может обслуживать несколько однотипных ПУ, но для подключения быстродействующих периферийных устройств часто применяются индивидуальные блоки управления. В свою очередь, некоторые ПУ могут подключаться одновременно к нескольким БУПУ. Это позволяет воспользоваться свободным трактом другого БУПУ при занятости данного БУПУ обслуживанием одного из подключённых к нему ПУ. Физически БУПУ может быть самостоятельным устройством или интегрирован с ПУ или каналом.
\textit{(Источник: Орлов, Цилькер, гл. 8, с. 382--383.)}

\paragraph{Режимы работы канала.}
В пределах канала ввода/вывода считается, что каждое ПУ подключено к своему подканалу. Подканалы имеют свои уникальные логические номера, с помощью которых канальная программа адресуется к конкретному ПУ. Физически подканал реализуется в виде участка памяти, в котором хранятся параметры операции ввода/вывода, выполняемой данным ПУ: текущие значения адреса и счётчика данных, код и указатели операции ввода/вывода, адрес следующего УСК и др. Для хранения этих параметров обычно используется локальная память канала.
\textit{(Источник: Орлов, Цилькер, гл. 8, с. 383.)}

Способ организации взаимодействия ПУ с каналом определяется соотношением быстродействия ОП и ПУ. По этому признаку ПУ образуют две группы: быстродействующие (НМД, НМЛ) со скоростью приёма и выдачи информации около 1 Мбайт/с и медленнодействующие (дисплеи, печатающие устройства и др.) со скоростями порядка 1 Кбайт/с и менее. С учётом производительности ПУ в КВВ реализуются два режима работы: мультиплексный (режим разделения времени) и монопольный. В мультиплексном режиме несколько периферийных устройств разделяют канал во времени; в монопольном режиме после установления связи между каналом и ПУ последнее монополизирует канал на всё время до завершения инициированной процессором канальной программы. Мультиплексный режим реализуется мультиплексным каналом (байт-мультиплексным или блок-мультиплексным), монопольный — селекторным каналом.
\textit{(Источник: Орлов, Цилькер, гл. 8, с. 383--384.)}

\paragraph{Канальная подсистема.}
В последовавших за IBM 360 семействах универсальных ВМ IBM 370 и особенно IBM 390 концепция системы ввода/вывода на базе каналов получила дальнейшее развитие и вылилась в так называемую канальную подсистему ввода/вывода (КПВВ). Главная идея заключается в интегрировании отдельных КВВ в единый специализированный процессор ввода/вывода с большим числом канальных трактов и подканалов. Блоки управления периферийными устройствами обычно подключаются к нескольким канальным трактам, что позволяет динамически менять путь пересылки информации с учётом текущей их загруженности. Кроме того, разные канальные тракты могут обладать различной пропускной способностью и при выборе трактов для подключения определённых ПУ может быть учтено их быстродействие. В КПВВ реализованы параллельные и последовательные канальные тракты; для работы с ПУ, соединёнными с КПВВ волоконно-оптическими линиями, используются последовательные канальные тракты, реализующие протокол ESCON. Основное преимущество КПВВ — динамическое перераспределение канальных трактов — позволяет применить оптимальную стратегию динамического распределения и использования этих ресурсов и достичь качественно нового уровня эффективности системы ввода/вывода.
\textit{(Источник: Орлов, Цилькер, гл. 8, с. 384--385.)}

\subsubsection{Краткое резюме подраздела}
Специализированные процессоры ориентированы на эффективное выполнение определённых классов задач. Матричные процессоры реализуют SIMD-параллелизм на уровне массива процессорных элементов; ассоциативные процессоры осуществляют доступ к данным по признаку и параллельную обработку; систолические структуры обеспечивают конвейерно-параллельную обработку потоков данных с многократным использованием каждого элемента. Процессоры и каналы ввода/вывода освобождают центральный процессор от управления периферийными устройствами, выполняя канальные программы и обеспечивая прямой доступ к памяти. Каждый класс специализированных процессоров имеет свои достоинства и недостатки, определяющие области их применения.
\textit{(Обобщение по: Орлов, Цилькер, гл. 8, с. 380--385; гл. 13, с. 526--556.)}

\subsection{Машины, обеспечивающие выполнение вычислений, управляемых потоком данных}

\noindent\textit{Основные источники: Орлов С.А., Цилькер Б.Я. Организация ЭВМ и систем: Учебник для вузов. 2-е изд. СПб.: Питер, 2011. Гл. 15, с. 579--597.}

\subsubsection{Механизмы управления вычислениями}

\paragraph{Традиционное управление последовательностью команд.}
Большинство современных ВМ построены в соответствии с фон-неймановской концепцией программного управления. Согласно этой концепции порядок вычислений определяет программа, представленная последовательностью команд. Эта последовательность сохраняется и при размещении программы в памяти. Команды программы, как правило, выполняются в порядке их следования в программе. Это обеспечивается наличием в каждом процессоре счётчика команд. Выполнение команд в каждом процессоре — поочерёдное и потому достаточно медленное. Для получения выигрыша программист или компилятор должны определить независимые команды, которые могут быть поданы на отдельные процессоры, причём так, чтобы коммуникационные издержки были не слишком велики.
\textit{(Источник: Орлов, Цилькер, гл. 15, с. 579.)}

Традиционные (фон-неймановские) вычислительные системы, управляемые с помощью счётчика команд, иногда называют вычислительными системами, управляемыми последовательностью команд (control flow computer). Данный термин применяется, когда нужно выделить этот тип ВС из альтернативных типов, где последовательность выполнения команд определяется не центральным устройством управления со счётчиком команд, а каким-либо иным способом. В вычислительных системах используют три механизма управления последовательностью вычислений:
\begin{enumerate}[label=\arabic*)]
\item команда выполняется, после того как выполнена предшествующая ей команда последовательности;
\item команда выполняется, когда становятся доступными её операнды;
\item команда выполняется, когда другим командам требуется результат её выполнения.
\end{enumerate}
Первый метод соответствует традиционному механизму с управлением последовательностью команд; второй механизм известен как управляемый данными (data driven) или потоковый (dataflow); третий вариант называют механизмом управления по запросу (demand driven) или редукционным.
\textit{(Источник: Орлов, Цилькер, гл. 15, с. 579--580.)}

\subsubsection{Вычислительная модель потоковой обработки}

\paragraph{Граф потоков данных.}
В потоковой вычислительной модели для описания вычислений используется ориентированный граф, иногда называемый графом потоков данных (dataflow graph). Этот граф состоит из вершин (узлов), отображающих операции, и рёбер или дуг, показывающих потоки данных между теми вершинами графа, которые они соединяют. Операция в вершине выполняется лишь когда по дугам в неё поступила вся необходимая информация. Обычно такая операция требует одного или двух операндов, а для условных операций необходимо наличие входного логического значения. Операция формирует один, два или один из двух возможных результатов. Таким образом, у каждой вершины может быть от одной до трёх входящих дуг и одна или две выходящих. После активации вершины и выполнения операции результат передаётся по дуге к ожидающей вершине. Процесс повторяется, пока не будут активированы все вершины и получен окончательный результат. Одновременно могут быть активированы несколько вершин, при этом параллелизм в вычислительной модели выявляется автоматически.
\textit{(Источник: Орлов, Цилькер, гл. 15, с. 580--581.)}

\paragraph{Токены и движение данных.}
Данные (операнды/результаты), перемещаемые вдоль дуг, содержатся в опознавательных информационных кадрах, маркерах специального формата — токенах (иначе фишках или маркерах доступа). На рис. 15.2 (см. учебник) показан простой потоковый граф для вычисления выражения \(f=a/b+b\times c\). Входами служат переменные \(a\), \(b\) и \(c\), изображённые вверху графа. Дуги между вершинами показывают информационную взаимосвязь операций. Направление вычислений — сверху вниз. Используются три вычислительные операции: сложение, умножение и деление. Заметим, что \(b\) требуется в двух вершинах. Вершина копирования dup предназначена для формирования дополнительной копии переменной \(b\). После поступления на граф входной информации токен, содержащий значение \(a\), направляется в вершину деления; токен с переменной \(b\) — в вершину копирования; токен с переменной \(c\) — в вершину умножения. Активирована может быть только вершина dup, поскольку у неё лишь один вход и на нём уже присутствует токен. Когда токен из вершины dup появится на её выходных дугах, узлы умножения и деления также получат все необходимые маркеры доступа и могут быть активированы. Последняя вершина ждёт завершения операций умножения и деления, то есть когда на её входе появятся все необходимые токены.
\textit{(Источник: Орлов, Цилькер, гл. 15, с. 581.)}

\paragraph{Примитивы узлов потокового графа.}
Практические вычисления требуют некоторых дополнительных возможностей, например при выполнении команд условного перехода. По этой причине в потоковых графах предусмотрены вершины-примитивы нескольких типов:
\begin{itemize}
\item \textbf{двухвходовая операционная вершина} — операции производятся над данными, поступающими с левой и правой входных дуг, а результат выводится через выходную дугу;
\item \textbf{одновходовая операционная вершина} — операции выполняются над входными данными, результат выводится через выходную дугу;
\item \textbf{вершина ветвления} — осуществляет копирование входных данных и их вывод через две выходных дуги (комбинацией таких узлов можно строить вершины ветвления на \(m\) выходов);
\item \textbf{вершина слияния} — данные поступают только с какого-нибудь одного из двух входов и без изменения подаются на выход (комбинируя такие узлы, можно строить вершины слияния с \(m\) входами);
\item \textbf{вершина управления}, которая существует в трёх вариантах:
  \begin{itemize}
  \item \textbf{TF-коммутатор} — если значение на управляющем входе истинно (T — True), то входные данные выводятся через левый выход, а при ложном значении (F — False) данные следуют через правый выход;
  \item \textbf{вентиль} — при истинном значении на входе управления данные выводятся через выходную дугу;
  \item \textbf{арбитр} — данные, поступившие на один из двух входов первыми, следуют через левую дугу, а прибывшие впоследствии — через правую выходную дугу (активация вершины происходит в момент прихода данных на какой-либо из входов).
  \end{itemize}
\end{itemize}
Процесс потоковой обработки может быть аналогичен конвейерному режиму: после обработки первого набора входных сигналов на вход графа может быть подан второй и т.д. Отличие состоит в том, что промежуточные результаты (токены) первого вычисления не обрабатываются совместно с промежуточными результатами второго и последующих вычислений.
\textit{(Источник: Орлов, Цилькер, гл. 15, с. 581--582.)}

\subsubsection{Архитектура потоковых вычислительных систем}

\paragraph{Представление программы в памяти.}
В потоковых ВС программа вычислений определена потоковым графом, который хранится в памяти системы в виде таблицы. На рис. 15.4 (см. учебник) показаны пример графа потоковой программы и содержание соответствующей ему таблицы. Каждая запись в таблице представляет одну вершину потокового графа. Порядок размещения записей в памяти несущественен, поскольку момент активации вершины определяется лишь наличием данных на всех её входах. Об этом можно судить по состоянию битов наличия (поле pres). В примере предполагается, что у вершины может быть не более двух входных дуг, поэтому достаточно двух битов наличия. Единичное состояние бита свидетельствует о поступлении данных на соответствующий вход вершины, а также о том, что они находятся в поле операнда, выделенное для этого входа. Если вершина имеет только одну входную дугу, другой бит сразу устанавливается в единицу. В каждой записи содержится информация, эквивалентная коду операции в обычной машинной команде (поле opc), а также указания о тех вершинах и их входах, куда должен быть передан результат (поле des). Отметим, что в полях операндов (opr1, opr2) должны располагаться значения операндов, а не ссылки на их расположение в памяти. На каждом шаге вычислений одновременно активируются все вершины, в записях которых все биты наличия содержат единицы.
\textit{(Источник: Орлов, Цилькер, гл. 15, с. 582--583.)}

\subsubsection{Статические потоковые вычислительные системы}

\paragraph{Правило активации и структура токена.}
Статическая потоковая архитектура, известная также под названием «единственный токен на дугу» (single-token-per-arc dataflow), была предложена Деннисом в 1975 году. В ней допускается присутствие на дуге графа не более одного токена. Это выражается в правиле активации вершины: вершина активируется, когда на всех её входных дугах присутствует по токену и ни на одном из её выходов токенов нет. Токен в статической потоковой системе представляет собой триаду: \(\langle v, \langle f, n \rangle, a \rangle\), где \(v\) — это данные, переносимые токеном, \(f\) — идентификатор функции, реализуемой текущим потоковым графом, \(n\) — номер целевой вершины, куда должен поступить данный токен, и, наконец, \(a\) — это номер дуги, по которой токен должен быть направлен к целевой вершине. Поле \(\langle f, n \rangle\) называют тегом, оно определяет предназначение токена определённой вершине. Для указания вершине о том, что её выходной токен уже востребован последующей вершиной (вершинами) графа, в ВС обычно прибегают к механизму подтверждения с квитированием связи.
\textit{(Источник: Орлов, Цилькер, гл. 15, с. 583--584.)}

\paragraph{Структура процессорного элемента.}
Типовую статическую потоковую архитектуру иллюстрирует рис. 15.6 (см. учебник). Память действий состоит из двух блоков памяти: памяти команд/данных и управляющей памяти. Вершина потокового графа представлена в памяти команд/данных кадром, содержащим поле кода операции, \(N\) полей операндов и \(K\) полей «вершина/дуга». Каждому кадру в памяти команд/данных соответствует кадр в управляющей памяти, содержащий биты наличия операндов и занятости дуг. Бит наличия операнда устанавливается в единицу, если этот операнд доступен, то есть если токен, содержащий данный операнд, уже поступил по входной дуге графа. Для каждого поля «вершина/дуга» установленный бит занятости означает, что выходная дуга, ассоциированная с данным полем, не содержит токена. Вершина графа, описанная в памяти команд/данных, может быть активирована (операция может быть выполнена), если все биты в соответствующем кадре памяти управления установлены в единицу. Когда данная ситуация распознаётся блоком обновления, он помещает пакет команды в очередь команд. Опираясь на пакет команды и содержимое памяти действий, блок выборки формирует пакет операции и направляет его в один из свободных функциональных блоков. После выполнения требуемой операции функциональный блок создаёт пакет результата и передаёт его в блок обновления, который в соответствии с полученным результатом обновляет содержимое памяти действий.
\textit{(Источник: Орлов, Цилькер, гл. 15, с. 584--585.)}

\paragraph{Достоинства и недостатки.}
Основное преимущество рассматриваемой модели потоковых вычислений заключается в упрощённом механизме обнаружения активированных узлов. К сожалению, статическая модель обладает множеством серьёзных недостатков, однако основной недостаток заключается в том, что данный механизм не допускает параллельного выполнения независимых итераций цикла. Например, операция сложения двух векторов может быть представлена циклическим процессом, где на каждой итерации суммируется пара одноимённых элементов векторов-операндов. В статической потоковой ВС такое суммирование должно выполняться последовательно в порядке увеличения индексов элементов. Это связано с тем, что, во-первых, на каждой входной дуге вершины сложения может находиться лишь по одному элементу, а во-вторых, из-за того, что в теге элемента отсутствует информация о его индексе в векторе. В качестве примеров статических потоковых ВС можно упомянуть: LAU System, TI’s Distributed Data Processor, DDMI Utah Data Driven Machine, NEC Image Pipelined Processor, Hughes Dataflow Multiprocessor.
\textit{(Источник: Орлов, Цилькер, гл. 15, с. 585.)}

\subsubsection{Динамические потоковые вычислительные системы}

\paragraph{Архитектура с помеченными токенами.}
Производительность потоковых систем существенно возрастает, если они в состоянии поддерживать дополнительный уровень параллелизма, соответствующий одновременному выполнению отдельных итераций цикла или параллельной обработке пар элементов в векторных операциях. Кроме того, в современных языках программирования активно используются так называемые реентерабельные процедуры, когда в памяти хранится только одна копия кода процедуры, но эта копия является повторно входимой (реентерабельной). Это означает, что к процедуре можно ещё раз обратиться, не дожидаясь завершения действий в соответствии с предыдущим входом в данную процедуру. Отсюда желательно, чтобы все обращения к реентерабельной процедуре также обрабатывались параллельно. Задача обеспечения дополнительного уровня параллелизма решается в динамических потоковых ВС и реализуется двумя вариантами архитектуры потоковой ВС: архитектуры с помеченными токенами и архитектуры с явно адресуемыми токенами.
\textit{(Источник: Орлов, Цилькер, гл. 15, с. 585.)}

Системы с помеченными токенами (tagged-token architecture) в известной мере свободны от основного недостатка статической модели. В них число токенов, одновременно присутствующих на дуге, не ограничивается. Для приведённого ранее примера циклического суммирования элементов векторов это открывает возможность выполнения итераций в произвольной последовательности, но требует учёта принадлежности токенов к одной и той же итерации. С этой целью токен должен содержать информацию о вычислительном контексте, в котором он используется, например о номере итерации цикла. Этот контекст называют «цветом значения», а токен соответственно называют «окрашенным», в силу чего метод имеет ещё одно название — метод окрашенных токенов. В модели с помеченными токенами структура токена сложнее, чем в статической модели: \(\langle v, \langle f, n, c, i \rangle, a \rangle\), где \(c\) определяет фрагмент кода или тело цикла в составе реализуемой функции \(f\), а \(i\) (индекс) представляет цвет токена. Каждая дуга потокового графа может рассматриваться как вместилище, способное содержать произвольное число токенов с различными тегами. Правило активирования вершины в модели с помеченными токенами имеет вид: вершина активируется, когда на всех её входных дугах присутствуют токены с идентичным цветом.
\textit{(Источник: Орлов, Цилькер, гл. 15, с. 586.)}

Типовая структура потоковой системы с помеченными токенами показана на рис. 15.7 (см. учебник). Для обнаружения одинаково окрашенных токенов (токенов с одинаковыми тегами) в процессорный элемент введён согласующий блок. Этот блок получает очередной токен из очереди токенов и проверяет, нет ли в памяти согласования его партнёра (токена с идентичным тегом). Если такой партнёр не обнаружен, принятый токен заносится в память согласования. Если же токен-партнёр уже хранится в памяти согласования, то согласующий блок удаляет его оттуда и направляет оба токена с совпавшими тегами в блок выборки. На основе общего тега блок выборки находит в памяти команд/данных соответствующую команду и формирует пакет операции, который затем направляет в функциональный блок. Функциональный блок выполняет операцию, создаёт токены результата и помещает их в очередь токенов. Примеры динамических потоковых вычислительных систем с архитектурой окрашенных токенов: система Гурда и Ватсона, созданная в Манчестерском университете в 1980 году, системы SIGMA-1, NTT’s Dataflow Processor Array System, DDDP Distributed Data Driven Processor, SDFA Stateless Data-Flow Architecture.
\textit{(Источник: Орлов, Цилькер, гл. 15, с. 586--587.)}

Основное преимущество динамических потоковых систем с помеченными токенами — повышенная производительность, достигаемая за счёт возможности присутствия на дуге множества токенов. При этом, однако, основной проблемой становится эффективная реализация блока, который собирает токены с совпадающим цветом. В плане производительности этой цели наилучшим образом отвечает ассоциативная память. К сожалению, такое решение является слишком дорогостоящим, поскольку число токенов, ожидающих совпадения тегов, как правило, достаточно велико. По этой причине в большинстве вычислительных систем вместо ассоциативных запоминающих устройств используются обычные адресные ЗУ. В частности, сравнение тегов в упомянутой выше манчестерской системе производится с привлечением хэширования, что несколько снижает быстродействие. В последнее время всё более популярной становится другая организация динамической потоковой ВС, позволяющая освободиться от ассоциативной памяти и известная как архитектура с явно адресуемыми токенами.
\textit{(Источник: Орлов, Цилькер, гл. 15, с. 587.)}

\paragraph{Архитектура с явно адресуемыми токенами.}
Значительным шагом в архитектуре потоковых ВС стало изобретение механизма явной адресации токенов (explicit token-store), имеющего и другое название — непосредственное согласование (direct matching). В основе этого механизма лежит наблюдение: все токены в одной и той же итерации цикла и в одном и том же вхождении в реентерабельную процедуру имеют идентичный тег (цвет). При инициализации очередной итерации цикла или очередном обращении к процедуре формируется так называемый кадр токенов, содержащий токены, относящиеся к данной итерации или данному обращению, то есть токены с одинаковыми тегами. Использование конкретных ячеек внутри кадра задаётся на этапе компиляции. Каждому кадру выделяется отдельная область в специальной памяти кадров (frame memory), причём раздача памяти под каждый кадр происходит уже на этапе выполнения программы.
\textit{(Источник: Орлов, Цилькер, гл. 15, с. 587.)}

В схеме с явной адресацией токенов любое вычисление полностью описывается указателем команды (IP, Instruction Pointer) и указателем кадра (FP, Frame Pointer). Токен выглядит следующим образом: \(\langle v, \langle FP, IP \rangle \rangle\). Команды, реализующие потоковый граф, хранятся в памяти команд и имеют формат: opc, i, dis. Здесь \(i\) (индекс в памяти кадров) определяет положение ячейки с нужным токеном внутри кадра, то есть какое число нужно добавить к FP, чтобы получить адрес интересующего токена. Поле dis указывает на местоположение команды, которой должен быть передан результат обработки данного токена. Адрес в этом поле также задан в виде смещения — числа, которое следует прибавить к текущему значению IP, чтобы получить исполнительный адрес команды назначения в памяти команд. Если потребителей токена несколько, в поле dis заносится несколько значений смещения. Простой пример кодирования потокового графа и токенов на его дугах показан на рис. 15.8 (см. учебник). Каждому слову в памяти кадров придан бит наличия, единичное значение которого удостоверяет, что в ячейке находится токен, ждущий согласования, то есть что одно из искомых значений операндов уже имеется. Когда на входную дугу вершины поступает токен, в ячейке памяти кадров проверяется бит наличия. Если бит наличия сброшен (ни один из пары токенов ещё не поступал), поле значения пришедшего токена заносится в анализируемую ячейку памяти кадров, а бит наличия устанавливается в единицу. Если токен приходит в вершину, для которой уже хранится значение, команда, представляющая данную вершину, может быть активирована и выполнена с операндами; значение извлекается из памяти кадров, бит наличия сбрасывается, и на функциональный блок передаётся пакет команды, содержащий операнды, код операции и адресат результата.
\textit{(Источник: Орлов, Цилькер, гл. 15, с. 587--589.)}

Типовая структура системы с явной адресацией токенов показана на рис. 15.9 (см. учебник). Отметим, что функция согласования токенов стала достаточно короткой операцией, что позволяет внедрить её в виде нескольких ступеней процессорного конвейера.
\textit{(Источник: Орлов, Цилькер, гл. 15, с. 589.)}

\subsubsection{Мультипотоковые вычислительные системы}

\paragraph{Идея мультипотоковой обработки.}
Рассмотренный ранее механизм обработки с управлением от потока данных функционирует на уровне команд и его относят к потоковой обработке низкого уровня (fine-grain dataflow). Данному подходу сопутствуют большие издержки при пересылке операндов, но, главное, не сохраняется принцип локальности (порядок использования команд и данных — произвольный), из-за чего теряются преимущества использования кэш-памяти. Справиться с этими проблемами помогает потоковая обработка на процедурном уровне, так называемая укрупнённая потоковая или мультипотоковая обработка (multithreading). Буквальный перевод английского термина означает потоковую обработку множества нитей.
\textit{(Источник: Орлов, Цилькер, гл. 15, с. 589--590.)}

Мультипотоковая модель совмещает локальность программы, характерную для фон-неймановской модели, с толерантностью к задержкам на переключение задач, характерным для потоковой архитектуры. Данная возможность обеспечивается тем, что вершина графа представляет собой не одну команду, а последовательность из нескольких команд, называемую нитью (thread). Именно поэтому мультипотоковую организацию часто называют крупнозернистой потоковой обработкой (coarse-grained dataflow). Мультипотоковая обработка сводится к потоковому выполнению нитей, в то время как внутри отдельной нити характер выполнения фон-неймановский. Порядок обработки нитей меняется динамически в процессе вычислений, а последовательность команд в пределах нити определена при компиляции статически. Структура мультипотоковой ВС представлена на рис. 15.10 (см. учебник). Существенное отличие мультипотоковой системы от обычной потоковой состоит в организации внутреннего управляющего конвейера, где последовательность выполнения команд задаётся счётчиком команд, как в фон-неймановских машинах. Иными словами, этот конвейер идентичен обычному конвейеру команд.
\textit{(Источник: Орлов, Цилькер, гл. 15, с. 590.)}

\paragraph{Формы мультипотоковой обработки.}
Существует две формы мультипотоковой обработки: без блокирования и с блокированием. В модели без блокирования выполнение нити не может быть начато, пока не получены все необходимые данные. Будучи запущенной, нить выполняется до конца без приостановки. В варианте с блокированием запуск нити может быть произведён до получения всех операндов. Когда требуется отсутствующий операнд, нить приостанавливается (блокируется), а возобновление выполнения откладывается на некоторое время. Процессор запоминает всю необходимую информацию о состоянии и загружает на выполнение другую готовую нить. Модель с блокированием обеспечивает более мягкий подход к формированию нитей (часто это выражается в возможности использования более длинных нитей) за счёт дополнительной аппаратуры для хранения блокированных нитей.
\textit{(Источник: Орлов, Цилькер, гл. 15, с. 590--591.)}

Возможна также и потоковая обработка переменного уровня, когда вершины соответствуют как простым операциям, так и сложным последовательным процедурам. Последний случай иногда называют комбинированной обработкой с потоками данных и потоками управления (combined dataflow/control flow).
\textit{(Источник: Орлов, Цилькер, гл. 15, с. 591.)}

\subsubsection{Вычислительные системы с управлением по запросу (редукционные ВС)}

\paragraph{Идея управления по запросу.}
В системах с управлением от потока данных каждая команда, для которой имеются все необходимые операнды, немедленно выполняется. Однако для получения окончательного результата многие из этих вычислений оказываются ненужными. Отсюда прагматичным представляется иной подход, когда вычисления инициируются не по готовности данных, а на основе запроса на данные. Такая организация вычислительного процесса носит название управления вычислениями по запросу (demand-driven control). В её основе, как и в потоковой модели (data-driven control), лежит представление вычислительного процесса в виде графа. В потоковой модели вершины вверху графа запускаются раньше, чем нижние. Это — нисходящая обработка. Механизм управления по запросу состоит в обработке вершин потокового графа снизу вверх (вершина запускается лишь когда требуется её результат). Данный процесс получил название редукции графа, а ВС, работающая в режиме снизу вверх, называется редукционной вычислительной системой.
\textit{(Источник: Орлов, Цилькер, гл. 15, с. 593.)}

Математическую основу редукционных ВС составляет лямбда-исчисление, а для написания программ под такие системы нужны так называемые функциональные языки программирования (FP, Haskell и др.). На функциональном языке все программы представляются в виде выражений, а процесс выполнения программы заключается в определении значений последних (это называется оценкой выражения). Оценка выражения производится посредством повторения операции выбора и упрощения тех частей выражения, которые можно свести от сложного к простому (такая часть выражения называется редексом (от REDuced EXpression — приведённое выражение), причём сам редекс также является отдельным выражением). Операция упрощения называется редукцией. Процесс редукции завершается, когда преобразованное редукцией выражение больше не содержит редекса. Выражение, не содержащее редекса, называется нормальной формой.
\textit{(Источник: Орлов, Цилькер, гл. 15, с. 593--594.)}

В редукционной ВС вычисления производятся по запросу на результат операции. Предположим, что вычисляется выражение \(a=(b+1)\times c-d/c\). В случае потоковых моделей процесс начинается с самых внутренних операций, а именно с параллельного вычисления \((b+1)\) и \(d/c\). Затем выполняется операция умножения \((b+1)\times c\) и, наконец, самая внешняя операция — вычитание. Такой род вычислений часто называют энергичными вычислениями (eager evaluation). При вычислениях, управляемых запросами, всё начинается с запроса на результат \(a\), который включает в себя запрос на вычисление выражений \((b+1)\times c\) и \(d/c\), а те, в свою очередь, формируют запрос на вычисление \(b+1\), то есть на операцию самого внутреннего уровня. Результат возвращается в порядке, обратном поступлению запросов. Отсюда название — ленивые вычисления (lazy evaluation), поскольку операции выполняются только тогда, когда их результат требуется другой операции. Редукционные вычисления естественно согласуются с концепцией функционального программирования, упрощающей распараллеливание программ.
\textit{(Источник: Орлов, Цилькер, гл. 15, с. 594.)}

\paragraph{Строчная и графовая модели редукции.}
Известны два типа моделей редукционных систем: строчная и графовая, отличающиеся тем, что именно передаётся в операцию (функцию) — скопированные значения данных или же только указатели мест хранения данных. В строчной редукционной модели каждая запросившая вершина получает отдельную копию выражения для собственной оценки. Длинное строковое выражение рекурсивным образом сокращается (редуцируется) до единственного значения. Каждый шаг редукции содержит операцию, сопровождаемую ссылкой на требуемые входные операнды. Операция приостанавливается, пока оцениваются входные параметры. На рис. 15.12 (см. учебник) показан процесс вычисления с помощью редукционной ВС значения выражения \(a=b-c\) (\(b=d+e\), \(c=f\times g\)) для \(d=1\), \(e=3\), \(f=5\), \(g=7\). Программа редукции состоит из распознавания редексов с последующей заменой их вычисленными значениями. Таким образом, вся программа в конечном итоге редуцируется до результата. На рис. 15.13 (см. учебник) показан процесс вычислений с помощью строчной редукции: если требуется значение \(a=(b+c)\times(b-c)\), то копируется граф программы, определяющий вычисление \(a\); запускается операция умножения; поскольку это вычисление невозможно без предварительного расчёта двух параметров, то запускаются вычисления сложения и вычитания, в результате чего образуется редуцированный граф (с ветвями «6» и «2»); результат получается путём дальнейшей редукции.
\textit{(Источник: Орлов, Цилькер, гл. 15, с. 595.)}

В графовой редукционной модели выражение представлено как ориентированный граф. Граф сокращается по результатам оценки ветвей и подграфов. В зависимости от запросов возможно параллельное оценивание и редукция различных частей графа или подграфов. Запросившему узлу, который управляет всеми ссылками на граф, возвращается указатель на результат редукции. «Обход» графа и изменение ссылок продолжаются, пока не будет получено значение результата, копия которого возвращается запросившей команде. Рисунок 15.14 (см. учебник) иллюстрирует пример вычисления \(a=(b+c)\times(b-c)\) с помощью графовой редукции. В этой модели для нахождения значения \(a\) граф программы не копируется, а лишь передаётся указатель на этот граф. При достижении вершины умножения направление переданного указателя меняется на противоположное, указывая место, куда будет выдаваться итоговый результат вычисления. Далее путём повторения операции смены направления текущего указателя на обратное получается граф, в котором изменены направления дуг для вершин сложения и вычитания. Теперь операции сложения и вычитания можно выполнять, граф редуцируется до окончательного результата.
\textit{(Источник: Орлов, Цилькер, гл. 15, с. 595--596.)}

\subsubsection{Сравнение потоковых и редукционных вычислений с фон-неймановской моделью}

\paragraph{Преимущества и недостатки.}
Потоковые и редукционные вычисления обладают рядом преимуществ по сравнению с традиционной фон-неймановской моделью:
\begin{itemize}
\item параллелизм выявляется автоматически, без явного указания программистом;
\item отсутствует счётчик команд, что упрощает управление;
\item естественным образом поддерживается функциональный стиль программирования;
\item в редукционных системах вычисления выполняются только по мере необходимости (ленивые вычисления), что может сократить объём бесполезной работы.
\end{itemize}
К недостаткам можно отнести:
\begin{itemize}
\item накладные расходы на пересылку токенов и согласование;
\item потерю локальности по обращению, что снижает эффективность кэш-памяти;
\item сложность эффективной реализации на современной элементной базе;
\item ограниченную применимость для задач с нерегулярной структурой.
\end{itemize}
Тем не менее идеи потоковых и редукционных вычислений оказали значительное влияние на развитие архитектуры параллельных ВС и языков программирования.
\textit{(Источник: Орлов, Цилькер, гл. 15, с. 579--597.)}

\subsubsection{Краткое резюме подраздела}
Машины, управляемые потоком данных, используют нетрадиционный механизм управления вычислениями: команда выполняется, когда становятся доступны её операнды. Вычислительная модель описывается графом потоков данных, вершины которого соответствуют операциям, а дуги — потокам данных (токенам). Статические потоковые системы допускают не более одного токена на дуге и не поддерживают параллельное выполнение итераций цикла; динамические системы (с помеченными или явно адресуемыми токенами) снимают это ограничение за счёт учёта контекста (цвета) токенов. Мультипотоковая обработка представляет собой крупнозернистую потоковую обработку, в которой вершины графа соответствуют последовательностям команд (нитям). Редукционные (управляемые по запросу) системы инициируют вычисления запросом на результат и реализуют ленивые вычисления; они используют строчную или графовую модель редукции. Потоковые и редукционные вычисления обладают как преимуществами (автоматический параллелизм, отсутствие счётчика команд), так и недостатками (накладные расходы, потеря локальности), и оказали значительное влияние на развитие параллельных архитектур.
\textit{(Обобщение по: Орлов, Цилькер, гл. 15, с. 579--597.)}

\subsection{Организация ввода-вывода, каналы и процессоры ввода-вывода, устройства сопряжения с объектами}

\noindent\textit{Основные источники: Орлов С.А., Цилькер Б.Я. Организация ЭВМ и систем: Учебник для вузов. 2-е изд. СПб.: Питер, 2011. Гл. 7, с. 332--363; гл. 8, с. 364--385.}

\subsubsection{Система ввода/вывода и её место в архитектуре ВМ}

\paragraph{Назначение системы ввода/вывода.}
Помимо центрального процессора (ЦП) и памяти, третьим ключевым элементом архитектуры ВМ является система ввода/вывода (СВВ), обеспечивающая обмен информацией между ядром ВМ и разнообразными периферийными устройствами (ПУ). Технические и программные средства СВВ несут ответственность за физическое и логическое сопряжение ядра вычислительной машины и ПУ. Технически система ввода/вывода в рамках ВМ реализуется совокупностью модулей ввода/вывода (МВВ). Модуль ввода/вывода выполняет сопряжение ПУ с ядром вычислительной машины и различные коммуникационные операции между ними. Две основные функции МВВ:
\begin{itemize}
\item обеспечение интерфейса с ЦП и памятью («большой» интерфейс);
\item обеспечение интерфейса с одним или несколькими периферийными устройствами («малый» интерфейс).
\end{itemize}
\textit{(Источник: Орлов, Цилькер, гл. 8, с. 364.)}

\paragraph{Способы подключения СВВ к ядру ВМ.}
Анализируя архитектуру известных ВМ, можно выделить три основных способа подключения СВВ к ядру вычислительной машины. В варианте с раздельными шинами памяти и ввода/вывода обмен информацией между ЦП и памятью физически отделён от ввода/вывода, поскольку обеспечивается полностью независимыми шинами. Это даёт возможность обращаться к памяти одновременно с выполнением ввода/вывода. Кроме того, данный архитектурный вариант позволяет специализировать каждую из шин, учесть формат пересылаемых данных, особенности синхронизации обмена и т.п. В частности, шина ввода/вывода, с учётом характеристик реальных ПУ, может иметь меньшую пропускную способность, что позволяет снизить затраты на её реализацию. Недостатком решения можно считать большое количество точек подключения к ЦП.
\textit{(Источник: Орлов, Цилькер, гл. 8, с. 364--365.)}

Второй вариант — с совместно используемыми линиями данных и адреса. Память и СВВ имеют общие для них линии адреса и линии данных, разделяя их во времени. В то же время управление памятью и СВВ, а также синхронизация их взаимодействия с процессором осуществляются независимо по раздельным линиям управления. Это позволяет учесть особенности процедур обращения к памяти и к модулям ввода/вывода и добиться наибольшей эффективности доступа к ячейкам памяти и периферийным устройствам.
\textit{(Источник: Орлов, Цилькер, гл. 8, с. 365.)}

Третий тип архитектуры ВМ предполагает подключение СВВ к системной шине на общих правах с процессором и памятью. Преимущества и недостатки такого подхода обсуждались при рассмотрении вопросов организации шин (глава 7). Потенциально возможен также вариант подключения периферийных устройств к системной шине напрямую, без использования МВВ, но против него можно выдвинуть сразу несколько аргументов. Во-первых, в этом случае ЦП пришлось бы оснащать универсальными схемами для управления любым ПУ; при большом разнообразии периферийных устройств, имеющих к тому же различные принципы действия, такие схемы оказываются чересчур сложными и избыточными. Во-вторых, пересылка данных при вводе и выводе происходит значительно медленнее, чем при обмене между ЦП и памятью, и было бы невыгодно задействовать для обмена информацией с ПУ высокоскоростную системную шину. И наконец, в ПУ часто используются иные форматы данных и длина слова, чем в ВМ, к которым они подключены.
\textit{(Источник: Орлов, Цилькер, гл. 8, с. 365.)}

\subsubsection{Адресное пространство системы ввода/вывода}

\paragraph{Совмещённое адресное пространство.}
Как и в случае памяти, операции ввода/вывода также предполагают наличие некоторой системы адресации, позволяющей выбрать один из модулей СВВ, а также одно из подключённых к нему периферийных устройств. Адрес модуля и ПУ является составной частью соответствующей команды. Адресное пространство ввода/вывода может быть совмещено с адресным пространством памяти или может быть выделенным. При совмещении адресного пространства для адресации модулей ввода/вывода отводится определённая область адресов. Обычно все операции с модулем ввода/вывода осуществляются с использованием входящих в него внутренних регистров: управления, состояния, данных. Фактически процедура ввода/вывода сводится к записи информации в одни регистры МВВ и считыванию её из других регистров. Это позволяет рассматривать регистры МВВ как ячейки основной памяти и работать с ними с помощью обычных команд обращения к памяти, при этом в системе команд ВМ вообще могут отсутствовать специальные команды ввода и вывода. Так, модификацию регистров МВВ можно производить непосредственно с помощью арифметических и логических команд. Адреса регистрам МВВ назначаются в области адресного пространства памяти, отведённой под систему ввода/вывода. Такой подход представляется вполне оправданным, если учесть, что ввод/вывод обычно составляет малую часть всех операций, выполняемых вычислительной машиной, чаще всего не более 1\% от общего числа команд в программе.
\textit{(Источник: Орлов, Цилькер, гл. 8, с. 365--366.)}

К достоинствам совмещённого адресного пространства можно отнести:
\begin{itemize}
\item унификацию операций обращения к памяти и периферийным устройствам;
\item возможность увеличения числа подключаемых периферийных устройств;
\item возможность внепроцессорного обмена данными между периферийными устройствами, если в системе команд есть команды пересылки между ячейками памяти;
\item возможность обмена информацией не только с аккумулятором, но и с любым регистром центрального процессора.
\end{itemize}
Недостатки совмещённого адресного пространства:
\begin{itemize}
\item сокращение области адресного пространства памяти;
\item усложнение декодирующих схем адресов в СВВ;
\item трудности распознавания операций передачи информации при вводе/выводе среди других операций.
\end{itemize}
Совмещённое адресное пространство используется в вычислительных машинах MIPS и SPARC.
\textit{(Источник: Орлов, Цилькер, гл. 8, с. 366.)}

\paragraph{Выделенное адресное пространство.}
В случае выделенного адресного пространства для обращения к модулям ввода/вывода применяются специальные команды и отдельная система адресов. Это позволяет разделить тракты для работы с памятью и тракты ввода/вывода, что даёт возможность совмещать во времени обмен с памятью и ввод/вывод. Кроме того, адресное пространство памяти может быть использовано по прямому назначению в полном объёме. В вычислительных машинах фирмы IBM и ВМ на базе процессоров фирмы Intel система ввода/вывода, как правило, организуется в соответствии с концепцией выделенного адресного пространства.
\textit{(Источник: Орлов, Цилькер, гл. 8, с. 366.)}

Достоинства выделенного адресного пространства:
\begin{itemize}
\item адрес периферийного устройства в команде ввода/вывода может быть коротким, благодаря чему сокращается длина таких команд и упрощаются схемы дешифрации адреса;
\item программы становятся более наглядными, так как операции ввода/вывода выполняются с помощью специальных команд;
\item разработка СВВ может проводиться отдельно от разработки памяти.
\end{itemize}
Недостатки выделенного адресного пространства:
\begin{itemize}
\item ввод/вывод производится только через аккумулятор центрального процессора;
\item малая разрядность адреса периферийного устройства ограничивает число подключаемых ПУ.
\end{itemize}
\textit{(Источник: Орлов, Цилькер, гл. 8, с. 367.)}

\subsubsection{Периферийные устройства}

\paragraph{Категории периферийных устройств.}
Связь ВМ с внешним миром осуществляется с помощью самых разнообразных периферийных устройств. Каждое ПУ подключается к модулю ввода/вывода посредством индивидуальной шины. Интерфейс, по которому организуется взаимодействие МВВ и ПУ, часто называют малым. Индивидуальная шина обеспечивает обмен данными и управляющими сигналами, а также информацией о состоянии участников обмена. Все множество ПУ можно свести к трём категориям:
\begin{enumerate}[label=\arabic*)]
\item для общения с пользователем;
\item для общения с ВМ;
\item для связи с удалёнными устройствами.
\end{enumerate}
Примерами первой группы служат видеотерминалы и принтеры. Ко второй группе причисляются внешние запоминающие устройства (магнитные и оптические диски, магнитные ленты и т.п.), датчики и исполнительные механизмы. Отметим двойственную роль внешних ЗУ, которые, с одной стороны, представляют собой часть памяти ВМ, а с другой — являются периферийными устройствами. Наконец, устройства третьей категории позволяют ВМ обмениваться информацией с удалёнными объектами, которые могут относиться к двум первым группам. В роли удалённых объектов могут выступать также другие ВМ.
\textit{(Источник: Орлов, Цилькер, гл. 8, с. 367.)}

\paragraph{Обобщённая структура периферийного устройства.}
Обобщённая структура ПУ включает интерфейс с МВВ, реализуемый в виде сигналов управления, состояния и данных. Данные представлены совокупностью битов, которые должны быть переданы в модуль ввода/вывода или получены из него. Сигналы управления определяют функцию, которая должна быть выполнена периферийным устройством. Это может быть стандартная для всех устройств функция — посылка данных в МВВ или получение данных из него, либо специфичная для данного типа ПУ функция, такая например, как позиционирование головки магнитного диска или перемотка магнитной ленты. Сигналы состояния характеризуют текущее состояние устройства, в частности включено ли ПУ и готово ли оно к передаче данных. Логика управления — это схемы, координирующие работу ПУ в соответствии с направлением передачи данных. Задачей преобразователя является трансформация информационных сигналов, имеющих самую различную физическую природу, в электрические сигналы, а также обратное преобразование. Обычно совместно с преобразователем используется буферная память, обеспечивающая временное хранение данных, пересылаемых между МВВ и ПУ.
\textit{(Источник: Орлов, Цилькер, гл. 8, с. 367--368.)}

\subsubsection{Модули ввода/вывода}

\paragraph{Функции модуля ввода/вывода.}
Модуль ввода/вывода в составе вычислительной машины отвечает за управление подключёнными к нему одним или несколькими ПУ и за обмен данными между этими устройствами, с одной стороны, и основной памятью или регистрами ЦП — с другой. Основные функции МВВ можно сформулировать следующим образом:
\begin{itemize}
\item локализация данных;
\item управление и синхронизация;
\item обмен информацией;
\item буферизация данных;
\item обнаружение ошибок.
\end{itemize}
\textit{(Источник: Орлов, Цилькер, гл. 8, с. 368.)}

\paragraph{Локализация данных.}
Под локализацией данных будем понимать возможность обращения к одному из ПУ, а также адресации данных на нём. Адрес ПУ обычно содержится в адресной части команд ввода/вывода. Как уже отмечалось, в состав СВВ могут входить несколько модулей ввода/вывода. К каждому МВВ может быть подключено несколько ПУ, поэтому каждому модулю назначается определённый диапазон адресов, независимо от того, является ли пространство адресов совмещённым или раздельным. Старшие разряды обеспечивают выбор одного из МВВ в рамках системы ввода/вывода. Младшие разряды адреса представляют собой уникальные адреса регистров данного модуля или подключённых к нему ПУ. Одной из функций МВВ является проверка «попадания» поступившего адреса в выделенный данному модулю диапазон адресов. При положительном ответе модуль должен обеспечить декодирование поступившего адреса и перенаправление информации к адресуемому объекту или от него. Для простейших периферийных устройств (клавиатура, принтер и т.п.) адрес ПУ однозначно определяет и расположение данных на этом устройстве. Для более сложных ПУ, таких как внешние запоминающие устройства, информация о местонахождении данных требует детализации. Так, для ЗУ на магнитном диске необходимо указать номер цилиндра, номер сектора и т.п. Эта часть адресной информации передаётся в МВВ не по шине адреса, а в виде служебных сообщений, пересылаемых по шине данных.
\textit{(Источник: Орлов, Цилькер, гл. 8, с. 368--369.)}

\paragraph{Управление и синхронизация.}
Функция управления и синхронизации заключается в том, что МВВ должен координировать перемещение данных между внутренними ресурсами ВМ и периферийными устройствами. При разработке системы управления и синхронизации модуля ввода/вывода необходимо учитывать целый ряд факторов. Прежде всего, нужно принимать во внимание, что центральный процессор может взаимодействовать одновременно с несколькими ПУ, причём быстродействие подключаемых к МВВ периферийных устройств варьируется в очень широких пределах — от нескольких байтов в секунду в терминалах до десятков миллионов байтов в секунду при обмене с магнитными дисками. Если в системе используются шины, каждое взаимодействие между ЦП и МВВ включает в себя одну или несколько процедур арбитража.
\textit{(Источник: Орлов, Цилькер, гл. 8, с. 369.)}

В отличие от обмена с памятью, процессы ввода/вывода и работа ЦП протекают не синхронно. Очередная порция информации может быть выдана на устройство вывода лишь тогда, когда это устройство готово их принять. Аналогично, ввод допустим только в случае доступности информации на устройстве ввода. Несинхронный характер процессов ввода/вывода предполагает обмен сигналами, который для двухпроводной системы синхронизации (применительно к операции вывода) можно описать следующим образом:
\begin{enumerate}[label=\arabic*)]
\item Центральный процессор с помощью сигнала ДД = 1 (данные достоверны) извещает о доступности данных, подлежащих выводу.
\item Приняв данные, устройство вывода сообщает процессору об их получении сигналом ДП = 1 (данные приняты).
\item Получив подтверждение, ЦП обнуляет сигнал ДД и снимает данные с шины, после чего может выставить на шину новые данные.
\item Обнаружив, что ДД = 0, устройство вывода, в свою очередь, устанавливает в нулевое состояние сигнал ДП, после чего оно готово для обработки принятых данных (и обрабатывает их до получения очередного сигнала ДД = 1).
\end{enumerate}
Таким образом, модуль ввода/вывода обязан снабдить центральный процессор информацией о собственной готовности к обмену, а также о готовности подключённых к модулю ПУ. Помимо этого, процессор должен обладать оперативной информацией и об иных событиях, происходящих в СВВ.
\textit{(Источник: Орлов, Цилькер, гл. 8, с. 369.)}

\paragraph{Обмен информацией.}
Основной функцией МВВ является обеспечение обмена информацией. Со стороны «большого» интерфейса — это обмен с ЦП, а со стороны «малого» интерфейса — обмен с ПУ. В таком плане требования к МВВ диктуются последовательностью операций, выполняемых процессором при вводе/выводе:
\begin{enumerate}[label=\arabic*)]
\item Выбор требуемого периферийного устройства.
\item Определение состояния МВВ и ПУ.
\item Выдача указания модулю ввода/вывода на подключение нужного ПУ к процессору.
\item Получение от МВВ подтверждения о подключении затребованного ПУ к процессору.
\item Распознавание сигнала готовности устройства к передаче очередной порции информации.
\item Приём (передача) порции информации.
\item Циклическое повторение двух предшествующих пунктов до завершения передачи информации в полном объёме.
\item Логическое отсоединение ПУ от процессора.
\end{enumerate}
С учётом описанной процедуры функция обмена информацией с ЦП включает в себя: распознавание команды (МВВ получает команды из ЦП в виде сигналов на шине управления); пересылку данных между МВВ и ЦП; извещение о состоянии (из-за того что ПУ — медленные устройства, важно знать состояние модуля ввода/вывода; так, в момент получения запроса на пересылку данных в центральный процессор МВВ может быть не готов выполнить эту пересылку, поскольку ещё не завершил предыдущую команду; этот факт должен быть сообщён процессору с помощью соответствующего сигнала; возможны также сигналы, уведомляющие о возникших ошибках); распознавание адреса (МВВ обязан распознавать адрес каждого ПУ, которым он управляет). Помимо обмена с процессором, МВВ должен обеспечивать функцию обмена информацией с ПУ. Такой обмен также включает в себя передачу данных, команд и информации о состоянии.
\textit{(Источник: Орлов, Цилькер, гл. 8, с. 370.)}

\paragraph{Буферизация данных.}
Несмотря на различия в скорости обмена информацией для разных ПУ, все они в этом плане значительно отстают от ЦП и памяти. Такое различие компенсируется за счёт буферизации. При выводе информации на ПУ данные пересылаются из основной памяти в МВВ с большой скоростью. В модуле эти данные буферизируются и затем направляются в ПУ со скоростью, свойственной последнему. При вводе из ПУ данные буферизируются так, чтобы не заставлять память работать в режиме медленной передачи. Таким образом, МВВ должен обеспечивать работу как со скоростью памяти, так и со скоростью ПУ.
\textit{(Источник: Орлов, Цилькер, гл. 8, с. 370.)}

\paragraph{Обнаружение ошибок.}
Ещё одной из важнейших функций МВВ является обнаружение ошибок, возникающих в процессе ввода/вывода. Центральный процессор следует оповещать о каждом случае обнаружения ошибки. Ошибки могут быть вызваны самыми разнообразными факторами:
\begin{itemize}
\item загрязнение и влага;
\item повышенная или пониженная температура окружающей среды;
\item электромагнитное облучение;
\item скачки напряжения питания;
\item старение элементной базы;
\item ошибки в системном программном обеспечении;
\item ошибки в пользовательском программном обеспечении.
\end{itemize}
Степень влияния каждого из этих факторов зависит от типа и конструкции МВВ и ПУ. Так, к загрязнению наиболее чувствительны оптические и механические элементы ПУ, в то время как работа электронных компонентов СВВ в большей степени зависит от температуры внешней среды, электромагнитного воздействия и стабильности питающего напряжения. Фактор старения характерен как для механических, так и для электронных элементов СВВ. В механических элементах он выражается в виде износа, следствием чего может быть неточное позиционирование головок считывания/записи на внешних запоминающих устройствах или неправильная подача бумаги в принтерах. Старение электронных элементов обычно выражается в изменении электрических параметров схем, приводящем к нарушению управления и синхронизации. Так, отклонения в параметрах электронных компонентов могут вызвать недопустимый «перекос» сигналов, передаваемых между ЦП и МВВ или внутри МВВ. Источником ошибок может стать и несовершенство системного программного обеспечения: непредвиденные последовательности команд или кодовые комбинации; некорректное распределение памяти; недостаточный размер буфера ввода/вывода; недостаточно продуманные и проверенные комбинации системных модулей. Среди ошибок, порождаемых пользовательским ПО, наиболее частыми являются нарушение последовательности выполнения программы и некорректные процедуры. Вероятность возникновения ошибки внутри процессора для современных ЦП оценивается величиной порядка \(10^{-18}\), в то время как для остальных составляющих ВМ она лежит в диапазоне \(10^{-8}\)--\(10^{-12}\). Способы обнаружения и исправления ошибок ввода/вывода практически не отличаются от рассмотренных применительно к памяти.
\textit{(Источник: Орлов, Цилькер, гл. 8, с. 371.)}

\paragraph{Структура модуля ввода/вывода.}
Структура МВВ в значительной мере зависит от числа и сложности периферийных устройств, которыми он управляет, однако в самом общем виде такой модуль можно представить в форме, показанной на рис. 8.4 (см. учебник). Связь модуля ввода/вывода с ядром ВМ осуществляется посредством системной или специализированной шины. С этой стороны МВВ реализуется так называемый «большой» интерфейс. Большие различия в архитектуре систем команд и шин ВМ являются причиной того, что со стороны «большого» интерфейса модули ввода/вывода достаточно трудно унифицировать, и часто МВВ, созданные для одних ВМ, не могут быть использованы в других. Тем не менее в структурном плане они достаточно схожи.
\textit{(Источник: Орлов, Цилькер, гл. 8, с. 372.)}

Данные, передаваемые между МВВ и ядром вычислительной машины, буферизируются в регистре данных, благодаря чему удаётся компенсировать различие в быстродействии ядра ВМ и периферийных устройств. Разрядность регистра, как правило, совпадает с количеством линий в шине, подключаемой со стороны «большого» интерфейса (обычно 4 или 8 байтов). В то же время большинство ПУ ориентировано на побайтовый обмен информацией. Побайтовая пересылка информации по «широкой» системной шине — крайне неэффективное решение, поэтому со стороны «малого» интерфейса регистр данных часто дополняют узлом упаковки/распаковки. Этот узел при вводе обеспечивает последовательное побайтовое заполнение регистра данных (упаковку), а при выводе — последовательную побайтовую выдачу содержимого регистра на ПУ (распаковку). В результате обмен данными через «большой» интерфейс выполняется за один такт. МВВ может содержать несколько регистров данных, что позволяет независимо хранить текущие данные каждого из периферийных устройств, подключённых к данному модулю ввода/вывода.
\textit{(Источник: Орлов, Цилькер, гл. 8, с. 372--373.)}

Помимо регистра данных, в составе МВВ имеются также регистр управления и регистр состояния (либо совмещённый регистр управления/состояния). В регистре управления (РУ) фиксируются поступившие из ЦП команды управления модулем или подключёнными к нему периферийными устройствами. Отдельные разряды регистра могут представлять такие команды, как очистка регистров МВВ, сброс ПУ, начало чтения, начало записи и т.п. В сложных МВВ присутствует несколько регистров управления, например регистр управляющих сигналов для модуля в целом и отдельные РУ для каждого из ПУ. Регистр состояния (РС) служит для хранения битов состояния МВВ и подключённых к нему ПУ. Содержимое определённого разряда регистра может характеризовать, например, готовность устройства ввода к приёму очередной порции данных, занятость устройства вывода или нахождение ПУ в автономном режиме (offline). В МВВ не исключается наличие и более одного регистра состояния.
\textit{(Источник: Орлов, Цилькер, гл. 8, с. 373.)}

Процедура ввода/вывода предполагает возможность работы с каждым регистром МВВ или периферийным устройством по отдельности. Такая возможность обеспечивается системой адресации. Каждому модулю в адресном пространстве ввода/вывода (совмещённом или раздельном) выделяется уникальный набор адресов, количество адресов в котором зависит от числа адресуемых элементов. Поступивший из ЦП адрес с помощью селектора адреса проверяется на принадлежность к диапазону, выделенному данному МВВ. В случае подтверждения дешифратор адреса выполняет раскодирование адреса, разрешая работу с соответствующим регистром модуля или ПУ. Узел управления вводом/выводом, по сути, играет роль местного устройства управления МВВ. На него возлагаются две задачи: обеспечение взаимодействия с ЦП и координация работы всех составляющих МВВ. Связь с ЦП реализуется посредством линий управления, по которым из ЦП в модуль поступают сигналы, служащие для синхронизации операций ввода и вывода. В обратном направлении передаются сигналы, информирующие о происходящих в модуле событиях, например сигналы прерывания. Часть линий управления может задействоваться модулем для арбитража. Вторая функция узла управления реализуется с помощью внутренних сигналов управления.
\textit{(Источник: Орлов, Цилькер, гл. 8, с. 373.)}

Со стороны «малого» интерфейса МВВ обеспечивает подключение периферийных устройств и взаимодействие с ними. Эта часть МВВ более унифицирована, поскольку периферийные устройства всегда «подгоняются» под один из стандартных протоколов. Каждое из периферийных устройств обслуживается своим узлом «малого» интерфейса, который реализует принятый для данного ПУ стандартный протокол взаимодействия. При управлении широким спектром ПУ модуль должен по возможности освобождать ЦП от знания деталей конкретных ПУ, так чтобы ЦП мог управлять любым устройством с помощью простых команд чтения и записи. МВВ при этом берёт на себя задачи синхронизации, согласования форматов данных и т.п. Модуль ввода/вывода, который берёт на себя детальное управление ПУ и общается с ЦП только с помощью команд высокого уровня, часто называют каналом ввода/вывода или процессором ввода/вывода. Наиболее примитивный МВВ, требующий детального управления со стороны ЦП, называют контроллером ввода/вывода или контроллером устройства. Как правило, контроллеры ввода/вывода типичны для микроЭВМ, а каналы ввода/вывода — для больших универсальных ВМ.
\textit{(Источник: Орлов, Цилькер, гл. 8, с. 373--374.)}

\subsubsection{Методы управления вводом/выводом}

\paragraph{Общая классификация.}
В ВМ находят применение три способа организации ввода/вывода (В/ВЫВ):
\begin{itemize}
\item ввод/вывод с опросом;
\item ввод/вывод по прерываниям;
\item прямой доступ к памяти.
\end{itemize}
В первых двух способах все действия, связанные с процедурой ввода/вывода, выполняются под полным контролем центрального процессора, причём информация, пересылаемая между ПУ и памятью, сначала попадает в процессор, после чего он передаёт эту информацию в пункт назначения.
\textit{(Источник: Орлов, Цилькер, гл. 8, с. 374.)}

\paragraph{Ввод/вывод с опросом.}
Наиболее простым методом управления вводом/выводом является ввод/вывод с опросом. Здесь ввод/вывод происходит под полным контролем центрального процессора и реализуется специальной процедурой ввода/вывода. ЦП с помощью команды ввода/вывода сообщает модулю ввода/вывода, а через него и периферийному устройству о предстоящей операции. Адрес модуля и ПУ, к которому производится обращение, указывается в адресной части команды ввода или вывода. Модуль исполняет затребованное действие, после чего устанавливает в единицу соответствующий бит в своём регистре состояния. Ничего другого, чтобы уведомить ЦП, модуль не предпринимает. Следовательно, для определения момента завершения операции или пересылки очередного элемента блока данных процессор должен периодически опрашивать и анализировать содержимое регистра состояния МВВ.
\textit{(Источник: Орлов, Цилькер, гл. 8, с. 374.)}

Иллюстрация процедуры ввода с опросом блока данных с устройства ввода приведена на рис. 8.5 (см. учебник). Данные читаются пословно. Для каждого читаемого слова ЦП должен оставаться в цикле проверки, пока не определит, что слово находится в регистре данных МВВ, то есть доступно для считывания. Процедура начинается с выдачи процессором команды ввода, в которой указан адрес конкретного МВВ и конкретного ПУ. Существуют четыре типа команд ввода/вывода, которые может получить МВВ: управление, проверка, чтение и запись.
\begin{itemize}
\item \textbf{Команды управления} используются для активизации ПУ и указания требуемой операции. Например, в устройство памяти на магнитной ленте может быть выдана команда перемотки или продвижения на одну запись. Для каждого типа ПУ характерны специфичные для него команды управления.
\item \textbf{Команда проверки} применяется для проверки различных ситуаций, возникающих в МВВ и ПУ в процессе ввода/вывода. С помощью таких команд ЦП способен выяснить, включено ли ПУ, готово ли оно к работе, завершена ли последняя операция ввода/вывода и не возникли ли в ходе её выполнения какие-либо ошибки. Действие команды сводится к установке или сбросу соответствующих разрядов регистра состояния МВВ.
\item \textbf{Команда чтения} побуждает модуль получить элемент данных из ПУ и занести его в регистр данных (РД). ЦП может получить этот элемент данных, запросив МВВ поместить его на шину данных.
\item \textbf{Команда записи} заставляет модуль принять элемент данных (байт или слово) с шины данных и переслать его в РД с последующей передачей в ПУ.
\end{itemize}
Если к МВВ подключено несколько ПУ, то в процедуре ввода/вывода нужно производить циклический опрос всех устройств, с которыми в данный момент выполняются операции В/ВЫВ.
\textit{(Источник: Орлов, Цилькер, гл. 8, с. 374--375.)}

Основной недостаток рассматриваемого способа ввода/вывода — неэффективное использование процессора. Для определения момента готовности МВВ к пересылке очередной порции информации ЦП должен циклически опрашивать модуль на наличие сигнала готовности и никаких иных полезных действий в это время не выполняет. Кроме того, пересылка даже одного слова требует выполнения нескольких команд. ЦП должен тратить время на анализ битов состояния МВВ, запись в МВВ битов управления, чтение или запись данных со скоростью, определяемой периферийным устройством. Всё это также отрицательно сказывается на эффективности ввода/вывода. Главным аргументом в пользу ввода/вывода с опросом является простота МВВ, поскольку основные функции по управлению вводом/выводом берёт на себя процессор. При одновременной работе с несколькими ПУ приоритет устройств легко изменить программными средствами (последовательностью опроса). Наконец, подключение к МВВ новых периферийных устройств или отключение ранее подключённых также реализуется без особых сложностей.
\textit{(Источник: Орлов, Цилькер, гл. 8, с. 375--376.)}

\paragraph{Ввод/вывод по прерываниям.}
Альтернативой предыдущему способу может быть вариант, когда ЦП выдаёт команду ввода/вывода, а затем продолжает делать другую полезную работу, не занимаясь анализом готовности ПУ. Когда ПУ готово к обмену данными, оно через МВВ извещает об этом процессор с помощью запроса на прерывание. ЦП осуществляет передачу очередного элемента данных, после чего возобновляет выполнение прерванной программы. Подобную возможность предоставляет ранее рассмотренная система прерывания программ вычислительной машины.
\textit{(Источник: Орлов, Цилькер, гл. 8, с. 376.)}

Рассмотрим процедуру ввода блока данных с использованием ввода/вывода по прерываниям (рис. 8.6, см. учебник). Оставим без внимания такие подробности, как сохранение и восстановление контекста, действия, выполняемые при завершении пересылки блока данных, а также в случае возникновения ошибок. Процедура реализуется следующим образом. ЦП выдаёт команду чтения, а затем продолжает выполнение других заданий, например другой программы. Получив команду, МВВ приступает к вводу элемента данных с ПУ. Когда считанное слово оказывается в регистре данных модуля, МВВ формирует на управляющей линии сигнал прерывания ЦП. Выставив запрос, МВВ помещает введённую информацию на шину данных, после чего он готов к следующей операции ввода/вывода. ЦП в конце каждого цикла команды проверяет наличие запросов прерывания. Когда от МВВ приходит такой сигнал, ЦП сохраняет контекст текущей программы и обрабатывает прерывание. В рассматриваемом случае ЦП читает слово из модуля, записывает его в память и выдаёт модулю команду на считывание очередного слова. Далее ЦП восстанавливает контекст прерванной программы и возобновляет её выполнение. Этот метод эффективнее ввода/вывода с опросом, поскольку устраняет ненужные ожидания, однако обработка прерывания занимает достаточно много времени ЦП. Кроме того, каждое слово, пересылаемое из памяти в модуль В/ВЫВ или в противоположном направлении, как и при методе с опросом, проходит через ЦП.
\textit{(Источник: Орлов, Цилькер, гл. 8, с. 376--377.)}

\paragraph{Прямой доступ к памяти.}
Хотя ввод/вывод по прерываниям эффективнее ввода/вывода с опросом, оба этих метода страдают двумя недостатками: темп передачи при вводе/выводе ограничен скоростью, с которой ЦП в состоянии опросить и обслужить устройство; ЦП вовлечён в управление передачей, для каждой пересылки он должен выполнить определённое количество команд. Когда пересылаются большие объёмы данных, требуется более эффективный способ ввода/вывода — прямой доступ к памяти (ПДП). ПДП предполагает наличие на системной шине дополнительного модуля — контроллера прямого доступа к памяти (КПДП), способного брать на себя функции ЦП по управлению системной шиной и обеспечивать прямую пересылку информации между ОП и ПУ без участия центрального процессора. В сущности, КПДП — это и есть модуль ввода/вывода, реализующий режим прямого доступа к памяти.
\textit{(Источник: Орлов, Цилькер, гл. 8, с. 377.)}

Если ЦП желает прочитать или записать блок данных, он прежде всего должен поместить в КПДП информацию, характеризующую предстоящее действие. Этот процесс называется инициализацией КПДП и включает в себя занесение в контроллер следующих четырёх параметров:
\begin{itemize}
\item вид запроса (чтение или запись);
\item адрес устройства ввода/вывода;
\item адрес начальной ячейки блока памяти, откуда будет извлекаться или куда будет вводиться информация;
\item количество слов, подлежащих чтению или записи.
\end{itemize}
Первый параметр определяет направление пересылки данных: из ОП в ПУ или наоборот. За исходную точку обычно принимается память, поэтому под чтением понимают считывание данных из ОП и выдачу их на устройство вывода, а под записью — приём данных из устройства ввода и запись в ОП. Вид запроса запоминается в схеме логики управления контроллера. К КПДП обычно могут быть подключены несколько ПУ, а адрес УВВ конкретизирует, какое из них должно участвовать в предстоящем обмене данными. Этот адрес запоминается в логике управления КПДП.
\textit{(Источник: Орлов, Цилькер, гл. 8, с. 377.)}

Третий параметр — адрес начальной ячейки — хранится в регистре адреса (РА) контроллера. После передачи каждого слова содержимое РА автоматически увеличивается на единицу, то есть в нём формируется адрес следующей ячейки ОП. Размер блока в словах заносится в счётчик данных (СД) контроллера. После передачи каждого слова содержимое СД автоматически уменьшается на единицу. Нулевое состояние СД свидетельствует о том, что пересылка блока данных завершена.
\textit{(Источник: Орлов, Цилькер, гл. 8, с. 378.)}

После инициализации процесс пересылки информации может быть начат в любой момент. Инициаторами обмена вправе выступать как ЦП, так и ПУ. Устройство, желающее начать В/ВЫВ, извещает об этом контроллер подачей соответствующего сигнала. Получив такой сигнал, КПДП выдаёт в ЦП сигнал Запрос ПДП. В ответ ЦП освобождает шины адреса и данных, а также те линии шины управления, по которым передаются сигналы, управляющие операциями на шине адреса (ША) и шине данных (ШД). К таким, прежде всего, относятся линии ЧТЗУ, ЗПЗУ, ВЫВ, ВВ и линия выдачи адреса на ША. Далее ЦП отвечает контроллеру сигналом Подтверждение ПДП, который для последнего означает, что ему делегированы права на управление системной шиной и можно приступать к пересылке данных.
\textit{(Источник: Орлов, Цилькер, гл. 8, с. 378.)}

Процесс пересылки каждого слова блока состоит из двух этапов. При выполнении операции чтения (ОП \(\to\) ПУ) на первом этапе КПДП выставляет на шину адреса содержимое РА (адрес текущей ячейки ОП) и формирует сигнал ЧТЗУ. Считанное из ячейки ОП слово помещается на шину данных. На втором этапе КПДП выставляет на ША адрес устройства вывода и формирует сигнал ВЫВ, который обеспечивает передачу слова с шины данных в ПУ. При выполнении операции записи (ПУ \(\to\) ОП) КПДП сначала выдаёт на шину адреса номер (адрес) устройства ввода и формирует сигнал ВВ, по которому введённые данные поступают на шину данных. На втором этапе КПДП помещает на ША адрес ячейки ОП, куда должны быть занесены данные, и выдаёт сигнал ЗПЗУ. Этим сигналом информация с ШД записывается в ячейку ОП.
\textit{(Источник: Орлов, Цилькер, гл. 8, с. 378--379.)}

Как при чтении, так и при записи происходит буферизация пересылаемого слова в регистре данных (РД) контроллера. Это необходимо для компенсации различий в скорости работы ОП и ПУ, в силу чего сигналы ВЫВ и ВВ формируются контроллером лишь при получении от ПУ подтверждения о готовности. Буферизация сводится к тому, что после первого этапа слово с ШД заносится в РД, а перед вторым — возвращается из РД на шину данных. После пересылки каждого слова логика управления прибавляет единицу к содержимому РА (формирует адрес следующей ячейки ОП) и уменьшает на единицу содержимое СД (ведёт подсчёт переданных слов). Когда пересылка завершена (при нулевом значении в СД), КПДП снимает сигнал Запрос ПДП, в ответ на что ЦП снимает сигнал Подтверждение ПДП и вновь берёт на себя управление системной шиной, то есть ЦП вовлечён в процесс ввода/вывода только в начале и конце передачи.
\textit{(Источник: Орлов, Цилькер, гл. 8, с. 379.)}

Эффективность ПДП зависит от организации совместного использования системной шины при пересылке блока. Здесь может применяться один из трёх режимов:
\begin{itemize}
\item \textbf{блочная пересылка}: КПДП полностью захватывает системную шину с момента начала пересылки и до момента завершения передачи всего блока; на весь этот период ЦП не имеет доступа к шине;
\item \textbf{режим пропуска цикла}: КПДП после передачи каждого слова на один цикл шины освобождает системную шину, предоставляя её на это время процессору; поскольку КПДП всё равно должен ждать готовности ПУ, это позволяет ЦП эффективно распорядиться данным обстоятельством;
\item \textbf{прозрачный режим}: КПДП имеет доступ к системной шине только в тех циклах, когда ЦП в ней не нуждается; это обеспечивает наиболее эффективную работу процессора, но может существенно замедлять операцию пересылки блока данных.
\end{itemize}
В отличие от обычного прерывания, в пределах цикла команды имеется несколько точек, где КПДП вправе захватить шину (рис. 8.8, см. учебник). Отметим, что это не прерывание: процессору не нужно запоминать контекст задачи.
\textit{(Источник: Орлов, Цилькер, гл. 8, с. 379--380.)}

Механизм ПДП может быть реализован различными путями. В первом варианте все ПУ совместно используют общую системную шину; КПДП работает как заменитель ЦП, и обмен данными между памятью и ПУ производится как ввод/вывод с опросом. Хотя этот вариант может быть достаточно дешёвым, эффективность его невысока: каждая пересылка требует двух циклов шины. Число необходимых циклов шины может быть уменьшено при объединении функций КПДП и ПУ: между КПДП и одним или несколькими ПУ есть другой тракт, не включающий системную шину. Логика ПДП может быть частью ПУ, либо это может быть отдельный КПДП, управляющий одним или несколькими периферийными устройствами. Допустим и ещё один шаг в том же направлении — соединение КПДП с ПУ посредством шины ввода/вывода. Это уменьшает число интерфейсов В/ВЫВ в КПДП и делает такую конфигурацию легко расширяемой. В двух последних вариантах системная шина задействуется КПДП только для обмена данными с памятью; обмен данными между КПДП и ПУ реализуется минуя системную шину.
\textit{(Источник: Орлов, Цилькер, гл. 8, с. 380.)}

\subsubsection{Каналы и процессоры ввода/вывода}

\paragraph{Назначение каналов.}
По мере развития систем В/ВЫВ их функции усложняются. Главная цель такого усложнения — максимальное высвобождение ЦП от управления процессами ввода/вывода. Некоторые пути решения этой задачи уже были рассмотрены. Следующими шагами в преодолении проблемы могут быть:
\begin{enumerate}[label=\arabic*)]
\item Расширение возможностей МВВ и предоставление ему прав процессора со специализированным набором команд, ориентированных на операции ввода/вывода. ЦП даёт указание такому процессору В/ВЫВ выполнить хранящуюся в памяти ВМ программу ввода/вывода. Процессор В/ВЫВ извлекает и исполняет команды этой программы без участия центрального процессора и прерывает ЦП только после завершения всей программы ввода/вывода.
\item Рассмотренному в пункте 1 процессору ввода/вывода придаётся собственная локальная память, при этом возможно управление множеством устройств В/ВЫВ с минимальным привлечением ЦП.
\end{enumerate}
В первом случае МВВ называют каналом ввода/вывода (КВВ), а во втором — процессором ввода/вывода. В принципе, различие между каналом и процессором ввода/вывода достаточно условно, поэтому в дальнейшем будем пользоваться термином «канал».
\textit{(Источник: Орлов, Цилькер, гл. 8, с. 380--381.)}

Концепция системы ввода/вывода с КВВ характерна для больших универсальных вычислительных машин (мэйнфреймов), где проблема эффективной организации ввода/вывода и максимального высвобождения центрального процессора в пользу его основной функции стоит наиболее остро. СВВ с каналами ввода/вывода была предложена и реализована в ВМ семейства IBM 360 и получила дальнейшее развитие в семействах IBM 370 и IBM 390.
\textit{(Источник: Орлов, Цилькер, гл. 8, с. 381.)}

\paragraph{Канальная программа и управляющие слова канала.}
В ВМ с каналами ввода/вывода центральный процессор практически не участвует в непосредственном управлении периферийными устройствами, делегируя эту задачу специализированному процессору, входящему в состав КВВ. Все функции ЦП сводятся к запуску и остановке операций в КВВ, а также проверке состояния канала и подключённых к нему ПУ. Для этих целей ЦП использует лишь несколько (от 4 до 8) команд ввода/вывода. Например, в IBM 360 таких команд четыре: «Начать ввод-вывод»; «Остановить ввод-вывод»; «Проверить ввод-вывод»; «Проверить канал».
\textit{(Источник: Орлов, Цилькер, гл. 8, с. 381.)}

КВВ реализует операции ввода/вывода путём выполнения так называемой канальной программы. Канальные программы для каждого ПУ, с которым предполагается обмен информацией, описывают нужную последовательность операций ввода-вывода и хранятся в основной памяти ВМ либо локальной памяти канала. Роль команд в канальных программах выполняют управляющие слова канала (УСК), структура которых отличается от структуры обычной машинной команды. Типовое УСК содержит:
\begin{itemize}
\item \textbf{код операции}, определяющий для КВВ и ПУ тип операции: «Записать» (вывод информации из ОП в ПУ), «Прочитать» (ввод информации из ПУ в ОП), «Управление» (перемещение головок НМД, магнитной ленты и т.п.);
\item \textbf{указатели} — дополнительные предписания, задающие более сложную последовательность операций В/ВЫВ, например при вводе пропускать отдельные записи или наоборот — с помощью одной команды вводить «разбросанный» по ОП массив как единый;
\item \textbf{адрес данных}, указывающий область памяти, используемую в операции ввода-вывода;
\item \textbf{счётчик данных}, хранящий значение длины передаваемого блока данных.
\end{itemize}
Кроме того, в УСК может содержаться идентификатор ПУ и информация о его уровне приоритета, указания по действиям, которые должны быть произведены при возникновении ошибок и т.п.
\textit{(Источник: Орлов, Цилькер, гл. 8, с. 381--382.)}

\paragraph{Иерархия управления вводом/выводом.}
Центральный процессор инициирует ввод/вывод, оповещая канал о необходимости выполнить канальную программу, находящуюся в ОП, и указывая начальный адрес этой программы в памяти ВМ. КВВ следует этим указаниям и управляет пересылкой данных. Пересылка информации каналом ведётся непосредственно между периферийным устройством (ПУ) и основной памятью, минуя центральный процессор. Таким образом, в вычислительной машине с КВВ управление вводом/выводом строится иерархическим образом. В операциях ввода/вывода участвуют три типа устройств:
\begin{itemize}
\item процессор (первый уровень управления);
\item канал ввода/вывода (второй уровень);
\item периферийное устройство (третий уровень).
\end{itemize}
Каждому типу устройств соответствует свой вид управляющей информации: процессору — команды ввода/вывода; каналу — управляющие слова канала; периферийному устройству — приказы.
\textit{(Источник: Орлов, Цилькер, гл. 8, с. 382.)}

\paragraph{Структура ВМ с канальной системой ввода/вывода.}
Обмен информацией между КВВ и основной памятью осуществляется посредством системной шины ВМ. ПУ подключаются к каналу не непосредственно, а через блоки управления периферийными устройствами (БУПУ). БУПУ принимает от канала приказы по управлению периферийным устройством (чтение, запись, перемещение носителя или магнитной головки и т.п.) и преобразует их в сигналы управления, свойственные данному типу ПУ. Обычно один БУПУ может обслуживать несколько однотипных ПУ, но для подключения быстродействующих периферийных устройств часто применяются индивидуальные блоки управления. В свою очередь, некоторые ПУ могут подключаться одновременно к нескольким БУПУ. Это позволяет воспользоваться свободным трактом другого БУПУ при занятости данного БУПУ обслуживанием одного из подключённых к нему ПУ. Физически БУПУ может быть самостоятельным устройством или интегрирован с ПУ или каналом.
\textit{(Источник: Орлов, Цилькер, гл. 8, с. 383.)}

Обмен информацией между БУПУ и КВВ обеспечивается так называемыми канальными трактами. Обычно каждое БУПУ связано с одним из канальных трактов, но возможно также подключение блока управления сразу к нескольким трактам, что даёт возможность избежать нежелательных задержек при занятости одного из них. В пределах канала ввода/вывода считается, что каждое ПУ подключено к своему подканалу. Подканалы имеют свои уникальные логические номера, с помощью которых канальная программа адресуется к конкретному ПУ. Физически подканал реализуется в виде участка памяти, в котором хранятся параметры операции ввода/вывода, выполняемой данным ПУ: текущие значения адреса и счётчика данных, код и указатели операции ввода/вывода, адрес следующего УСК и др. Для хранения этих параметров обычно используется локальная память канала.
\textit{(Источник: Орлов, Цилькер, гл. 8, с. 383.)}

Обмен информацией между ПУ и ОП, как уже упоминалось, реализуется в режиме прямого доступа к памяти, при этом для взаимодействия ЦП и канала задействованы сигналы Запрос ПДП и Подтверждение ПДП. Чтобы известить ЦП об окончании текущей канальной программы или об ошибках, возникших при её выполнении, КВВ выдаёт в ЦП сигнал Запрос прерывания. В свою очередь, ЦП может привлечь внимание канала сигналом Внимание.
\textit{(Источник: Орлов, Цилькер, гл. 8, с. 383.)}

\paragraph{Мультиплексный и монопольный режимы.}
Способ организации взаимодействия ПУ с каналом определяется соотношением быстродействия ОП и ПУ. По этому признаку ПУ образуют две группы: быстродействующие (накопители на магнитных дисках (НМД), накопители на магнитных лентах (НМЛ)) со скоростью приёма и выдачи информации около 1 Мбайт/с и медленнодействующие (дисплеи, печатающие устройства и др.) со скоростями порядка 1 Кбайт/с и менее. Быстродействие основной памяти обычно значительно выше. С учётом производительности ПУ в КВВ реализуются два режима работы: мультиплексный (режим разделения времени) и монопольный.
\textit{(Источник: Орлов, Цилькер, гл. 8, с. 383.)}

В мультиплексном режиме несколько периферийных устройств разделяют канал во времени, при этом каждое из параллельно работающих с каналом ПУ связывается с КВВ на короткие промежутки времени только после того, как ПУ будет готово к приёму или выдаче очередной порции информации (байта, группы байтов и т.д.). Такая схема принята в мультиплексном канале ввода/вывода. Если в течение сеанса связи пересылается один байт или несколько байтов, образующих одно машинное слово, канал называется байт-мультиплексным. Канал, в котором в пределах сеанса связи пересылка данных выполняется поблочно, носит название блок-мультиплексного. В монопольном режиме после установления связи между каналом и ПУ последнее монополизирует канал на всё время до завершения инициированной процессором канальной программы и всех предусмотренных этой программой пересылок данных между ПУ и ОП. На всё время выполнения канальной программы канал оказывается недоступным для других ПУ. Данную процедуру обеспечивает селекторный канал ввода/вывода. Отметим, что в блок-мультиплексном канале в рамках сеанса связи пересылка блока осуществляется в монопольном режиме.
\textit{(Источник: Орлов, Цилькер, гл. 8, с. 383--384.)}

\paragraph{Канальная подсистема.}
В последовавших за IBM 360 семействах универсальных ВМ IBM 370 и особенно IBM 390 концепция системы ввода/вывода на базе каналов получила дальнейшее развитие и вылилась в так называемую канальную подсистему ввода/вывода (КПВВ). Главная идея заключается в интегрировании отдельных КВВ в единый специализированный процессор ввода/вывода с большим числом канальных трактов и подканалов. Блоки управления периферийными устройствами обычно подключаются к нескольким канальным трактам, что позволяет динамически менять путь пересылки информации с учётом текущей их загруженности. Кроме того, разные канальные тракты могут обладать различной пропускной способностью и при выборе трактов для подключения определённых ПУ может быть учтено их быстродействие.
\textit{(Источник: Орлов, Цилькер, гл. 8, с. 384.)}

Одной из наиболее совершенных канальных подсистем обладают ВМ семейства IBM 390. В ней предусмотрено использование до 65\,536 подканалов и до 256 канальных трактов. Реализованы два типа канальных трактов: параллельный и последовательный. Параллельные канальные тракты по своим возможностям и принципу действия аналогичны рассмотренным ранее мультиплексному и селекторному каналам, но в отличие от них являются универсальными, то есть могут работать в байт-мультиплексном, блок-мультиплексном и монопольном режимах. Такие канальные тракты в КПВВ называют параллельными, поскольку они обеспечивают пересылку информации параллельным кодом. Для работы с ПУ, соединёнными с КПВВ волоконно-оптическими линиями, используются последовательные канальные тракты, реализующие протокол ESCON (Enterprise Systems Connection Architecture). Последовательный канальный тракт рассчитан на передачу информации только в последовательном коде и только в монопольном режиме. Для подключения блоков управления периферийными устройствами (БУПУ) к ESCON-тракту служат специальные устройства, называемые ESCON-директорами. Каждое такое устройство может обеспечить одновременное подключение до 60 БУПУ и одновременную передачу информации от 30 из них со скоростью до 10 Мбайт/с. Кроме того, в КПВВ предусмотрены специальные коммуникационные канальные тракты для подключения к сетям ВМ, модемам, другим системам.
\textit{(Источник: Орлов, Цилькер, гл. 8, с. 384--385.)}

В принципе, основное преимущество КПВВ — динамическое перераспределение канальных трактов — в какой-то мере может быть реализовано и в рамках каждого отдельного канала. Однако объединение всех канальных ресурсов в единую канальную подсистему позволяет применить оптимальную стратегию динамического распределения и использования этих ресурсов и благодаря этому достичь качественно нового уровня эффективности системы ввода/вывода.
\textit{(Источник: Орлов, Цилькер, гл. 8, с. 385.)}

\subsubsection{Устройства сопряжения с объектами}

\paragraph{Назначение и состав.}
Устройства сопряжения с объектами обеспечивают связь ВМ с датчиками, исполнительными механизмами и другими объектами управления. Они выполняют преобразование физических величин в электрические сигналы и обратно (аналого-цифровые и цифро-аналоговые преобразователи), дискретный ввод/вывод, гальваническую развязку, усиление и формирование сигналов, буферизацию данных, формирование запросов прерывания и взаимодействие с системной шиной или шиной ввода/вывода. В состав таких устройств могут входить порты ввода/вывода, адаптеры, блоки управления периферийными устройствами (БУПУ), драйверы, регистры данных и состояния, схемы синхронизации.
\textit{(Источник: Орлов, Цилькер, гл. 7, с. 332--363; гл. 8, с. 364--385.)}

\paragraph{Порты ввода/вывода.}
Портом называют схему, ответственную за передачу информации из периферийного устройства ввода в аккумулятор АЛУ (порт ввода) или из аккумулятора на периферийное устройство вывода (порт вывода). Схема обеспечивает электрическое и логическое сопряжение ВМ с подключённым к нему периферийным устройством. В модуле ввода/вывода рассматриваемой ВМ предполагается, что каждое ПУ подключается к своему порту. Каждый порт имеет уникальный номер, который указывается в адресной части команд ввода/вывода. Дешифратор номера порта ввода/вывода (ДВВ) обеспечивает преобразование номера порта в сигнал, разрешающий операцию ввода или вывода в соответствующем порту. Непосредственно ввод (вывод) происходит при поступлении из МПА сигнала Вв или Выв.
\textit{(Источник: Орлов, Цилькер, гл. 3, с. 123--124; гл. 8, с. 368.)}

\paragraph{Блоки управления периферийными устройствами.}
БУПУ принимает от канала приказы по управлению периферийным устройством (чтение, запись, перемещение носителя или магнитной головки и т.п.) и преобразует их в сигналы управления, свойственные данному типу ПУ. Обычно один БУПУ может обслуживать несколько однотипных ПУ, но для подключения быстродействующих периферийных устройств часто применяются индивидуальные блоки управления. В свою очередь, некоторые ПУ могут подключаться одновременно к нескольким БУПУ. Это позволяет воспользоваться свободным трактом другого БУПУ при занятости данного БУПУ обслуживанием одного из подключённых к нему ПУ. Физически БУПУ может быть самостоятельным устройством или интегрирован с ПУ или каналом.
\textit{(Источник: Орлов, Цилькер, гл. 8, с. 383.)}

\paragraph{Адаптеры и интерфейсы.}
Для подключения периферии и объектов управления используются системные шины, шины ввода/вывода и стандартные интерфейсы. Среди шин «малого» интерфейса наибольшее распространение получили USB, FireWire, Bluetooth, IrDA. Шина USB (Universal Serial Bus) ориентирована на среднескоростные и низкоскоростные периферийные устройства; для подключения ПУ используется четырёхпроводный кабель, при этом два провода используются для приёма и передачи данных, а два оставшихся — для питания ПУ. К одному контроллеру шины USB можно подсоединить до 127 устройств по топологии «звезда». FireWire — стандартная технология шины последовательной передачи данных для соединения компьютера с периферией; ПУ подключаются с помощью 6-проводной линии, где 4 провода используются для передачи информации, а два — для подачи питания. Интерфейс Bluetooth — технология беспроводной передачи данных, работающая на частоте 2,4 ГГц; в зависимости от степени мощности устанавливается связь в пределах 10 или 100 метров. Интерфейс IrDA (Infrared Data Association) позволяет соединяться с периферийным оборудованием без кабеля при помощи инфракрасного излучения с длиной волны 880 нм на коротком расстоянии до одного метра в режиме «точка-точка».
\textit{(Источник: Орлов, Цилькер, гл. 7, с. 360--362; гл. 8, с. 364--385.)}

\paragraph{Сопряжение с объектами управления.}
В системах управления ВМ взаимодействует с датчиками и исполнительными механизмами. Устройства сопряжения выполняют функции аналого-цифрового и цифро-аналогового преобразования, дискретного ввода/вывода, гальванической развязки, формирования сигналов управления и приёма сигналов состояния. Они подключаются к системной шине или шине ввода/вывода через стандартные интерфейсы и обеспечивают автоматизированный обмен данными между объектом управления и вычислительной системой. Такие устройства могут быть выполнены в виде отдельных плат расширения, модулей ввода/вывода или специализированных контроллеров.
\textit{(Источник: Орлов, Цилькер, гл. 7, с. 332--363; гл. 8, с. 364--385.)}

\subsubsection{Краткое резюме подраздела}
Система ввода/вывода связывает ядро ВМ с периферийными устройствами и реализуется совокупностью модулей ввода/вывода. Адресное пространство ввода/вывода может быть совмещённым с памятью или выделенным. Модули ввода/вывода выполняют локализацию данных, управление и синхронизацию, обмен информацией, буферизацию и обнаружение ошибок. Основные методы управления вводом/выводом — опрос, прерывания и прямой доступ к памяти (ПДП). ПДП освобождает процессор от участия в пересылке каждого слова; используются режимы блочной пересылки, пропуска цикла и прозрачный режим. В больших ВМ применяются каналы и процессоры ввода/вывода, выполняющие канальные программы из управляющих слов канала; управление вводом/выводом становится иерархическим (процессор — канал — периферийное устройство). Мультиплексные и селекторные каналы обеспечивают соответственно разделение времени и монопольное обслуживание периферийных устройств. Канальная подсистема объединяет канальные тракты и подканалы и позволяет динамически перераспределять ресурсы ввода/вывода. Устройства сопряжения с объектами (порты, БУПУ, адаптеры, интерфейсы) обеспечивают подключение периферии, датчиков и исполнительных механизмов к вычислительной системе.
\textit{(Обобщение по: Орлов, Цилькер, гл. 7, с. 332--363; гл. 8, с. 364--385.)}

\subsection{Краткое резюме}

Современная ВМ строится на принципах фон Неймана: двоичное кодирование, программное управление, однородность и адресуемость памяти. Архитектура системы команд определяет видимые программисту средства: команды, форматы данных, способы адресации и регистры; исторически выделяются CISC-, RISC-, VLIW- и EPIC-подходы. Память организуется иерархически; кэш-память сокращает среднее время доступа за счёт локальности, а виртуальная память обеспечивает иллюзию большого адресного пространства и защиту. Страничная организация использует страницы фиксированного размера и страничные таблицы; сегментная организация --- логические сегменты переменной длины; сегментно-страничная схема сочетает оба подхода. Конвейеризация команд и арифметических операций, суперскалярная обработка, внеочередное выполнение, переименование регистров и предсказание переходов повышают производительность процессора. Векторные и SIMD-команды обрабатывают группы данных одной командой. Специализированные процессоры (матричные, ассоциативные, систолические, процессоры ввода/вывода) ориентированы на определённые классы задач. Потоковые и редукционные машины используют нетрадиционные механизмы управления вычислениями: выполнение по готовности данных или по запросу результата. Система ввода/вывода связывает ядро ВМ с периферией; для этого применяются опрос, прерывания, прямой доступ к памяти, каналы и процессоры ввода/вывода, а также устройства сопряжения с объектами.
\textit{(Обобщение по: Орлов, Цилькер, гл. 1--2, 6, 8--9, 13, 15.)}

\noindent\textit{Источники к разделу:}
\begin{itemize}
\item Орлов С.А., Цилькер Б.Я. Организация ЭВМ и систем: Учебник для вузов. 2-е изд. СПб.: Питер, 2011. 688 с.
\item Гл. 1, с. 20--53; гл. 2, с. 55--117; гл. 6, с. 230--331; гл. 7, с. 332--363; гл. 8, с. 364--385; гл. 9, с. 386--449; гл. 13, с. 526--556; гл. 15, с. 579--597.
\end{itemize}

\end{document}